\documentclass[UTF8,11pt,a4paper]{ctexart}
\usepackage{amsmath,amssymb,amsthm,geometry,enumitem,hyperref,array,longtable,booktabs,ragged2e}
\geometry{margin=2.1cm}
\setlist{itemsep=2pt,topsep=3pt}
\hypersetup{colorlinks=true,linkcolor=blue}
\setlength{\emergencystretch}{2em}
\ctexset{section={beforeskip=1.8ex,afterskip=0.8ex},subsection={beforeskip=1.2ex,afterskip=0.5ex}}
\setlength{\abovedisplayskip}{4pt}
\setlength{\belowdisplayskip}{4pt}
\setlength{\abovedisplayshortskip}{2pt}
\setlength{\belowdisplayshortskip}{2pt}
\newcommand{\APS}[2]{\par\noindent\textbf{中文。} #1\par\noindent\textbf{English.} #2\par\smallskip}
\newcommand{\EN}[1]{\par\noindent \textbf{English.} #1\par}
\newcommand{\FORMULA}[2]{\par\noindent\textbf{公式名称 / Formula name.} #1 / #2\par}
\newcommand{\READ}[1]{\par\noindent \textbf{公式英文读法 / How to read it.} #1\par}
\newcommand{\EX}[2]{\par\noindent \textbf{例子 / Example.} #1\par\EN{#2}}
\newcommand{\ASK}[4]{\par\noindent \textbf{可能问题 / Possible question.} #1\par\EN{#2}\noindent \textbf{回答要点 / Key answer.} #3\par\EN{#4}\medskip}
\theoremstyle{definition}
\newtheorem*{definition}{定义\,/\,Definition}
\newtheorem*{theorem}{定理\,/\,Theorem}
\newtheorem*{proposition}{命题\,/\,Proposition}
\newtheorem*{lemma}{引理\,/\,Lemma}
\newtheorem*{corollary}{推论\,/\,Corollary}
\newtheorem*{remark}{注记\,/\,Remark}
\renewcommand{\proofname}{证明\,/\,Proof}
\newcommand{\R}{\mathbb{R}}
\newcommand{\C}{\mathbb{C}}
\newcommand{\Z}{\mathbb{Z}}
\newcommand{\N}{\mathbb{N}}
\newcommand{\dd}{\mathrm{d}}
\newcommand{\vv}[1]{\mathbf{#1}}
\pagestyle{plain}
\title{解析几何 APS 复习讲义\\[4pt]\large 吕林根、许子道《解析几何》第五版 第 4--5 章\\[2pt]\normalsize 二次曲面与二次曲线的一般理论}
\author{Analytic Geometry: Bilingual Review}
\date{}
\begin{document}\maketitle

\section*{课程概述 | Course overview}
\APS{解析几何是用代数方程研究几何图形、又用几何直观解释代数方程的学科。本讲义覆盖吕林根、许子道《解析几何》（第五版）第 4 章``柱面、锥面、旋转曲面与二次曲面''与第 5 章``二次曲线的一般理论''：前者是空间解析几何的曲面总论，后者是平面二次曲线的完整代数--几何理论。两章的共同思想是：先由几何生成方式（一族直线、一族纬圆、一族椭圆或一族平行弦）给出曲面的定义与方程，再从方程出发用``截口''或``不变量''把图形还原出来。}{Analytic geometry studies geometric figures through algebraic equations and, conversely, interprets algebraic equations geometrically. These notes cover Chapters 4 and 5 of Lü Lingen and Xu Zidao, \emph{Analytic Geometry}, 5th edition: ``Cylinders, cones, surfaces of revolution and quadric surfaces'' and ``The general theory of conics''. The first is a survey of surfaces in space, the second a complete algebraic--geometric theory of plane conics. Both share one idea: first obtain the definition and equation of a surface from its geometric generation (a family of lines, of latitude circles, of ellipses, or of parallel chords), then recover the figure from the equation by slicing or by invariants.}
\APS{第 4 章按``直纹曲面 $\to$ 旋转曲面 $\to$ 二次曲面''展开。柱面与锥面是最基本的直纹曲面：一族平行直线沿准线滑动生成柱面，一族过定点的直线沿准线滑动生成锥面；旋转曲面由一条母线绕定轴旋转生成，核心技巧是``保留与旋转轴同名的坐标，把另一坐标换成其余两坐标的平方和开方''。随后用``平行截割法''逐一研究椭球面、单叶与双叶双曲面、椭圆抛物面与双曲抛物面：先看对称性、顶点、与坐标面的截线（主截线、主抛物线），再用一族平行平面去截，由截口曲线的变化读出整个曲面。最后证明单叶双曲面与双曲抛物面都是直纹曲面，并给出两族直母线方程与两个结构定理（异族共面或相交、同族异面）。第 5 章先把二次曲线写成带对称矩阵的一般方程，用二次型 $\Phi(X,Y)$ 刻划方向，用 $F_1,F_2$ 刻划点，由此得到渐近方向、中心、渐近线、切线、直径、主直径与主方向；再用坐标变换把方程化为三种简化方程、九种标准方程，最后用不变量 $I_1,I_2,I_3$ 与半不变量 $K_1$ 直接判定曲线类型，完成二次曲线的度量分类。}{Chapter 4 runs through ``ruled surfaces $\to$ surfaces of revolution $\to$ quadrics''. Cylinders and cones are the basic ruled surfaces: a family of parallel lines sliding along a directrix generates a cylinder, a family of lines through a fixed point generates a cone; a surface of revolution is generated by rotating a generator about an axis, and the working rule is to keep the coordinate named after the axis and replace the other coordinate by the square root of the sum of the squares of the remaining two. Then the ``method of parallel sections'' studies the ellipsoid, the hyperboloids, the elliptic and hyperbolic paraboloids one by one: symmetry, vertices and traces with the coordinate planes come first, and slicing by a family of parallel planes reveals the whole surface. Finally the one-sheeted hyperboloid and the hyperbolic paraboloid are shown to be ruled, with two families of generators and two structural theorems (generators from different families are coplanar or intersect; those from the same family are skew). Chapter 5 writes a conic as a general equation with a symmetric matrix, describes directions by the quadratic form $\Phi(X,Y)$ and points by $F_1,F_2$, and derives asymptotic directions, centre, asymptotes, tangents, diameters, principal diameters and directions; coordinate changes reduce the equation to three simplified and nine standard forms, and the invariants $I_1,I_2,I_3$ and seminvariant $K_1$ classify conics without computing the transformation.}
\APS{阅读约定：定义、定理、推论的编号与教材一致（如``定义 4.5.1''、``(4.5-2)''即书中第 4.5 节第 2 个公式），便于逐条对照原书；教材为扫描版，书内页码 $+8$ 等于 PDF 物理页码（第 4 章 $96$--$127$ 页即 PDF 第 $104$--$135$ 页，第 5 章 $128$--$169$ 页即 PDF 第 $136$--$177$ 页）。每节按``导读 $\to$ 定义 $\to$ 定理 $\to$ 证明或思路 $\to$ 算例 $\to$ 可能问题''组织，最后两节给出全书中英术语表与章节级参考文献。书中的图一律略去，改用文字说明截面与顶点位置，读者应自行对照原书插图。}{Conventions: the numbers of definitions, theorems and corollaries follow the textbook (e.g.\ ``Definition 4.5.1'' and ``(4.5-2)''), so every item can be checked against the book; the textbook is a scan and its page number plus 8 gives the PDF page (Chapter 4, pages 96--127, is PDF 104--135; Chapter 5, pages 128--169, is PDF 136--177). Each section is organised as overview $\to$ definitions $\to$ theorems $\to$ proof or idea $\to$ worked example $\to$ possible questions. The last two sections give a bilingual glossary and chapter-level references. Figures are omitted; cross-sections and vertex positions are described in words instead.}

\subsection*{主线地图 | Road map}
\begin{center}\small
\begin{tabular}{@{}>{\raggedright\arraybackslash}p{1.5cm}>{\raggedright\arraybackslash}p{4.6cm}>{\raggedright\arraybackslash}p{9.3cm}@{}}
\toprule
节 & 核心对象 Object & 关键结论 Key conclusion \\
\midrule
4.1 & 柱面 cylinder & 只含两元的三元方程表示柱面（定理 4.1.1）；射影柱面把空间曲线压成两条平面曲线 \\
4.2 & 锥面 cone & 齐次方程表示顶点在原点的锥面（定理 4.2.1）；二次锥面 $x^2/a^2+y^2/b^2-z^2/c^2=0$ \\
4.3 & 旋转曲面 surface of revolution & 替换规律；五种标准旋转曲面；环面 \\
4.4 & 椭球面 ellipsoid & 对称性与半轴；平行截割法；参数方程 \\
4.5 & 双曲面 hyperboloids & 单叶的腰椭圆；双叶的两叶 $z\ge c$ 与 $z\le -c$ \\
4.6 & 抛物面 paraboloids & 主抛物线；椭圆抛物面无心，双曲抛物面是马鞍面 \\
4.7 & 直母线 generators & 两族直母线方程；异族共面（相交）、同族异面（定理 4.7.1--4.7.2） \\
5.1 & 一般方程与记号 notation & $F=xF_1+yF_2+F_3$；矩阵 $A,A^*$；$I_1,I_2,I_3,K_1$ \\
5.2 & 渐近方向、中心、渐近线 & $\Phi(X,Y)=0$ 定渐近方向；中心方程组 (5.2-2)；按 $I_2$ 分类 \\
5.3 & 切线 tangent & 切线方程 (5.3-3)(5.3-4)；奇点 $F_1=F_2=0$ \\
5.4 & 直径 diameter & 平行弦中点轨迹是直线 $XF_1+YF_2=0$；共轭方向 (5.4-5) \\
5.5 & 主直径与主方向 & 特征方程 $\lambda^2-I_1\lambda+I_2=0$；主方向正交 \\
5.6 & 化简与分类 & 三种简化方程；九种标准方程（定理 5.6.2） \\
5.7 & 不变量 invariants & $I_1,I_2,I_3$ 不变、$K_1$ 半不变；九类判定表（定理 5.7.6） \\
\bottomrule
\end{tabular}
\end{center}

\subsection*{记号约定 | Notation}
\APS{第 4 章用 $a,b,c$（或 $p$）表示正常数，$(x,y,z)$ 为空间直角坐标。第 5 章把二次曲线的一般方程写成
\[
F(x,y)\equiv a_{11}x^2+2a_{12}xy+a_{22}y^2+2a_{13}x+2a_{23}y+a_{33}=0,
\]
其中 $xy$ 项系数写 $2a_{12}$ 是为了配方与求导方便（书中小字注）。称
\[
F_1=a_{11}x+a_{12}y+a_{13},\quad F_2=a_{12}x+a_{22}y+a_{23},\quad F_3=a_{13}x+a_{23}y+a_{33},
\]
\[
\Phi(x,y)=a_{11}x^2+2a_{12}xy+a_{22}y^2,
\]
矩阵 $A=\begin{pmatrix}a_{11}&a_{12}&a_{13}\\ a_{12}&a_{22}&a_{23}\\ a_{13}&a_{23}&a_{33}\end{pmatrix}$，$A^*=\begin{pmatrix}a_{11}&a_{12}\\ a_{12}&a_{22}\end{pmatrix}$；不变量与半不变量为
\[
I_1=a_{11}+a_{22},\qquad I_2={\begin{vmatrix}a_{11}&a_{12}\\ a_{12}&a_{22}\end{vmatrix}},\qquad I_3=\det A,
\]
\[
K_1={\begin{vmatrix}a_{11}&a_{13}\\ a_{13}&a_{33}\end{vmatrix}}+{\begin{vmatrix}a_{22}&a_{23}\\ a_{23}&a_{33}\end{vmatrix}}.
\]
注意 $F_1,F_2$ 恰是 $F$ 对 $x,y$ 偏导数的一半：$F_1=\tfrac12 F_x$，$F_2=\tfrac12 F_y$，而 $F_3$ 是``把二次项删掉后的线性部分''。}{Chapter 4 uses positive constants $a,b,c$ (or $p$) and rectangular coordinates $(x,y,z)$. Chapter 5 writes a conic as displayed, where the coefficient of $xy$ is written $2a_{12}$ for convenience in completing squares and differentiating. Define $F_1,F_2,F_3,\Phi$ and the matrices $A,A^*$ as displayed; the invariants are $I_1=a_{11}+a_{22}$, $I_2=a_{11}a_{22}-a_{12}^2$, $I_3=\det A$, and $K_1=(a_{11}a_{33}-a_{13}^2)+(a_{22}a_{33}-a_{23}^2)$. Note $F_1=\tfrac12F_x$, $F_2=\tfrac12F_y$, while $F_3$ is the linear part left after deleting the quadratic terms.}

\subsection*{预备：虚元素 | Preliminaries: complex elements}
\APS{第 5 章一开始引进虚元素：一对有序复数 $(x,y)$ 仍看作平面上的点（虚点），坐标均为共轭复数的两点称为共轭虚点，实点与虚点统称复点；直线与二次曲线可能交于两个共轭虚点。这样做的目的只有一个：让``直线与二次曲线总有两个交点（重数计算）''这一代数事实恒成立，从而使切线、直径、渐近线、共轭方向等概念有统一的代数表述。本讲义只讨论实系数方程，但允许交点、方向、切点为复的。}{Chapter 5 begins with complex elements: an ordered pair of complex numbers $(x,y)$ is still viewed as a point (an imaginary point) of the plane; two points with conjugate coordinates form a pair of conjugate imaginary points, and real and imaginary points are both called complex points; a line and a conic may meet in two conjugate imaginary points. The purpose is to make the algebraic statement ``a line meets a conic in two points counted with multiplicity'' always true, so that tangency, diameters, asymptotes and conjugate directions have a uniform algebraic description. Only real equations are studied here, but intersection points, directions and points of tangency may be complex.}

\section{\label{sec:1}柱面 | Cylinders}
\APS{柱面是最容易``由运动生成''的曲面：一条直线保持方向不变，沿一条定曲线平移，扫出的曲面就是柱面。本节要回答三个问题：怎样从准线与母线方向写出柱面方程；反过来，什么样的方程一定表示柱面（定理 4.1.1）；以及怎样用柱面把空间曲线``压''成平面曲线（射影柱面）。这三个问题是第 4 章所有曲面讨论的范式：先由生成方式定方程，再由方程读图形。}{A cylinder is the easiest surface to generate by motion: a line moving without changing its direction and sweeping along a fixed curve. This section answers three questions: how to write the equation from the directrix and the direction of the generators; conversely, which equations necessarily represent cylinders (Theorem 4.1.1); and how to flatten a space curve into plane curves by cylinders (projecting cylinders). These three questions set the pattern for all surfaces in Chapter 4: first obtain the equation from the generation, then read the figure from the equation.}

\begin{definition}[4.1.1 柱面 / cylinder]\label{def:cylinder}
在空间，由平行于定方向且与一条定曲线相交的一族平行直线所生成的曲面叫做\textbf{柱面 / cylinder}，定方向叫做柱面的\textbf{方向 / direction}，定曲线叫做柱面的\textbf{准线 / directrix}，那族平行直线中的每一条直线都叫做柱面的\textbf{母线 / generator}。
\EN{In space, the surface generated by a family of parallel lines that are parallel to a fixed direction and meet a fixed curve is called a cylinder; the fixed direction is the direction of the cylinder, the fixed curve is its directrix, and each line of the family is a generator.}
\end{definition}

\APS{求柱面方程的一般步骤。设准线为两曲面交线
\[
\Gamma:\begin{cases}F_1(x,y,z)=0,\\ F_2(x,y,z)=0,\end{cases}\tag{1}
\]
母线的方向数为 $X,Y,Z$。若 $M_1(x_1,y_1,z_1)$ 为准线上任意一点，则过 $M_1$ 的母线为
\[
\frac{x-x_1}{X}=\frac{y-y_1}{Y}=\frac{z-z_1}{Z},\tag{2}
\]
且 $M_1$ 在准线上，故
\[
F_1(x_1,y_1,z_1)=0,\qquad F_2(x_1,y_1,z_1)=0.\tag{3}
\]
(2) 与 (3) 共四个等式，从中消去参数 $x_1,y_1,z_1$，最后得到一个三元方程 $F(x,y,z)=0$，这就是以 (1) 为准线、母线方向数为 $X,Y,Z$ 的柱面方程。}{General method. Let the directrix be the intersection (1) and let the generators have direction numbers $X,Y,Z$. If $M_1(x_1,y_1,z_1)$ is any point of the directrix, the generator through it is (2), and $M_1$ satisfies (3). Eliminating the three parameters $x_1,y_1,z_1$ from the four equations (2)--(3) gives one equation $F(x,y,z)=0$, the equation of the cylinder.}

\EX{（例 1，书中第 96--97 页）柱面的准线方程为
\[
\Gamma:\begin{cases}x^2+y^2+z^2=1,\\ 2x^2+2y^2+z^2=2,\end{cases}
\]
而母线的方向数是 $-1,0,1$，求这柱面的方程。}{（Example 1, pp.~96--97）The directrix is the displayed curve and the generators have direction numbers $(-1,0,1)$; find the cylinder.}
设 $M_1(x_1,y_1,z_1)$ 是准线上的点，那么过 $M_1$ 的母线为
\[
\frac{x-x_1}{-1}=\frac{y-y_1}{0}=\frac{z-z_1}{1},
\]
且有
\[
x_1^2+y_1^2+z_1^2=1,\tag{4}
\]
\[
2x_1^2+2y_1^2+z_1^2=2.\tag{5}
\]
再设公共比值为 $t$，那么由母线的分母 $0$ 知 $y_1=y$，且
\[
x_1=x+t,\qquad y_1=y,\qquad z_1=z-t.\tag{6}
\]
把 (6) 代入 (4) 及 (5) 得
\[
(x+t)^2+y^2+(z-t)^2=1,\tag{7}
\]
\[
2(x+t)^2+2y^2+(z-t)^2=2.\tag{8}
\]
以 $2$ 乘 (7) 再减去 (8)，得
\[
(z-t)^2=0,
\]
所以
\[
t=z.\tag{9}
\]
把 (9) 代入 (7) 或 (8)，即得所求的柱面方程
\[
(x+z)^2+y^2=1,
\]
即
\[
x^2+y^2+z^2+2xz-1=0.
\]
\EN{Set the common ratio equal to $t$; since the second denominator is $0$ we have $y_1=y$, and $x_1=x+t$, $z_1=z-t$. Substituting into the two equations of the directrix and eliminating $(z-t)^2$ gives $t=z$; back-substitution gives $(x+z)^2+y^2=1$, i.e.\ $x^2+y^2+z^2+2xz-1=0$.}

\begin{theorem}[4.1.1 只含两元的三元方程表示柱面 / an equation missing one variable]\label{thm:cylinder}
在空间直角坐标系中，只含两个元（坐标）的三元方程所表示的曲面是一个柱面，它的母线平行于所缺元（坐标）的同名坐标轴。
\EN{In a rectangular coordinate system in space, a surface given by a three-variable equation that contains only two of the variables is a cylinder whose generators are parallel to the coordinate axis named by the missing variable.}
\end{theorem}

\begin{proof}[证明 / Proof]
只需证 $F(x,y)=0$ 表示的曲面是一个母线平行于 $z$ 轴的柱面。取该曲面与坐标面 $z=0$ 的交线
\[
\Gamma:\begin{cases}F(x,y)=0,\\ z=0\end{cases}\tag{12}
\]
为准线，以 $z$ 轴的方向 $(0,0,1)$ 为母线方向建立柱面。设 $M_1(x_1,y_1,0)$ 为准线 (12) 上的任意点，则过 $M_1$ 的母线为
\[
x=x_1,\qquad y=y_1,\qquad z=z_1+t,\tag{13}
\]
即消去 $t$ 后 $x=x_1$，$y=y_1$，而 $z$ 任意。又 $M_1$ 在准线 (12) 上，故 $F(x_1,y_1)=0$。把 $x_1=x,y_1=y$ 代入即得 $F(x,y)=0$，这正是原方程，所以原方程表示的正是这个柱面。同理 $G(y,z)=0$ 与 $H(x,z)=0$ 分别表示母线平行于 $x$ 轴与 $y$ 轴的柱面。
\EN{It suffices to treat $F(x,y)=0$. Take its trace on $z=0$ as directrix and the $z$-direction as generator direction. Any point $M_1(x_1,y_1,0)$ of the directrix gives the generator $x=x_1$, $y=y_1$ with $z$ free; since $F(x_1,y_1)=0$, substituting $x_1=x,y_1=y$ returns the original equation. Hence the surface is the cylinder. The cases $G(y,z)=0$ and $H(x,z)=0$ are analogous.}
\end{proof}

\APS{注记：这个定理在空间仿射坐标系中也成立（证明只用到``平行''与``坐标轴''，与度量无关）。由定理 4.1.1，下列方程都表示柱面：
\[
\frac{x^2}{a^2}+\frac{y^2}{b^2}=1,\qquad \frac{x^2}{a^2}-\frac{y^2}{b^2}=1,\qquad y^2=2px.\tag{4.1-1--3}
\]
它们与 $xOy$ 坐标面的交线分别是椭圆、双曲线与抛物线，所以依次叫做\textbf{椭圆柱面 / elliptic cylinder}、\textbf{双曲柱面 / hyperbolic cylinder}、\textbf{抛物柱面 / parabolic cylinder}；方程都是二次的，故统称\textbf{二次柱面 / quadric cylinders}。}{Remark: the theorem also holds in affine coordinates, since the proof uses only parallelism. Hence the three displayed equations define cylinders whose traces on $xOy$ are an ellipse, a hyperbola and a parabola: the elliptic, hyperbolic and parabolic cylinders, collectively the quadric cylinders.}

\subsection{空间曲线的射影柱面 | Projecting cylinders of a space curve}
\APS{设空间曲线为
\[
\Gamma:\begin{cases}F_1(x,y,z)=0,\\ F_2(x,y,z)=0,\\ F_3(x,y,z)=0.\end{cases}\tag{15}
\]
从 (15) 中依次消去一个元，可得三个含两元的方程 $F_{xy}(x,y)=0$、$F_{xz}(x,z)=0$、$F_{yz}(y,z)=0$。任取其中两个（例如 $F_{xy}=0$ 与 $F_{xz}=0$）组成方程组，它与 (15) 等价，表示同一条曲线，所以曲面 $F_{xy}(x,y)=0$ 与 $F_{xz}(x,z)=0$ 都通过 $\Gamma$。由定理 4.1.1，$F_{xy}(x,y)=0$ 是母线平行于 $z$ 轴的柱面，在直角坐标系下其母线垂直于 $xOy$ 面，称它为空间曲线 (15) 对 $xOy$ 坐标面射影的\textbf{射影柱面 / projecting cylinder}，而曲线
\[
\begin{cases}F_{xy}(x,y)=0,\\ z=0\end{cases}
\]
叫做 $\Gamma$ 在 $xOy$ 坐标面上的\textbf{射影曲线 / projection}。同理有对 $xOz$、$yOz$ 面的射影柱面与射影曲线。}{Let the curve be (15). Eliminating one variable at a time gives three two-variable equations; any two of them form a system equivalent to (15), so each defines a projecting cylinder through the curve. By Theorem 4.1.1 the surface $F_{xy}(x,y)=0$ is a cylinder with generators parallel to the $z$-axis, perpendicular to $xOy$; its trace on $z=0$ is the projection of the curve on $xOy$. The other two coordinate planes are analogous.}

\EX{（书中第 100 页）从曲线
\[
\begin{cases}2x^2+z^2+4y=4z,\\ x^2+3z^2-8y=12z\end{cases}
\]
分别消去 $y$ 及 $z$，得射影柱面
\[
\begin{cases}x^2+z^2=4z,\\ x^2+4y=0.\end{cases}
\]
消去 $y$ 时把第一式乘 $2$ 与第二式相加（$8y$ 相消）得 $5x^2+5z^2=20z$，即 $x^2+z^2=4z$；它是一个准线为 $xOz$ 面上圆 $x^2+(z-2)^2=4$、母线平行于 $y$ 轴的圆柱面。消去 $z$ 时把第一式乘 $3$ 减去第二式得 $5x^2+20y=0$，即 $x^2+4y=0$；它是一个准线为 $xOy$ 面上抛物线 $x^2=-4y$、母线平行于 $z$ 轴的抛物柱面。原曲线就是这两个柱面的交线。}{（p.~100）Adding twice the first equation to the second eliminates $y$ and gives $x^2+z^2=4z$, a circular cylinder with generators parallel to $y$; subtracting the second from three times the first eliminates $z$ and gives $x^2+4y=0$, a parabolic cylinder with generators parallel to $z$. The original curve is the intersection of these two cylinders.}

\APS{补充算例（书中例 2 的思路）：已知圆柱面的轴为 $\dfrac{x}{1}=\dfrac{y-1}{-2}=\dfrac{z+1}{-2}$，点 $(1,-2,1)$ 在此圆柱面上，求它的方程。圆柱面上的点 $M(x,y,z)$ 到轴的距离等于圆柱面的半径。轴上的定点 $P_0(0,1,-1)$ 到已知点 $(1,-2,1)$ 的距离为 $\sqrt{14}$，而 $P_0$ 到过已知点且垂直于轴的平面 $x-2y-2z-3=0$ 的距离为 $1$，故半径 $d=\sqrt{14-1^2}=\sqrt{13}$。用点到直线的距离公式 $\dfrac{|P_0M\times \vv{v}|}{|\vv{v}|}=d$，其中 $\vv{v}=(1,-2,-2)$，两边平方整理得 $8x^2+5y^2+5z^2+4xy+4xz-8yz-18y+18z-99=0$；这就是所求圆柱面的方程（书中用准线圆消参法得到同一结果）。}{Supplementary example (idea of Example 2): the axis is the displayed line and the point $(1,-2,1)$ lies on the cylinder. The distance from $M$ to the axis is the radius; the fixed point $P_0(0,1,-1)$ on the axis is at distance $\sqrt{14}$ from $(1,-2,1)$, and its distance to the plane $x-2y-2z-3=0$ is $1$, so $d=\sqrt{13}$. Squaring $|P_0M\times v|/|v|=d$ with $v=(1,-2,-2)$ gives the displayed quadric, the same result as the book's elimination method.}

\ASK{为什么``只含两个坐标的三元方程''一定是柱面？准线是否唯一？}{Why is an equation missing one coordinate necessarily a cylinder? Is the directrix unique?}{因为方程 $F(x,y)=0$ 不限制 $z$：若 $M(x,y,z)$ 在曲面上，则对任意 $t$，$(x,y,z+t)$ 也在曲面上，即整条平行于 $z$ 轴的直线都在曲面上；这些直线正是母线，它们与 $xOy$ 面的交线可取作准线。准线不唯一：任何与所有母线都相交的曲线都可以作准线，例如用斜平面 $z=kx$ 去截，得到的椭圆（相交的二次曲线）也可以作准线，但用它求方程更繁。}{Because $F(x,y)=0$ puts no restriction on $z$: if a point lies on the surface, so does the whole line parallel to the $z$-axis through it, and these lines are the generators; their trace on $xOy$ serves as a directrix. The directrix is not unique: any curve meeting every generator once may be used, e.g.\ the ellipse cut out by an oblique plane, but the resulting computation is longer.}

\section{\label{sec:2}锥面 | Cones}
\APS{锥面与柱面的对偶：把``平行直线族''换成``共点直线族''，柱面就成了锥面。公共点叫顶点，因此锥面的方程关于顶点有``齐次性''。本节先给出由顶点与准线求锥面的方法，再算一个二次锥面与一个圆锥面，最后证明``齐次方程必为以原点为顶点的锥面''（定理 4.2.1）及其平移推论。}{A cone is dual to a cylinder: replace the family of parallel lines by a family of concurrent lines. The common point is the vertex, so the equation is homogeneous about it. The section gives the method from vertex and directrix, computes a quadratic cone and a circular cone, and proves that a homogeneous equation always defines a cone with vertex at the origin.}

\begin{definition}[4.2.1 锥面 / cone]\label{def:cone}
在空间，通过一定点且与定曲线相交的一族直线所生成的曲面叫做\textbf{锥面 / cone}，这些直线都叫做锥面的\textbf{母线 / generators}，那个定点叫做锥面的\textbf{顶点 / vertex}，定曲线叫做锥面的\textbf{准线 / directrix}。
\EN{In space, the surface generated by a family of lines passing through a fixed point and meeting a fixed curve is called a cone; the lines are the generators, the fixed point is the vertex and the fixed curve is the directrix.}
\end{definition}

\APS{求锥面方程的一般方法。设准线为
\[
\Gamma:\begin{cases}F_1(x,y,z)=0,\\ F_2(x,y,z)=0,\end{cases}\tag{1}
\]
顶点为 $A(x_0,y_0,z_0)$，$M_1(x_1,y_1,z_1)$ 为准线上任意点，则母线 $AM_1$ 为
\[
\frac{x-x_0}{x_1-x_0}=\frac{y-y_0}{y_1-y_0}=\frac{z-z_0}{z_1-z_0},\tag{2}
\]
且 $M_1$ 满足 (1)。从 (1) 与 (2) 共四个等式中消去 $x_1,y_1,z_1$，得三元方程 $F(x,y,z)=0$，即以 (1) 为准线、$A$ 为顶点的锥面方程。}{General method. Let the directrix be (1) and the vertex $A(x_0,y_0,z_0)$; the generator through $M_1$ of the directrix is (2). Eliminating $x_1,y_1,z_1$ from (1)--(2) gives the equation of the cone.}

\EX{（例 1，书中第 102 页）锥面的顶点在原点，且准线为
\[
\Gamma:\begin{cases}\dfrac{x^2}{a^2}+\dfrac{y^2}{b^2}=1,\\ z=c,\end{cases}
\]
求锥面的方程。}{（Example 1, p.~102）The vertex is the origin and the directrix is the displayed ellipse in the plane $z=c$; find the cone.}
设 $M_1(x_1,y_1,z_1)$ 为准线上任意点，则过 $M_1$ 的母线为
\[
\frac{x}{x_1}=\frac{y}{y_1}=\frac{z}{z_1},\tag{4}
\]
且有
\[
\frac{x_1^2}{a^2}+\frac{y_1^2}{b^2}=1,\tag{5}
\]
\[
z_1=c.\tag{6}
\]
由 (4)、(6) 得
\[
x_1=c\frac{x}{z},\qquad y_1=c\frac{y}{z}.\tag{7}
\]
把 (7) 代入 (5)，得所求锥面方程
\[
\frac{c^2x^2}{a^2z^2}+\frac{c^2y^2}{b^2z^2}=1,
\]
或改写为
\[
\frac{x^2}{a^2}+\frac{y^2}{b^2}-\frac{z^2}{c^2}=0.\tag{4.2-1}
\]
这个锥面叫做\textbf{二次锥面 / quadratic cone}。
\EN{Parametrize the generator by (4); from $z_1=c$ we get $x_1=cx/z$, $y_1=cy/z$; substitution into the directrix gives $x^2/a^2+y^2/b^2-z^2/c^2=0$, the quadratic cone.}

\APS{显然，锥面的准线不是唯一的：和一切母线都相交的每一条曲线都可以作为它的准线。}{Clearly the directrix is not unique: any curve meeting every generator once can serve as the directrix.}

\EX{（例 2，书中第 102--103 页）已知圆锥面的顶点为 $(1,2,3)$，轴垂直于平面 $2x+2y-z+1=0$，母线与轴成 $30^\circ$ 角，试求这圆锥面的方程。}{（Example 2, pp.~102--103）The vertex is $(1,2,3)$, the axis is perpendicular to $2x+2y-z+1=0$, and every generator makes a $30^\circ$ angle with the axis; find the cone.}
设 $M(x,y,z)$ 为任一母线上的点，则母线的方向向量为
\[
\vv{v}=\{x-1,\,y-2,\,z-3\},
\]
而轴的方向即平面的法方向 $\vv{n}=\{2,2,-1\}$。根据题意，
\[
\frac{\vv{v}\cdot\vv{n}}{|\vv{v}|\,|\vv{n}|}=\pm\cos 30^\circ,
\]
即
\[
\frac{2(x-1)+2(y-2)-(z-3)}{\sqrt{(x-1)^2+(y-2)^2+(z-3)^2}\cdot 3}=\pm\frac{\sqrt3}{2}.
\]
两边平方、去分母并化简，得
\[
11(x-1)^2+11(y-2)^2+23(z-3)^2-32(x-1)(y-2)+16(x-1)(z-3)+16(y-2)(z-3)=0.
\]
\EN{The generator direction $v=(x-1,y-2,z-3)$ makes a $30^\circ$ angle with the normal $n=(2,2,-1)$. Squaring $\cos^2 30^\circ=3/4$ and clearing denominators gives the displayed quadratic equation of the circular cone.}

\begin{theorem}[4.2.1 齐次方程表示锥面 / homogeneous equations are cones]\label{thm:cone}
一个关于 $x,y,z$ 的齐次方程总表示顶点在坐标原点的锥面。
\EN{A homogeneous equation in $x,y,z$ always represents a cone with vertex at the origin.}
\end{theorem}

\begin{proof}[证明 / Proof]
设关于 $x,y,z$ 的齐次方程为
\[
F(x,y,z)=0.\tag{8}
\]
按齐次函数的定义，对任意 $t$ 有 $F(tx,ty,tz)=t^\lambda F(x,y,z)$。令 $t=0$ 得 $F(0,0,0)=0$，故曲面过原点。再设非原点的点 $M_0(x_0,y_0,z_0)$ 满足 (8)，即 $F(x_0,y_0,z_0)=0$。直线 $OM_0$ 上的点可写成 $(tx_0,ty_0,tz_0)$，代入 (8) 得
\[
F(tx_0,ty_0,tz_0)=t^\lambda F(x_0,y_0,z_0)=0,
\]
所以整条直线 $OM_0$ 都在曲面上。因此曲面 (8) 由这种通过原点的直线组成，即以原点为顶点的锥面。
\EN{For a homogeneous $F$ we have $F(tx,ty,tz)=t^\lambda F(x,y,z)$. Taking $t=0$ gives $F(0,0,0)=0$, so the origin lies on the surface. If $M_0\ne O$ satisfies (8), then every point $tM_0$ of the line $OM_0$ also satisfies it; hence the surface is a union of lines through the origin, i.e.\ a cone with vertex the origin.}
\end{proof}

\APS{注记：在特殊情况下，齐次方程可能只表示一个原点，即除原点外没有别的实点，例如 $x^2+y^2+z^2=0$；这样的曲面常叫做具有实顶点的\textbf{虚锥面 / imaginary cone}。}{Remark: in degenerate cases a homogeneous equation may represent only the origin, e.g.\ $x^2+y^2+z^2=0$; such a surface is called an imaginary cone with a real vertex.}

\begin{corollary}[顶点平移 / shifted vertex]\label{cor:coneshift}
关于 $x-x_0,\ y-y_0,\ z-z_0$ 的齐次方程表示顶点在 $(x_0,y_0,z_0)$ 的锥面。
\EN{A homogeneous equation in $x-x_0$, $y-y_0$, $z-z_0$ represents a cone with vertex $(x_0,y_0,z_0)$.}
\end{corollary}

\APS{推论可由平移 $x'=x-x_0$，$y'=y-y_0$，$z'=z-z_0$ 直接得到（书中小字注：利用 §1.5 习题第 4 题容易证明）。}{The corollary follows from the translation $x'=x-x_0$, $y'=y-y_0$, $z'=z-z_0$.}

\ASK{怎样判定一个方程表示锥面？圆锥面与一般锥面求法的区别？}{How do we recognize a cone from its equation? What is special about a circular cone?}{把方程写成关于 $x-x_0,y-y_0,z-z_0$ 的多项式：若它是齐次的，则由定理 4.2.1 及推论，曲面是以 $(x_0,y_0,z_0)$ 为顶点的锥面；若还含有一次项等非齐次项，则一般不是锥面。圆锥面是母线与轴夹角为定值的锥面，除上述一般方法外，还可用``母线与轴夹角余弦等于常数''直接列式（如例 2），不必先求准线。}{Write the equation as a polynomial in $x-x_0,y-y_0,z-z_0$: if it is homogeneous, Theorem 4.2.1 and the corollary give a cone with vertex $(x_0,y_0,z_0)$; otherwise it is not a cone in general. A circular cone is characterised by a constant angle between generator and axis, so one may equate the cosine to a constant directly, without first finding a directrix.}

\section{\label{sec:3}旋转曲面 | Surfaces of revolution}
\APS{旋转曲面的生成方式是``一条曲线绕一条定直线旋转一周''。它比柱面、锥面多了一个度量约束——旋转保持到轴的距离不变，所以方程推导中会出现``到轴上一点的距离相等''与``垂直于轴''两式。真正用于计算的是最后那条替换规律：在坐标面内的母线上保留与轴同名的坐标，把另一坐标换成其余两坐标的平方和开方。掌握这条规律，五种标准旋转曲面（旋转椭球面两种、旋转双曲面两种、旋转抛物面）与环面都能立刻写出。}{A surface of revolution is generated by rotating a curve about a fixed line. It has one metric constraint more than cylinders and cones---rotation preserves the distance to the axis---so the derivation uses ``equal distances to a point of the axis'' and ``perpendicular to the axis''. The working rule is the substitution rule: in a plane curve taken as generator, keep the coordinate named after the axis and replace the other coordinate by the square root of the sum of squares of the remaining two. With this rule, the five standard surfaces of revolution and the torus are immediate.}

\begin{definition}[4.3.1 旋转曲面 / surface of revolution]\label{def:revolution}
在空间，一条曲线 $\Gamma$ 绕着定直线 $l$ 旋转一周所生成的曲面叫做\textbf{旋转曲面 / surface of revolution}，或称回转曲面。曲线 $\Gamma$ 叫做旋转曲面的\textbf{母线 / generator}，定直线 $l$ 叫做旋转曲面的\textbf{旋转轴 / axis}，简称为轴。
\EN{The surface generated by rotating a curve $\Gamma$ about a fixed line $l$ through a full turn is a surface of revolution; $\Gamma$ is the generator and $l$ is the axis.}
\end{definition}

\APS{母线 $\Gamma$ 上的任意点 $M_1$ 在旋转时形成一个圆，这个圆也就是通过点 $M_1$ 且垂直于轴 $l$ 的平面与旋转曲面的交线，把它叫做\textbf{纬圆}（或纬线）。在通过旋转轴 $l$ 的平面上，以 $l$ 为界的每个半平面都与曲面交成一条曲线，这些曲线显然在旋转中都能彼此重合，这曲线叫做旋转面的\textbf{经线 / meridian}。}{Every point $M_1$ of the generator traces a circle, the section of the surface by the plane through $M_1$ perpendicular to the axis; this circle is a \emph{latitude circle}. Each half-plane bounded by the axis cuts the surface in a curve, and all these curves coincide under rotation; such a curve is a \emph{meridian}.}

\subsection{旋转曲面方程的推导 | Deriving the equation}
\APS{设母线为
\[
\Gamma:\begin{cases}F_1(x,y,z)=0,\\ F_2(x,y,z)=0,\end{cases}\tag{1}
\]
旋转轴为直线
\[
l:\ \frac{x-x_0}{X}=\frac{y-y_0}{Y}=\frac{z-z_0}{Z},\tag{2}
\]
其中 $P_0(x_0,y_0,z_0)$ 为轴 $l$ 上的一个定点，$X,Y,Z$ 为轴的方向数。设 $M_1(x_1,y_1,z_1)$ 为母线 $\Gamma$ 上的任意点，则过 $M_1$ 的纬圆可以看成是过 $M_1$ 且垂直于轴 $l$ 的平面与以 $P_0$ 为球心、$|P_0M_1|$ 为半径的球面的交线，所以过 $M_1$ 的纬圆方程为
\[
X(x-x_1)+Y(y-y_1)+Z(z-z_1)=0,\tag{3}
\]
\[
(x-x_0)^2+(y-y_0)^2+(z-z_0)^2=(x_1-x_0)^2+(y_1-y_0)^2+(z_1-z_0)^2.\tag{4}
\]
又因 $M_1$ 在母线 $\Gamma$ 上，所以
\[
F_1(x_1,y_1,z_1)=0,\qquad F_2(x_1,y_1,z_1)=0.\tag{5--6}
\]
从 (3)--(6) 四个等式中消去 $x_1,y_1,z_1$，最后得三元方程 $F(x,y,z)=0$，这就是以 (1) 为母线、(2) 为轴的旋转曲面方程。}{Let the generator be (1) and the axis (2). For $M_1$ on the generator, the latitude circle is the intersection of the plane (3) through $M_1$ perpendicular to the axis with the sphere (4) centred at $P_0$ through $M_1$; together with (5)--(6) and elimination of $x_1,y_1,z_1$ this gives the equation of the surface.}

\EX{（例 1，书中第 105 页）求直线 $\dfrac{x}{2}=\dfrac{y}{1}=\dfrac{z-1}{0}$ 绕直线 $x=y=z$ 旋转所得的旋转曲面的方程。}{（Example 1, p.~105）Rotate the line $x/2=y/1=(z-1)/0$ about the line $x=y=z$; find the surface.}
设 $M_1(x_1,y_1,z_1)$ 是母线上的任意点。因为旋转轴通过原点，所以过 $M_1$ 的纬圆方程是
\[
(x-x_1)+(y-y_1)+(z-z_1)=0,\tag{7}
\]
\[
x^2+y^2+z^2=x_1^2+y_1^2+z_1^2.\tag{8}
\]
由于 $M_1$ 在母线上，所以又有 $\dfrac{x_1}{2}=\dfrac{y_1}{1}=\dfrac{z_1-1}{0}$，即
\[
x_1=2y_1,\qquad z_1=1.\tag{9}
\]
由 (7)、(8)、(9) 三式消去 $x_1,y_1,z_1$：由 (7) 得 $x_1+y_1+z_1=x+y+z$，即 $3y_1+1=x+y+z$，故 $y_1=\frac13(x+y+z-1)$，$x_1=\frac23(x+y+z-1)$，$z_1=1$。代入 (8) 得
\[
x^2+y^2+z^2-1=\frac59(x+y+z-1)^2,
\]
即
\[
2(x^2+y^2+z^2)-5(xy+xz+yz)+5(x+y+z)-7=0.
\]
\EN{Since the axis passes through the origin, the latitude circle is given by (7)--(8); from the generator $x_1=2y_1$, $z_1=1$. Eliminating $x_1,y_1,z_1$ gives $x^2+y^2+z^2-1=\frac59(x+y+z-1)^2$, i.e.\ the displayed quadratic surface.}

\subsection{替换规律 | The substitution rule}
\APS{由于旋转曲面的经线总可以作为最初的母线来产生旋转曲面，因此为了方便，今后总是取旋转曲面的某一条经线（显然是平面曲线）作为母线。在直角坐标系下导出旋转曲面的方程时，又常把母线所在平面取作坐标面、旋转轴取作坐标轴，这时旋转曲面的方程具有特殊的形式。}{Since a meridian may serve as the original generator, we henceforth take a plane meridian as generator, put its plane as a coordinate plane and the axis as a coordinate axis.}

\APS{如图 4--6，设母线为
\[
\Gamma:\begin{cases}F(y,z)=0,\\ x=0,\end{cases}\tag{10}
\]
旋转轴为 $y$ 轴：
\[
\frac{x}{0}=\frac{y}{1}=\frac{z}{0}.\tag{11}
\]
若 $M_1(0,y_1,z_1)$ 为母线 $\Gamma$ 上的任意点，则过 $M_1$ 的纬圆为
\[
y-y_1=0,\tag{12}
\]
\[
x^2+y^2+z^2=y_1^2+z_1^2,\tag{13}
\]
且有
\[
F(y_1,z_1)=0.\tag{14}
\]
从 (12)、(13)、(14) 消去参数 $y_1,z_1$，得绕 $y$ 轴旋转的旋转曲面方程
\[
F\bigl(y,\ \pm\sqrt{x^2+z^2}\bigr)=0.\tag{15}
\]
同样，把曲线 $\Gamma$ 绕 $z$ 轴旋转所得的旋转曲面的方程是
\[
F\bigl(\pm\sqrt{x^2+y^2},\ z\bigr)=0.\tag{16}
\]
于是得到如下的\textbf{替换规律 / substitution rule}：当坐标面上的曲线 $\Gamma$ 绕此坐标面里的一个坐标轴旋转时，为了求出旋转曲面的方程，只要将曲线 $\Gamma$ 在坐标面里的方程\textbf{保留和旋转轴同名的坐标}，而以\textbf{其他两个坐标平方和的平方根}来代替方程中的另一坐标。}{With the meridian (10) in the $yOz$ plane and the $y$-axis as axis, the latitude circle is (12)--(13) plus (14); eliminating gives $F(y,\pm\sqrt{x^2+z^2})=0$. Rotation about the $z$-axis gives $F(\pm\sqrt{x^2+y^2},z)=0$. Rule: keep the coordinate named after the axis and replace the other coordinate by the square root of the sum of squares of the remaining two.}

\subsection{五种标准旋转曲面与环面 | Five standard surfaces and the torus}
\EX{（例 2，书中第 106 页）将椭圆
\[
\Gamma:\begin{cases}\dfrac{x^2}{a^2}+\dfrac{y^2}{b^2}=1\quad(a>b),\\ z=0\end{cases}\tag{17}
\]
分别绕长轴（即 $x$ 轴）与短轴（即 $y$ 轴）旋转，求所得旋转曲面的方程。}{（Example 2, p.~106）Rotate the ellipse (17) about its major axis (the $x$-axis) and about its minor axis (the $y$-axis).}
因为旋转轴是 $x$ 轴，同名坐标就是 $x$。在方程 $\frac{x^2}{a^2}+\frac{y^2}{b^2}=1$ 中保留坐标 $x$ 不变，用 $\pm\sqrt{y^2+z^2}$ 代 $y$，便将椭圆 (17) 绕长轴旋转的曲面方程写成
\[
\frac{x^2}{a^2}+\frac{y^2}{b^2}+\frac{z^2}{b^2}=1,\tag{4.3-1}
\]
叫做\textbf{长形旋转椭球面 / prolate spheroid}。同样，将椭圆 (17) 绕短轴（即 $y$ 轴）旋转的曲面方程为
\[
\frac{x^2}{a^2}+\frac{y^2}{b^2}+\frac{z^2}{a^2}=1,\tag{4.3-2}
\]
叫做\textbf{扁形旋转椭球面 / oblate spheroid}。
\EN{About the $x$-axis: keep $x$, replace $y$ by $\pm\sqrt{y^2+z^2}$, giving (4.3-1). About the $y$-axis similarly gives (4.3-2).}

\EX{（例 3，书中第 106--107 页）将双曲线
\[
\Gamma:\begin{cases}\dfrac{y^2}{b^2}-\dfrac{z^2}{c^2}=1,\\ x=0\end{cases}\tag{18}
\]
绕虚轴（即 $z$ 轴）旋转的旋转曲面方程为
\[
\frac{x^2}{b^2}+\frac{y^2}{b^2}-\frac{z^2}{c^2}=1,\tag{4.3-3}
\]
叫做\textbf{单叶旋转双曲面 / hyperboloid of one sheet of revolution}；绕实轴（即 $y$ 轴）旋转的旋转曲面方程为
\[
\frac{y^2}{b^2}-\frac{x^2}{c^2}-\frac{z^2}{c^2}=1,\tag{4.3-4}
\]
叫做\textbf{双叶旋转双曲面 / hyperboloid of two sheets of revolution}。}{（Example 3）Rotating (18) about the imaginary axis (the $z$-axis) gives (4.3-3), a one-sheeted hyperboloid of revolution; about the real axis (the $y$-axis) gives (4.3-4), a two-sheeted one.}

\EX{（例 4，书中第 107 页）将抛物线
\[
\Gamma:\begin{cases}y^2=2pz,\\ x=0\end{cases}\tag{19}
\]
绕它的对称轴（$z$ 轴）旋转的旋转曲面方程为
\[
x^2+y^2=2pz,\tag{4.3-5}
\]
曲面 (4.3-5) 叫做\textbf{旋转抛物面 / paraboloid of revolution}。}{（Example 4）Rotating the parabola (19) about its axis gives $x^2+y^2=2pz$, the paraboloid of revolution.}

\EX{（例 5，书中第 107--108 页）将圆
\[
\Gamma:\begin{cases}(y-b)^2+z^2=a^2\quad(b>a>0),\\ x=0\end{cases}\tag{20}
\]
绕 $z$ 轴旋转，求所得旋转曲面的方程。}{（Example 5, pp.~107--108）Rotate the circle (20) about the $z$-axis.}
因为绕 $z$ 轴旋转，所以在方程 $(y-b)^2+z^2=a^2$ 中保留 $z$ 不变，而 $y$ 用 $\pm\sqrt{x^2+y^2}$ 代入，得
\[
\bigl(\pm\sqrt{x^2+y^2}-b\bigr)^2+z^2=a^2,
\]
即
\[
x^2+y^2+z^2+b^2-a^2=\pm 2b\sqrt{x^2+y^2},
\]
或
\[
\bigl(x^2+y^2+z^2+b^2-a^2\bigr)^2=4b^2\bigl(x^2+y^2\bigr).
\]
这样的曲面叫做\textbf{环面 / torus}，它的形状像救生圈。
\EN{Rotation about the $z$-axis replaces $y$ by $\pm\sqrt{x^2+y^2}$, which gives the torus $(x^2+y^2+z^2+b^2-a^2)^2=4b^2(x^2+y^2)$.}

\ASK{为什么旋转曲面的方程可以只由``保留同名坐标、另一坐标换成平方和开方''得到？}{Why does the substitution rule produce the equation of the surface of revolution?}{因为旋转把点 $(x_1,y_1,z_1)$ 变成同一纬圆上的点，而纬圆由两条性质刻划：到轴的距离（的平方）不变，以及沿轴方向的坐标（与轴同名的坐标）不变。若母线在 $yOz$ 面内、轴为 $y$ 轴，则 $y$ 坐标不变，而 $x^2+z^2=y_1^2$ 等于原点到母线上点的``垂直距离的平方''，于是把母线方程中的 $z_1$ 换成 $\pm\sqrt{x^2+z^2}$ 即可。一般情形可先平移、旋转坐标把轴化为坐标轴。}{Rotation moves a point along its latitude circle, which is characterised by two facts: the distance to the axis is unchanged and the coordinate along the axis is unchanged. For a meridian in $yOz$ with the $y$-axis as axis, $y$ is fixed and $x^2+z^2$ equals the squared distance from the point of the meridian to the axis, so one replaces the other coordinate of the meridian equation by $\pm\sqrt{x^2+z^2}$.}

\section{\label{sec:4}椭球面 | The ellipsoid}
\APS{椭球面是二次曲面中最``有界''的一个：它完全落在长方体 $|x|\le a$，$|y|\le b$，$|z|\le c$ 内，关于三个坐标面都对称。研究它的标准流程是本节要示范的``平行截割法''：先用三个坐标面截出主截线（三个椭圆），再用平行于坐标面的平面去截，看截口如何随高度变化；截口椭圆的半轴与两端点在主截线上的滑动，就把整个曲面``撑''出来了。}{The ellipsoid is the most ``bounded'' quadric: it lies in the box $|x|\le a$, $|y|\le b$, $|z|\le c$ and is symmetric about the three coordinate planes. The method demonstrated here is that of parallel sections: cut by the coordinate planes to get the principal ellipses, then cut by planes parallel to a coordinate plane and watch how the section changes with height; the semi-axes and endpoints sliding along the principal ellipses ``span'' the whole surface.}

\begin{definition}[4.4.1 椭球面 / ellipsoid]\label{def:ellipsoid}
在直角坐标系下，由方程
\[
\frac{x^2}{a^2}+\frac{y^2}{b^2}+\frac{z^2}{c^2}=1\tag{4.4-1}
\]
所表示的曲面叫做\textbf{椭球面 / ellipsoid}，或称椭圆面，方程 (4.4-1) 叫做椭球面的\textbf{标准方程 / standard equation}，其中 $a,b,c$ 为任意的正常数，通常假定 $a\ge b\ge c$。
\EN{In rectangular coordinates the surface given by (4.4-1), with positive constants $a,b,c$ and usually $a\ge b\ge c$, is the ellipsoid, and (4.4-1) is its standard equation.}
\end{definition}

\APS{对称性。因为方程 (4.4-1) 仅含坐标的平方项，可见当 $(x,y,z)$ 满足 (4.4-1) 时，点 $(\pm x,\pm y,\pm z)$ 也一定满足，其中正负号可任意选取，所以椭球面 (4.4-1) 关于三坐标面、三坐标轴与坐标原点都对称。椭球面的对称平面、对称轴与对称中心分别叫做它的\textbf{主平面 / principal planes}、\textbf{主轴 / principal axes}与\textbf{中心 / centre}。}{Symmetry: only squares occur, so with $(x,y,z)$ all sign changes $(\pm x,\pm y,\pm z)$ also lie on the surface; it is symmetric about the three coordinate planes, the three axes and the origin, called its principal planes, principal axes and centre.}

\APS{顶点与轴。椭球面 (4.4-1) 与它的三对称轴即坐标轴的交点分别为 $(\pm a,0,0)$、$(0,\pm b,0)$、$(0,0,\pm c)$，这六个点叫做椭球面的\textbf{顶点 / vertices}。同一条对称轴上的两顶点间的线段以及它们的长度 $2a,2b,2c$ 叫做椭球面的\textbf{轴 / axes}；轴的一半，即中心与各顶点间的线段及它们的长度 $a,b,c$ 叫做椭球面的\textbf{半轴 / semi-axes}。当 $a>b>c$ 时，$2a,2b,2c$ 分别叫做\textbf{长轴、中轴、短轴}，$a,b,c$ 分别叫做\textbf{长半轴、中半轴、短半轴}。任何两轴相等的椭球面一定是旋转椭球面，而三轴相等的椭球面就是球面：当 $a>b=c$ 时 (4.4-1) 变成 (4.3-1)，是长形旋转椭球面；当 $a=c>b$ 时变成 (4.3-2)，是扁形旋转椭球面；当 $a=b=c$ 时变成球面。三轴不等的椭球面叫\textbf{三轴椭球面 / triaxial ellipsoid}。}{Vertices: the six intersections $(\pm a,0,0),(0,\pm b,0),(0,0,\pm c)$; the segments $2a,2b,2c$ are the axes and $a,b,c$ the semi-axes. If $a>b>c$ they are the major, mean and minor axes and semi-axes. Two equal semi-axes give a spheroid, three equal a sphere; the triaxial case is the general ellipsoid.}

\APS{有界性。因为椭球面上任何一点的坐标 $(x,y,z)$ 总有
\[
|x|\le a,\qquad |y|\le b,\qquad |z|\le c,
\]
因此椭球面完全被封闭在一个长方体的内部，这个长方体由六个平面 $x=\pm a$，$y=\pm b$，$z=\pm c$ 所组成。}{Boundedness: $|x|\le a$, $|y|\le b$, $|z|\le c$, so the ellipsoid lies in the box bounded by $x=\pm a$, $y=\pm b$, $z=\pm c$.}

\APS{\textbf{平行截割法 / method of parallel sections}：为了看出曲面的大致形状，考虑曲面与一组平行平面的交线，当这些平面曲线的形状都清楚时，曲面的大致形状也就看出来了；为方便起见，常取与坐标面平行的一组平面。用坐标面 $z=0$，$y=0$，$x=0$ 分别截割椭球面 (4.4-1)，所得截口都是椭圆，它们的方程分别是
\[
\begin{cases}\dfrac{x^2}{a^2}+\dfrac{y^2}{b^2}=1,\\ z=0;\end{cases}\qquad
\begin{cases}\dfrac{x^2}{a^2}+\dfrac{z^2}{c^2}=1,\\ y=0;\end{cases}\qquad
\begin{cases}\dfrac{y^2}{b^2}+\dfrac{z^2}{c^2}=1,\\ x=0.\end{cases}\tag{1--3}
\]
椭圆 (1)、(2)、(3) 叫做椭球面 (4.4-1) 的\textbf{主截线 / principal sections}（或主椭圆）。}{Method of parallel sections: sections by $z=0$, $y=0$, $x=0$ are the ellipses (1)--(3), the principal sections.}

\APS{取平行于 $xOy$ 坐标面的一组平行平面 $z=h$ 截割 (4.4-1)，得到的截线方程是
\[
\begin{cases}\dfrac{x^2}{a^2}+\dfrac{y^2}{b^2}=1-\dfrac{h^2}{c^2},\\ z=h.\end{cases}\tag{4}
\]
当 $|h|>c$ 时 (4) 无图形，这表示平面 $z=h$ 与椭球面不相交；当 $|h|=c$ 时 (4) 的图形是平面 $z=h$ 上的一个点 $(0,0,c)$ 或 $(0,0,-c)$；当 $|h|<c$ 时 (4) 的图形是一个椭圆，这个椭圆的两半轴分别是
\[
a\sqrt{1-\frac{h^2}{c^2}}\quad\text{与}\quad b\sqrt{1-\frac{h^2}{c^2}},
\]
两轴的端点分别是 $\bigl(\pm a\sqrt{1-\frac{h^2}{c^2}},0,h\bigr)$ 与 $\bigl(0,\pm b\sqrt{1-\frac{h^2}{c^2}},h\bigr)$，容易知道两轴的端点分别在椭圆 (2) 与 (3) 上。这样，椭球面 (4.4-1) 可以看成是由一个椭圆的变动（大小位置都改变）而产生的，这个椭圆在变动中保持所在平面与 $xOy$ 面平行，且两轴的端点分别在另外两个定椭圆 (2) 与 (3) 上滑动。}{For $|h|>c$ there is no section, for $|h|=c$ a single point, and for $|h|<c$ an ellipse with semi-axes $a\sqrt{1-h^2/c^2}$ and $b\sqrt{1-h^2/c^2}$, whose endpoints slide along (2) and (3). Thus the ellipsoid is swept by a varying ellipse parallel to $xOy$.}

\APS{参数方程。椭球面的方程除了用标准方程 (4.4-1) 表达外，有时也用参数方程
\[
\begin{cases}x=a\cos\theta\cos\varphi,\\ y=b\cos\theta\sin\varphi,\\ z=c\sin\theta,\end{cases}\tag{4.4-2}
\]
来表达，其中 $\theta\ \bigl(-\frac{\pi}{2}\le\theta\le\frac{\pi}{2}\bigr)$，$\varphi\ (0\le\varphi<2\pi)$ 为参数。从 (4.4-2) 中消去参数 $\theta,\varphi$，就得到 (4.4-1)。}{Parametrization (4.4-2) with $-\pi/2\le\theta\le\pi/2$ and $0\le\varphi<2\pi$; elimination of the parameters returns (4.4-1).}

\EX{（书中第 110 页）已知椭球面的轴与坐标轴重合，且通过椭圆 $\dfrac{x^2}{9}+\dfrac{y^2}{16}=1$，$z=0$ 与点 $M(1,2,\sqrt{23})$，求这个椭球面的方程。}{（p.~110）The axes of the ellipsoid coincide with the coordinate axes, and it passes through the ellipse $x^2/9+y^2/16=1$ in $z=0$ and through $M(1,2,\sqrt{23})$; find it.}
因为所求椭球面的轴与三坐标轴重合，所以设所求椭球面的方程为
\[
\frac{x^2}{a^2}+\frac{y^2}{b^2}+\frac{z^2}{c^2}=1.
\]
它与 $xOy$ 面的交线为椭圆
\[
\begin{cases}\dfrac{x^2}{a^2}+\dfrac{y^2}{b^2}=1,\\ z=0,\end{cases}
\]
与已知椭圆比较知
\[
a^2=9,\qquad b^2=16.
\]
又因为椭球面通过点 $M(1,2,\sqrt{23})$，所以又有
\[
\frac19+\frac{4}{16}+\frac{23}{c^2}=1,
\]
即 $\frac{23}{c^2}=\frac{23}{36}$，所以
\[
c^2=36.
\]
因此所求椭球面的方程为
\[
\frac{x^2}{9}+\frac{y^2}{16}+\frac{z^2}{36}=1.
\]
\EN{Write the standard equation. Its trace on $z=0$ must equal the given ellipse, so $a^2=9$, $b^2=16$. Passing through $M$ gives $1/9+4/16+23/c^2=1$, hence $c^2=36$ and the ellipsoid is $x^2/9+y^2/16+z^2/36=1$.}

\ASK{椭球面与球面的关系？为什么三轴不等的椭球面不能由旋转得到？}{How does the ellipsoid relate to the sphere? Why is a triaxial ellipsoid not a surface of revolution?}{球面是 $a=b=c$ 的椭球面；旋转椭球面是恰有两轴相等的椭球面（$a=b$ 时绕 $z$ 轴旋转，$a=c$ 时绕 $y$ 轴旋转，$b=c$ 时绕 $x$ 轴旋转）。三轴不等时，过任意轴的截面都不是圆，旋转要求存在一条轴使所有垂直于它的截面都是圆，故不可能由旋转产生。}{The sphere is the case $a=b=c$; a spheroid has exactly two equal semi-axes and is obtained by rotation about the axis corresponding to the third. If the three semi-axes differ, no perpendicular section is a circle, so no rotation can produce it.}

\section{\label{sec:5}双曲面 | Hyperboloids}
\APS{双曲面是椭球面把某一项右边的 $1$ 改成 $-1$ 得到的：改一个符号得单叶双曲面，改两个符号（或等价地右边为 $-1$）得双叶双曲面。单叶的关键特征是``腰椭圆''：它与 $z$ 轴不相交，却由一族越来越大的椭圆组成，并且含两族直母线（§4.7）；双叶的关键特征是曲面分成 $z\ge c$ 与 $z\le -c$ 两叶。本节仍用平行截割法，重点看 $z=h$ 与 $y=h$ 两种截线的形状变化。}{A hyperboloid is obtained from the ellipsoid by changing the sign of one term: one sign change gives the one-sheeted hyperboloid, two sign changes (equivalently right side $-1$) the two-sheeted one. The key feature of one sheet is the waist ellipse: the surface misses the $z$-axis yet consists of ellipses growing without bound, and it carries two families of generators (§4.7). The two-sheeted one splits into the sheets $z\ge c$ and $z\le -c$. Again we use parallel sections, watching the sections $z=h$ and $y=h$.}

\subsection{单叶双曲面 | Hyperboloid of one sheet}
\begin{definition}[4.5.1 单叶双曲面 / hyperboloid of one sheet]\label{def:one-sheet}
在直角坐标系下，由方程
\[
\frac{x^2}{a^2}+\frac{y^2}{b^2}-\frac{z^2}{c^2}=1\tag{4.5-1}
\]
所表示的曲面叫做\textbf{单叶双曲面 / hyperboloid of one sheet}，方程 (4.5-1) 叫做单叶双曲面的标准方程，其中 $a,b,c$ 是任意的正常数。
\EN{In rectangular coordinates the surface (4.5-1), with positive $a,b,c$, is the hyperboloid of one sheet and (4.5-1) is its standard equation.}
\end{definition}

\APS{显然，单叶双曲面 (4.5-1) 与椭球面 (4.4-1) 一样，关于三坐标面、三坐标轴以及坐标原点都对称。它与 $z$ 轴不相交，与 $x$ 轴、$y$ 轴分别交于点 $(\pm a,0,0)$ 与 $(0,\pm b,0)$，这四点叫做单叶双曲面的\textbf{顶点 / vertices}。}{Like the ellipsoid it is symmetric about the three coordinate planes, the three axes and the origin; it misses the $z$-axis and meets the $x$- and $y$-axes at $(\pm a,0,0)$ and $(0,\pm b,0)$, the four vertices.}

\APS{用三个坐标面 $z=0$，$y=0$，$x=0$ 分别截割 (4.5-1)，所得截线顺次为
\[
\begin{cases}\dfrac{x^2}{a^2}+\dfrac{y^2}{b^2}=1,\\ z=0;\end{cases}\qquad
\begin{cases}\dfrac{x^2}{a^2}-\dfrac{z^2}{c^2}=1,\\ y=0;\end{cases}\qquad
\begin{cases}\dfrac{y^2}{b^2}-\dfrac{z^2}{c^2}=1,\\ x=0.\end{cases}\tag{1--3}
\]
(1) 为 $xOy$ 面上的椭圆，叫做单叶双曲面的\textbf{腰椭圆 / waist ellipse}；(2) 与 (3) 分别为 $xOz$ 面与 $yOz$ 面上的双曲线，这两条双曲线有着共同的虚轴与虚轴长。}{The traces are the waist ellipse (1) in $xOy$ and the hyperbolas (2), (3) in $xOz$ and $yOz$ with a common conjugate axis.}

\APS{用一组平行平面 $z=h$（$h$ 为任意实数）来截割 (4.5-1)，得到椭圆
\[
\begin{cases}\dfrac{x^2}{a^2}+\dfrac{y^2}{b^2}=1+\dfrac{h^2}{c^2},\\ z=h,\end{cases}\tag{4}
\]
它的两半轴分别是 $a\sqrt{1+\frac{h^2}{c^2}}$ 与 $b\sqrt{1+\frac{h^2}{c^2}}$，两轴的端点分别为
\[
\Bigl(\pm a\sqrt{1+\tfrac{h^2}{c^2}},\,0,\,h\Bigr)\quad\text{与}\quad\Bigl(0,\,\pm b\sqrt{1+\tfrac{h^2}{c^2}},\,h\Bigr),
\]
容易知道这两对端点分别在双曲线 (2) 与 (3) 上。这样，单叶双曲面可以看成是由一个椭圆的变动（大小位置都改变）而产生的，这个椭圆在变动中保持所在的平面与 $xOy$ 面平行，且两对顶点分别沿着两个定双曲线 (2) 与 (3) 滑动。}{Sections $z=h$ are ellipses (4) with semi-axes $a\sqrt{1+h^2/c^2}$, $b\sqrt{1+h^2/c^2}$ whose endpoints slide along (2) and (3); hence the surface is swept by an ellipse parallel to $xOy$.}

\APS{用平行于 $xOz$ 面的平面 $y=h$ 来截割 (4.5-1)，截线方程为
\[
\begin{cases}\dfrac{x^2}{a^2}-\dfrac{z^2}{c^2}=1-\dfrac{h^2}{b^2},\\ y=h.\end{cases}\tag{5}
\]
当 $|h|<b$ 时，截线 (5) 为双曲线，它的实轴平行于 $x$ 轴，实半轴长为 $\frac{a}{b}\sqrt{b^2-h^2}$，虚轴平行于 $z$ 轴，虚半轴长为 $\frac{c}{b}\sqrt{b^2-h^2}$，且双曲线 (5) 的顶点 $\bigl(\pm\frac{a}{b}\sqrt{b^2-h^2},h,0\bigr)$ 在腰椭圆 (1) 上。当 $|h|>b$ 时，截线 (5) 仍为双曲线，但它的实轴平行于 $z$ 轴，实半轴长为 $\frac{c}{b}\sqrt{h^2-b^2}$，虚轴平行于 $x$ 轴，虚半轴长为 $\frac{a}{b}\sqrt{h^2-b^2}$，而且它的顶点 $\bigl(0,h,\pm\frac{c}{b}\sqrt{h^2-b^2}\bigr)$ 在双曲线 (3) 上。当 $|h|=b$ 时，(5) 变成
\[
\begin{cases}\dfrac{x^2}{a^2}-\dfrac{z^2}{c^2}=0,\\ y=b\end{cases}\quad\text{或}\quad
\begin{cases}\dfrac{x^2}{a^2}-\dfrac{z^2}{c^2}=0,\\ y=-b,\end{cases}
\]
这是两条直线
\[
\begin{cases}\dfrac{x}{a}\pm\dfrac{z}{c}=0,\\ y=b\end{cases}\quad\text{或}\quad
\begin{cases}\dfrac{x}{a}\pm\dfrac{z}{c}=0,\\ y=-b.\end{cases}
\]
如果 $h=b$，那么两条直线交于点 $(0,b,0)$；如果 $h=-b$，那么交于 $(0,-b,0)$。用平行于 $yOz$ 面的平面来截割，结果与用平行于 $xOz$ 面的平面完全类似。}{For $|h|<b$ the section (5) is a hyperbola with real axis parallel to $x$ and vertices on the waist ellipse; for $|h|>b$ the real axis is parallel to $z$ and the vertices lie on (3); for $|h|=b$ the section degenerates into a pair of lines meeting at $(0,\pm b,0)$. The case of planes parallel to $yOz$ is analogous.}

\APS{在方程 (4.5-1) 中，如果 $a=b$，那么它就写成单叶旋转双曲面 (4.3-3)。方程
\[
\frac{x^2}{a^2}-\frac{y^2}{b^2}+\frac{z^2}{c^2}=1\qquad\text{与}\qquad-\frac{x^2}{a^2}+\frac{y^2}{b^2}+\frac{z^2}{c^2}=1
\]
所表示的图形，也都是单叶双曲面（只是``腰''的方向不同）。}{If $a=b$, (4.5-1) is the one-sheeted hyperboloid of revolution (4.3-3); the displayed sign variants are also one-sheeted hyperboloids with the waist in another coordinate plane.}

\subsection{双叶双曲面 | Hyperboloid of two sheets}
\begin{definition}[4.5.2 双叶双曲面 / hyperboloid of two sheets]\label{def:two-sheet}
在直角坐标系下，由方程
\[
\frac{x^2}{a^2}+\frac{y^2}{b^2}-\frac{z^2}{c^2}=-1\tag{4.5-2}
\]
所表示的图形叫做\textbf{双叶双曲面 / hyperboloid of two sheets}，方程 (4.5-2) 叫做双叶双曲面的标准方程，其中 $a,b,c$ 是任意的正常数。
\EN{In rectangular coordinates the surface (4.5-2), with positive $a,b,c$, is the hyperboloid of two sheets.}
\end{definition}

\APS{因为方程 (4.5-2) 仅含坐标的平方项，因此这个曲面关于三坐标面、三坐标轴以及坐标原点都对称；而且曲面与 $x$ 轴、$y$ 轴都不相交，只与 $z$ 轴相交于两点 $(0,0,\pm c)$，这两点叫做双叶双曲面 (4.5-2) 的\textbf{顶点 / vertices}。从 (4.5-2) 容易知道，曲面上的点恒有 $z^2\ge c^2$，因此曲面分成两叶 $z\ge c$ 与 $z\le -c$。坐标面 $z=0$ 与曲面 (4.5-2) 不相交，而其他两个坐标面 $y=0$ 与 $x=0$ 分别交曲面于两条双曲线
\[
\begin{cases}\dfrac{z^2}{c^2}-\dfrac{x^2}{a^2}=1,\\ y=0;\end{cases}\qquad
\begin{cases}\dfrac{z^2}{c^2}-\dfrac{y^2}{b^2}=1,\\ x=0.\end{cases}\tag{6--7}
\]
}{Only squares occur, so the surface is symmetric about the three coordinate planes, the three axes and the origin; it misses the $x$- and $y$-axes and meets the $z$-axis at $(0,0,\pm c)$, the vertices. Since $z^2\ge c^2$, the surface splits into the sheets $z\ge c$ and $z\le -c$; $z=0$ misses it, while $y=0$ and $x=0$ cut the hyperbolas (6--7).}

\APS{如果用一组平行于 $xOy$ 面的平行平面 $z=h\ (|h|\ge c)$ 来截割曲面 (4.5-2)，得截线方程为
\[
\begin{cases}\dfrac{x^2}{a^2}+\dfrac{y^2}{b^2}=\dfrac{h^2}{c^2}-1,\\ z=h.\end{cases}\tag{8}
\]
当 $|h|=c$ 时截得的图形为一点；当 $|h|>c$ 时截线为椭圆，它的两半轴为
\[
a\sqrt{\frac{h^2}{c^2}-1}\quad\text{与}\quad b\sqrt{\frac{h^2}{c^2}-1},
\]
这时椭圆 (8) 的两轴的端点 $\bigl(\pm a\sqrt{\frac{h^2}{c^2}-1},0,h\bigr)$ 与 $\bigl(0,\pm b\sqrt{\frac{h^2}{c^2}-1},h\bigr)$ 分别在双曲线 (6) 与 (7) 上。因此双叶双曲面可以看成是由一个椭圆的变动（大小位置都改变）而产生的，这个椭圆在变动中保持所在平面平行于 $xOy$ 面，且两轴的端点分别沿着双曲线 (6)、(7) 滑动。}{For $|h|\ge c$ the section (8) is a point when $|h|=c$ and an ellipse with semi-axes $a\sqrt{h^2/c^2-1}$, $b\sqrt{h^2/c^2-1}$ when $|h|>c$, its endpoints sliding along (6) and (7); the surface is swept by an ellipse parallel to $xOy$.}

\APS{在方程 (4.5-2) 中，如果 $a=b$，那么截线 (8) 为一圆，曲面就是一个双叶旋转双曲面。方程
\[
\frac{x^2}{a^2}-\frac{y^2}{b^2}+\frac{z^2}{c^2}=-1\qquad\text{与}\qquad-\frac{x^2}{a^2}+\frac{y^2}{b^2}+\frac{z^2}{c^2}=-1
\]
所表示的图形，也都是双叶双曲面。单叶双曲面与双叶双曲面统称为\textbf{双曲面 / hyperboloids}。}{If $a=b$ the sections (8) are circles and the surface is a two-sheeted hyperboloid of revolution; the displayed sign variants are also two-sheeted hyperboloids. One- and two-sheeted hyperboloids are collectively called hyperboloids.}

\EX{（书中第 114--115 页）用一组平行平面 $z=h$（$h$ 为任意实数）截割单叶双曲面
\[
\frac{x^2}{a^2}+\frac{y^2}{b^2}-\frac{z^2}{c^2}=1\quad(a>b),
\]
得一族椭圆，求这些椭圆焦点的轨迹。}{（pp.~114--115）Cut the one-sheeted hyperboloid with $a>b$ by the planes $z=h$; the sections are ellipses. Find the locus of their foci.}
这一族椭圆的方程为
\[
\frac{x^2}{a^2\bigl(1+\frac{h^2}{c^2}\bigr)}+\frac{y^2}{b^2\bigl(1+\frac{h^2}{c^2}\bigr)}=1,
\]
因为 $a>b$，所以椭圆的长半轴为 $a\sqrt{1+\frac{h^2}{c^2}}$（沿 $x$ 轴方向），短半轴为 $b\sqrt{1+\frac{h^2}{c^2}}$，从而椭圆焦点的坐标为
\[
\Bigl(\pm\sqrt{(a^2-b^2)\Bigl(1+\frac{h^2}{c^2}\Bigr)},\ 0,\ h\Bigr).
\]
设焦点坐标为 $(x,0,z)$，则 $z=h$，$x^2=(a^2-b^2)\bigl(1+\frac{z^2}{c^2}\bigr)$，消去参数 $h$ 得
\[
\frac{x^2}{a^2-b^2}-\frac{z^2}{c^2}=1.
\]
显然这族椭圆焦点的轨迹是一条在坐标面 $xOz$ 上的双曲线，双曲线的实轴为 $x$ 轴，虚轴为 $z$ 轴。
\EN{The sections are ellipses with semi-axes $a\sqrt{1+h^2/c^2}$ and $b\sqrt{1+h^2/c^2}$. Since $a>b$, the major axis is along $x$, and the foci are at $(\pm\sqrt{(a^2-b^2)(1+h^2/c^2)},0,h)$. Put $z=h$ and eliminate $h$: the foci satisfy $x^2/(a^2-b^2)-z^2/c^2=1$, a hyperbola in the $xOz$ plane with real axis the $x$-axis.}

\ASK{单叶双曲面与双叶双曲面在截线形状上有何本质差别？}{What is the essential difference between the sections of the two hyperboloids?}{单叶双曲面被 $z=h$ 截出的永远是椭圆（$h$ 任意），并且这些椭圆无界增大；被 $y=h$ 截时，$|h|<b$ 得实轴沿 $x$ 的双曲线、$|h|>b$ 得实轴沿 $z$ 的双曲线、$|h|=b$ 退化为两条相交直线——这些直线正是 §4.7 的直母线。双叶双曲面被 $z=h$ 截时，$|h|<c$ 完全无截口，$|h|=c$ 截得一点，$|h|>c$ 才得到椭圆；它不能由一族直线生成，也没有实直母线。}{Sections of the one-sheeted hyperboloid by $z=h$ are always ellipses, growing without bound; sections by $y=h$ are hyperbolas (with real axis along $x$ for $|h|<b$, along $z$ for $|h|>b$) and degenerate into two intersecting lines for $|h|=b$, which are its generators (§4.7). The two-sheeted hyperboloid has no section for $|h|<c$, a single point for $|h|=c$, and ellipses only for $|h|>c$; it is not ruled and has no real generators.}

\section{\label{sec:6}抛物面 | Paraboloids}
\APS{抛物面是无心二次曲面：椭圆抛物面与双曲抛物面都没有对称中心，却都有一条对称轴和两个互相垂直的对称面。两者都由抛物线``平移''生成：椭圆抛物面的两条主抛物线开口同向，双曲抛物面开口相反，所以后者形如马鞍。本节用主抛物线与截线展示两种曲面的结构，并算一个球面与旋转抛物面的交线。}{Paraboloids are the quadrics without a centre: the elliptic and hyperbolic paraboloids have no centre of symmetry but do have one axis and two mutually perpendicular planes of symmetry. Both are swept by translating a parabola: the two principal parabolas of the elliptic paraboloid open in the same direction, those of the hyperbolic one in opposite directions, which produces the saddle shape. The section uses principal parabolas and sections to display both structures, and computes one intersection with a sphere.}

\subsection{椭圆抛物面 | Elliptic paraboloid}
\begin{definition}[4.6.1 椭圆抛物面 / elliptic paraboloid]\label{def:ell-parab}
在直角坐标系下，由方程
\[
\frac{x^2}{a^2}+\frac{y^2}{b^2}=2z\tag{4.6-1}
\]
所表示的曲面叫做\textbf{椭圆抛物面 / elliptic paraboloid}，方程 (4.6-1) 叫做椭圆抛物面的标准方程，其中 $a,b$ 是任意的正常数。
\EN{In rectangular coordinates the surface (4.6-1), with positive $a,b$, is the elliptic paraboloid and (4.6-1) its standard equation.}
\end{definition}

\APS{显然椭圆抛物面 (4.6-1) 关于 $xOz$ 与 $yOz$ 坐标面对称，也关于 $z$ 轴对称，但是它没有对称中心；它与对称轴交于点 $(0,0,0)$，这点叫做椭圆抛物面 (4.6-1) 的\textbf{顶点 / vertex}。从方程 (4.6-1) 知
\[
z=\frac12\Bigl(\frac{x^2}{a^2}+\frac{y^2}{b^2}\Bigr)\ge 0,
\]
所以曲面全部在 $xOy$ 平面的一侧，即在 $z\ge 0$ 的一侧。}{It is symmetric about the planes $xOz$, $yOz$ and the $z$-axis, but has no centre; the origin is its vertex. Since $z\ge0$, the surface lies in the half-space $z\ge0$.}

\APS{用坐标面 $y=0$ 及 $x=0$ 截割曲面 (4.6-1)，分别得抛物线
\[
\begin{cases}x^2=2a^2z,\\ y=0\end{cases}\tag{1}
\]
与
\[
\begin{cases}y^2=2b^2z,\\ x=0.\end{cases}\tag{2}
\]
这两条抛物线叫做椭圆抛物面 (4.6-1) 的\textbf{主抛物线 / principal parabolas}，它们有着共同的轴与相同的开口方向，即开口方向都与 $z$ 轴的正向一致。用坐标面 $xOy$ 来截曲面 (4.6-1) 只得一点 $(0,0,0)$，但用平行于 $xOy$ 面的平面 $z=h\ (h>0)$ 来截曲面 (4.6-1)，截线总是椭圆
\[
\begin{cases}\dfrac{x^2}{2a^2h}+\dfrac{y^2}{2b^2h}=1,\\ z=h,\end{cases}\tag{3}
\]
这个椭圆的两对顶点分别为 $(\pm a\sqrt{2h},0,h)$、$(0,\pm b\sqrt{2h},h)$，它们分别在主抛物线 (1) 与 (2) 上。因此椭圆抛物面 (4.6-1) 可以看成是由一个椭圆的变动（大小位置都改变）而产生的，这个椭圆在变动中保持所在平面平行于 $xOy$ 面，且两对顶点分别在抛物线 (1) 与 (2) 上滑动。}{$y=0$ and $x=0$ give the principal parabolas (1), (2), with the same axis and opening towards positive $z$. The section by $z=h>0$ is the ellipse (3) with vertices $(\pm a\sqrt{2h},0,h)$, $(0,\pm b\sqrt{2h},h)$ sliding along (1) and (2).}

\APS{用平行于 $xOz$ 面的平面 $y=t$ 截割椭圆抛物面 (4.6-1) 得抛物线
\[
\begin{cases}x^2=2a^2\Bigl(z-\dfrac{t^2}{2b^2}\Bigr),\\ y=t.\end{cases}\tag{4}
\]
显然抛物线 (4) 与主抛物线 (1) 全等（两个有相同焦参数 $p$ 的抛物线是全等的），且它所在的平面平行于主抛物线 (1) 所在的平面并有相同的开口方向；此外，抛物线 (4) 的顶点 $\bigl(0,t,\frac{t^2}{2b^2}\bigr)$ 位于主抛物线 (2) 上。于是得到下面的结论：如果取两个这样的抛物线，它们所在的平面互相垂直，它们的顶点和轴都重合，而且两抛物线有相同的开口方向，让其中一条抛物线平行于自己（即与抛物线所在的平面平行）且使其顶点在另一个抛物线上滑动，那么前一抛物线的运动轨迹便是一个椭圆抛物面。}{Sections by $y=t$ are parabolas (4) congruent to (1), with the same opening direction, whose vertices $(0,t,t^2/2b^2)$ lie on (2). Hence: take two orthogonal parabolas with a common vertex and axis and the same opening direction; translate one parallel to itself with its vertex sliding on the other; it sweeps an elliptic paraboloid.}

\APS{在方程 (4.6-1) 中，如果 $a=b$，那么方程变为 (4.3-5)，即
\[
x^2+y^2=2a^2z,
\]
这时截线 (3) 为一圆，曲面就成为旋转抛物面。}{If $a=b$, (4.6-1) becomes $x^2+y^2=2a^2z$ and the sections (3) are circles, giving the paraboloid of revolution.}

\subsection{双曲抛物面 | Hyperbolic paraboloid}
\begin{definition}[4.6.2 双曲抛物面 / hyperbolic paraboloid]\label{def:hyp-parab}
在直角坐标系下，由方程
\[
\frac{x^2}{a^2}-\frac{y^2}{b^2}=2z\tag{4.6-2}
\]
所表示的曲面叫做\textbf{双曲抛物面 / hyperbolic paraboloid}，方程 (4.6-2) 叫做双曲抛物面的标准方程，其中 $a,b$ 为任意的正常数。
\EN{In rectangular coordinates the surface (4.6-2), with positive $a,b$, is the hyperbolic paraboloid.}
\end{definition}

\APS{显然曲面 (4.6-2) 关于 $xOz$ 面、$yOz$ 面与 $z$ 轴对称，但是它没有对称中心。用坐标面 $z=0$ 去截割曲面 (4.6-2)，就得
\[
\begin{cases}\dfrac{x^2}{a^2}-\dfrac{y^2}{b^2}=0,\\ z=0,\end{cases}\tag{5}
\]
这是一对相交于原点的直线
\[
\begin{cases}\dfrac{x}{a}-\dfrac{y}{b}=0,\\ z=0\end{cases}\quad\text{与}\quad
\begin{cases}\dfrac{x}{a}+\dfrac{y}{b}=0,\\ z=0.\end{cases}\tag{5$'$}
\]
其次用坐标面 $y=0$ 与 $x=0$ 来截割曲面 (4.6-2)，分别得两抛物线
\[
\begin{cases}x^2=2a^2z,\\ y=0\end{cases}\tag{6}
\]
与
\[
\begin{cases}y^2=-2b^2z,\\ x=0.\end{cases}\tag{7}
\]
这两条抛物线叫做双曲抛物面 (4.6-2) 的\textbf{主抛物线}，它们所在的平面相互垂直，有相同的顶点与对称轴，但两抛物线的开口方向不同：抛物线 (6) 沿 $z$ 轴方向开口，而抛物线 (7) 的开口方向却与 $z$ 轴方向相反。}{It is symmetric about $xOz$, $yOz$ and the $z$-axis but has no centre. The plane $z=0$ cuts the pair of lines (5'), and $y=0$, $x=0$ cut the principal parabolas (6), (7), which are orthogonal, share vertex and axis, but open in opposite directions.}

\APS{如果用平行于 $xOy$ 面的平面 $z=h\ (h\ne 0)$ 来截割曲面 (4.6-2)，截线总是双曲线
\[
\begin{cases}\dfrac{x^2}{2a^2h}-\dfrac{y^2}{2b^2h}=1,\\ z=h.\end{cases}\tag{8}
\]
当 $h>0$ 时，双曲线 (8) 的实轴与 $x$ 轴平行，虚轴与 $y$ 轴平行，顶点 $(\pm a\sqrt{2h},0,h)$ 在主抛物线 (6) 上；当 $h<0$ 时，双曲线 (8) 的实轴与 $y$ 轴平行，虚轴与 $x$ 轴平行，顶点 $(0,\pm b\sqrt{-2h},h)$ 在主抛物线 (7) 上。因此，曲面 (4.6-2) 被 $xOy$ 平面分割成上下两部分，上半部沿 $x$ 轴的两个方向上升，下半部沿 $y$ 轴的两个方向下降，曲面的大体形状像一只马鞍子，所以双曲抛物面 (4.6-2) 也叫做\textbf{马鞍曲面 / saddle surface}。}{Sections $z=h\ne0$ are hyperbolas (8): for $h>0$ with real axis along $x$ and vertices on (6), for $h<0$ with real axis along $y$ and vertices on (7). Hence the surface rises along the $x$-direction and falls along the $y$-direction: a saddle.}

\APS{用平行于 $xOz$ 面的一组平行平面 $y=t$ 来截割曲面 (4.6-2)，截线为抛物线
\[
\begin{cases}x^2=2a^2\Bigl(z+\dfrac{t^2}{2b^2}\Bigr),\\ y=t.\end{cases}\tag{9}
\]
不论 $t$ 取怎样的实数，所截得的抛物线 (9) 总与主抛物线 (6) 全等，且所在平面平行于这个主抛物线所在的平面；它的顶点 $\bigl(0,t,-\frac{t^2}{2b^2}\bigr)$ 在另一主抛物线 (7) 上。于是得到结论：如果取两个这样的抛物线，它们所在的平面互相垂直，有公共的顶点与轴，而两抛物线的开口方向相反，让其中的一个抛物线平行于自己且使其顶点在另一抛物线上滑动，那么前一抛物线的运动轨迹便是一个双曲抛物面。椭圆抛物面与双曲抛物面统称为抛物面，它们都没有对称中心，所以又叫做\textbf{无心二次曲面 / quadrics without centre}。}{Sections $y=t$ are parabolas (9) congruent to (6) with vertices $(0,t,-t^2/2b^2)$ on (7); translating the one along the other with opposite openings sweeps the hyperbolic paraboloid. Elliptic and hyperbolic paraboloids are the paraboloids, and both are centraless quadrics.}

\EX{（例 1，书中第 118--119 页）作出球面 $x^2+y^2+z^2=8$ 与旋转抛物面 $x^2+y^2=2z$ 的交线。}{（Example 1, pp.~118--119）Describe the intersection of the sphere $x^2+y^2+z^2=8$ with the paraboloid of revolution $x^2+y^2=2z$.}
两曲面的交线为
\[
\begin{cases}x^2+y^2+z^2=8,\\ x^2+y^2=2z.\end{cases}\tag{1--2}
\]
把 (2) 代入 (1) 得
\[
z^2+2z-8=0,
\]
即
\[
(z+4)(z-2)=0,
\]
所以 $z=-4$ 或 $z=2$。由 (2) 知 $z\ge 0$，所以取 $z=2$，因此交线方程可改写为
\[
\begin{cases}x^2+y^2=4,\\ z=2,\end{cases}
\]
这是平面 $z=2$ 上的一个圆，圆心为 $(0,0,2)$，半径为 $2$。
\EN{The intersection is given by (1)--(2). Substituting $x^2+y^2=2z$ gives $z^2+2z-8=0$, i.e.\ $(z+4)(z-2)=0$; since $z\ge0$ we take $z=2$, so the intersection is the circle $x^2+y^2=4$ in the plane $z=2$, with centre $(0,0,2)$ and radius $2$.}

\ASK{椭圆抛物面与双曲抛物面的主抛物线各有什么特点？怎样从截线判断抛物面的类型？}{What distinguishes the principal parabolas of the two paraboloids? How do sections identify the type?}{椭圆抛物面 $x^2/a^2+y^2/b^2=2z$ 的两条主抛物线 $x^2=2a^2z$、$y^2=2b^2z$ 开口同向（都朝 $z$ 正方向），水平截线恒为椭圆；双曲抛物面 $x^2/a^2-y^2/b^2=2z$ 的两条主抛物线 $x^2=2a^2z$、$y^2=-2b^2z$ 开口反向，水平截线为双曲线，$z=0$ 截出两条相交直线。因此判断时看水平截线：椭圆 $\to$ 椭圆抛物面（$a=b$ 时是旋转抛物面），双曲线（且 $z=0$ 退化为一对相交直线）$\to$ 双曲抛物面。}{For the elliptic paraboloid the two principal parabolas open in the same direction and all horizontal sections are ellipses; for the hyperbolic one the principal parabolas open in opposite directions, horizontal sections are hyperbolas, and $z=0$ degenerates into two intersecting lines. So elliptic sections mean elliptic paraboloid (a paraboloid of revolution if $a=b$), while hyperbolic sections plus the degenerate $z=0$ mean hyperbolic paraboloid.}

\section{\label{sec:7}直母线：单叶双曲面与双曲抛物面 | Generators of ruled quadrics}
\APS{本节是全章的高潮：单叶双曲面与双曲抛物面不仅``含有''直线，而且整个曲面可以由一族直线扫出——它们是直纹曲面。方法极其漂亮：把单叶双曲面的方程化成两个一次式的乘积相等 $(A)(B)=(C)(D)$，再令 $A = uC$、$B = \frac1u D$，就得到一族直线；交换写法又得另一族。最后用四点共面的行列式判据证明异族直母线共面（双曲抛物面则相交）、同族直母线异面。}{This section is the climax of the chapter: the one-sheeted hyperboloid and the hyperbolic paraboloid not only contain lines but are swept by families of lines; they are ruled surfaces. The device is elegant: write the equation as an equality of products of linear factors $(A)(B)=(C)(D)$, then set $A=uC$, $B=u^{-1}D$ to get one family, and interchange to get the other. A four-point coplanarity determinant then proves that generators from different families are coplanar (and meet, for the hyperbolic paraboloid) while those from the same family are skew.}

\APS{柱面与锥面都可以由一族直线所生成，这种由一族直线所生成的曲面叫做\textbf{直纹曲面 / ruled surface}，而生成曲面的那族直线叫做这曲面的\textbf{一族直母线 / a family of generators}。下面证明单叶双曲面与双曲抛物面也都是直纹曲面。}{A surface swept by a family of lines is a ruled surface, and the family is a family of generators. Both the one-sheeted hyperboloid and the hyperbolic paraboloid are ruled.}

\subsection{单叶双曲面的两族直母线 | Two families on the one-sheeted hyperboloid}
\APS{首先考察单叶双曲面
\[
\frac{x^2}{a^2}+\frac{y^2}{b^2}-\frac{z^2}{c^2}=1,\tag{1}
\]
其中 $a,b,c$ 为正的常数。把 (1) 改写为
\[
\frac{x^2}{a^2}-\frac{z^2}{c^2}=1-\frac{y^2}{b^2},
\]
或者
\[
\Bigl(\frac{x}{a}+\frac{z}{c}\Bigr)\Bigl(\frac{x}{a}-\frac{z}{c}\Bigr)=\Bigl(1+\frac{y}{b}\Bigr)\Bigl(1-\frac{y}{b}\Bigr).\tag{2}
\]
现在引进不等于零的参数 $u$，并考察由上式得来的方程组
\[
\begin{cases}\dfrac{x}{a}+\dfrac{z}{c}=u\Bigl(1+\dfrac{y}{b}\Bigr),\\[4pt] \dfrac{x}{a}-\dfrac{z}{c}=\dfrac1u\Bigl(1-\dfrac{y}{b}\Bigr),\end{cases}\tag{3}
\]
以及
\[
\begin{cases}\dfrac{x}{a}+\dfrac{z}{c}=0,\\ 1-\dfrac{y}{b}=0,\end{cases}\tag{4}
\]
\[
\begin{cases}\dfrac{x}{a}-\dfrac{z}{c}=0,\\ 1+\dfrac{y}{b}=0.\end{cases}\tag{4$'$}
\]
方程组 (4) 与 (4$'$) 实际上是 (3) 中当参数 $u\to 0$ 和 $u\to\infty$ 时的两种极限情形。显然不论 $u$ 取何值，(3) 以及 (4)、(4$'$) 都表示直线，把 (3)、(4)、(4$'$) 合起来组成的一族直线叫做 $u$ 族直线。}{Start from (1) and factor it as (2); introduce $u\ne0$ and consider the system (3), together with the limiting cases (4), (4$'$). Each member is a line, and the whole collection is the $u$-family.}

\APS{证明 $u$ 族直线生成曲面 (1)：容易知道 $u$ 族直线中的任何一条直线上的点都在曲面 (1) 上，这是因为当 $u\ne0$ 时，由 (3) 两边相乘即得 (2) 从而得 (1)，所以 (3) 所表示的直线上的点都在曲面 (1) 上；而满足 (4) 与 (4$'$) 的点显然满足 (2)，从而满足 (1)，因此直线 (4) 与 (4$'$) 上的点也都在曲面 (1) 上。反过来，设 $(x_0,y_0,z_0)$ 是曲面 (1) 上的点，从而有
\[
\Bigl(\frac{x_0}{a}+\frac{z_0}{c}\Bigr)\Bigl(\frac{x_0}{a}-\frac{z_0}{c}\Bigr)=\Bigl(1+\frac{y_0}{b}\Bigr)\Bigl(1-\frac{y_0}{b}\Bigr).\tag{5}
\]
显然 $1+\frac{y_0}{b}$ 与 $1-\frac{y_0}{b}$ 不能同时为零，因此不失一般性，假设 $1+\frac{y_0}{b}\ne0$。如果 $\frac{x_0}{a}-\frac{z_0}{c}\ne0$，那么取 $u$ 的值使得 $\frac{x_0}{a}+\frac{z_0}{c}=u\bigl(1+\frac{y_0}{b}\bigr)$，由 (5) 便得 $\frac{x_0}{a}-\frac{z_0}{c}=\frac1u\bigl(1-\frac{y_0}{b}\bigr)$，所以点 $(x_0,y_0,z_0)$ 在直线 (3) 上；如果 $\frac{x_0}{a}-\frac{z_0}{c}=0$，那么由 (5) 知必有 $1-\frac{y_0}{b}=0$，所以点 $(x_0,y_0,z_0)$ 在直线 (4) 上。因此曲面 (1) 上的任一点一定在 $u$ 族直线中的某一条直线上，即曲面 (1) 由 $u$ 族直线生成，$u$ 族直线是单叶双曲面 (1) 的一族直母线。}{Every point of a line of the $u$-family satisfies (1) (multiply the two equations of (3)), and the lines (4), (4$'$) satisfy (2) hence (1). Conversely, any point $(x_0,y_0,z_0)$ of (1) satisfies (5); since the two factors in $y_0$ cannot both vanish, either one solves for $u$ and finds the point on a line (3), or the degenerate case puts it on (4). So the $u$-family sweeps the surface.}

\APS{同样可以证明由直线
\[
\begin{cases}\dfrac{x}{a}+\dfrac{z}{c}=v\Bigl(1-\dfrac{y}{b}\Bigr),\\[4pt] \dfrac{x}{a}-\dfrac{z}{c}=\dfrac1v\Bigl(1+\dfrac{y}{b}\Bigr)\end{cases}\tag{6}
\]
（其中 $v$ 为不等于零的任意实数）与另两直线（相当于 (6) 中当 $v\to0$ 和 $v\to\infty$ 的情形）
\[
\begin{cases}\dfrac{x}{a}+\dfrac{z}{c}=0,\\ 1+\dfrac{y}{b}=0,\end{cases}\tag{7}
\]
\[
\begin{cases}\dfrac{x}{a}-\dfrac{z}{c}=0,\\ 1-\dfrac{y}{b}=0\end{cases}\tag{7$'$}
\]
合在一起组成的直线族是单叶双曲面 (1) 的另一族直母线，称为 $v$ 族直母线。}{The same argument with (6) and its limits (7), (7$'$) gives the $v$-family.}

\begin{corollary}\label{cor:onegen}
对于单叶双曲面上的点，两族直母线中各有一条直母线通过这点。
\EN{Through every point of the one-sheeted hyperboloid there passes exactly one generator from each family.}
\end{corollary}

\APS{为了避免取极限，常把单叶双曲面 (1) 的 $u$ 族直母线写成
\[
\begin{cases}w\Bigl(\dfrac{x}{a}+\dfrac{z}{c}\Bigr)=u\Bigl(1+\dfrac{y}{b}\Bigr),\\[6pt] u\Bigl(\dfrac{x}{a}-\dfrac{z}{c}\Bigr)=w\Bigl(1-\dfrac{y}{b}\Bigr),\end{cases}\tag{4.7-1}
\]
其中 $u,w$ 不同时为零：当 $u\ne0,w\ne0$ 时各式除以 $w$ 即化为 (3)；当 $u=0$ 时化为 (4)；当 $w=0$ 时化为 (4$'$)。而 $v$ 族直母线写成
\[
\begin{cases}t\Bigl(\dfrac{x}{a}+\dfrac{z}{c}\Bigr)=v\Bigl(1-\dfrac{y}{b}\Bigr),\\[6pt] v\Bigl(\dfrac{x}{a}-\dfrac{z}{c}\Bigr)=t\Bigl(1+\dfrac{y}{b}\Bigr),\end{cases}\tag{4.7-2}
\]
其中 $v,t$ 不同时为零。必须指出，(4.7-1) 与 (4.7-2) 中的直线分别只依赖于 $u:w$ 与 $v:t$ 的值。}{To avoid limits, the $u$-family is written as (4.7-1) with $u,w$ not both zero, and the $v$-family as (4.7-2) with $v,t$ not both zero; each line depends only on the ratio $u:w$ resp.\ $v:t$.}

\subsection{双曲抛物面的两族直母线 | Two families on the hyperbolic paraboloid}
\APS{对于双曲抛物面
\[
\frac{x^2}{a^2}-\frac{y^2}{b^2}=2z,
\]
同样可以证明它也有两族直母线，它们的方程分别是
\[
\begin{cases}\dfrac{x}{a}+\dfrac{y}{b}=2u,\\[6pt] u\Bigl(\dfrac{x}{a}-\dfrac{y}{b}\Bigr)=z,\end{cases}\tag{4.7-3}
\]
与
\[
\begin{cases}\dfrac{x}{a}-\dfrac{y}{b}=2v,\\[6pt] v\Bigl(\dfrac{x}{a}+\dfrac{y}{b}\Bigr)=z.\end{cases}\tag{4.7-4}
\]
并且也有下面的推论：对于双曲抛物面上的点，两族直母线中各有一条直母线通过这一点。}{The hyperbolic paraboloid also has two families, given by (4.7-3) and (4.7-4), and through each of its points passes one generator of each family.}

\APS{注记（书中脚注）：对于双曲抛物面的直母线族方程不用双参数，这是因为若将 (4.7-3) 改写为双参数形式
\[
\begin{cases}w\Bigl(\dfrac{x}{a}+\dfrac{y}{b}\Bigr)=2u,\\ u\Bigl(\dfrac{x}{a}-\dfrac{y}{b}\Bigr)=wz,\end{cases}
\]
那么当 $w=0$ 时必有 $u=0$，所以如果要写成这种形式，必须附加条件 $w\ne0$，但这样它实质上就是 (4.7-3) 了。}{Remark: no two-parameter form is used for the hyperbolic paraboloid, since setting $w=0$ in the two-parameter form forces $u=0$; the form is thus equivalent to (4.7-3) only when $w\ne0$.}

\begin{theorem}[4.7.1 异族直母线的位置关系 / generators of different families]\label{thm:gens-cross}
单叶双曲面上异族的任意两直母线必共面，而双曲抛物面上异族的任意两直母线必相交。
\EN{On the one-sheeted hyperboloid any two generators from different families are coplanar; on the hyperbolic paraboloid any two generators from different families intersect.}
\end{theorem}

\begin{proof}[证明 / Proof]
只证前半部分。由 (4.7-1) 与 (4.7-2) 的四个方程的系数和常数项所组成的行列式为
\[
\begin{vmatrix}
\dfrac{w}{a} & -\dfrac{u}{b} & \dfrac{w}{c} & -u\\[6pt]
\dfrac{u}{a} & \dfrac{w}{b} & -\dfrac{u}{c} & -w\\[6pt]
\dfrac{t}{a} & \dfrac{v}{b} & \dfrac{t}{c} & -v\\[6pt]
\dfrac{v}{a} & -\dfrac{t}{b} & -\dfrac{v}{c} & -t
\end{vmatrix}
=-\frac{1}{abc}
\begin{vmatrix}
w & -u & w & u\\
u & w & -u & w\\
t & v & t & v\\
v & -t & -v & t
\end{vmatrix}
=-\frac{4}{abc}(uwvt-wuvt)=0.
\]
根据 §3.8 的例 3（四条直线共面的行列式判据）知道这两直线一定是共面的，所以单叶双曲面上异族的两直母线必共面。
\EN{Only the first claim is proved. The determinant of the coefficients and constant terms of the four equations (4.7-1), (4.7-2) reduces to $-(4/abc)(uwvt-wuvt)=0$, so by the coplanarity criterion of §3.8, Example 3, the two lines are coplanar.}
\end{proof}

\begin{theorem}[4.7.2 同族直母线的位置关系 / generators of the same family]\label{thm:gens-skew}
单叶双曲面或双曲抛物面上同族的任意两直母线总是异面直线，而且双曲抛物面同族的全体直母线平行于同一平面。
\EN{On the one-sheeted hyperboloid or the hyperbolic paraboloid any two generators from the same family are skew lines; moreover all generators of one family on the hyperbolic paraboloid are parallel to one fixed plane.}
\end{theorem}

\APS{定理 4.7.2 的证明书中留给读者：对单叶双曲面，同族两条直线的方向向量成比例而位置不同，故平行；又同族直线不能相交（否则交点处两族各有一条直母线，两直线重合，与族的定义矛盾），故为异面。对双曲抛物面，同族直线（如 (4.7-3) 族）的方向向量都平行于平面 $x/a+y/b=0$ 的方向，故全体平行于同一平面。}{The book leaves the proof as an exercise: on the one-sheeted hyperboloid generators of one family have parallel direction vectors, hence cannot meet (otherwise two generators would coincide), so they are skew; on the hyperbolic paraboloid all lines of (4.7-3) are parallel to one and the same plane.}

\EX{（书中第 124 页）求过单叶双曲面 $\dfrac{x^2}{9}+\dfrac{y^2}{4}-\dfrac{z^2}{16}=1$ 上的点 $(6,2,8)$ 的直母线方程。}{（p.~124）Find the generators through the point $(6,2,8)$ of the hyperboloid $x^2/9+y^2/4-z^2/16=1$.}
该单叶双曲面的两族直母线方程是
\[
\begin{cases}w\Bigl(\dfrac{x}{3}+\dfrac{z}{4}\Bigr)=u\Bigl(1+\dfrac{y}{2}\Bigr),\\[6pt] u\Bigl(\dfrac{x}{3}-\dfrac{z}{4}\Bigr)=w\Bigl(1-\dfrac{y}{2}\Bigr),\end{cases}
\quad\text{与}\quad
\begin{cases}t\Bigl(\dfrac{x}{3}+\dfrac{z}{4}\Bigr)=v\Bigl(1-\dfrac{y}{2}\Bigr),\\[6pt] v\Bigl(\dfrac{x}{3}-\dfrac{z}{4}\Bigr)=t\Bigl(1+\dfrac{y}{2}\Bigr).\end{cases}
\]
把点 $(6,2,8)$ 分别代入上面两组方程，求得
\[
w:u=1:2\quad\text{与}\quad t=0,
\]
代入直母线族方程，得过 $(6,2,8)$ 的两条直母线分别为
\[
\begin{cases}4x-12y+3z-24=0,\\ 4x+3y-3z-6=0,\end{cases}
\quad\text{与}\quad
\begin{cases}y-2=0,\\ 4x-3z=0.\end{cases}
\]
\EN{The two families are displayed above. Substituting $(6,2,8)$ into them gives $w:u=1:2$ for the first family and $t=0$ for the second, hence the two generators displayed.}

\ASK{为什么单叶双曲面能由直线生成，而双叶双曲面不能？}{Why is the one-sheeted hyperboloid ruled while the two-sheeted one is not?}{核心是方程能否分解为两个一次式之积相等：单叶的 $x^2/a^2+y^2/b^2-z^2/c^2=1$ 可写成 $(x/a+z/c)(x/a-z/c)=(1+y/b)(1-y/b)$，于是令两侧成比例即得直线族；双叶的 $x^2/a^2+y^2/b^2-z^2/c^2=-1$ 即 $z^2/c^2-x^2/a^2-y^2/b^2=1$，左端不能分解为两个实一次式之积（它是正定二次型之差且右边为 $1$ 时对应双曲面，但 $z^2-x^2-y^2$ 的零锥是实锥面，双叶双曲面在锥面内部，实直母线不存在）。另外从截线看，单叶双曲面 $|h|=b$ 的退化为两条直线，就是直母线的痕迹。}{The key is whether the equation can be written as an equality of products of two linear factors: for one sheet $(x/a+z/c)(x/a-z/c)=(1+y/b)(1-y/b)$ gives lines by proportionality; for two sheets the corresponding factorization has no real linear factors, and the surface lies inside the real cone $z^2/c^2=x^2/a^2+y^2/b^2$, so no real generators exist. The degenerate sections of the one-sheeted surface at $|h|=b$ are traces of generators.}

\section{\label{sec:8}第四章口试问答 | Chapter 4 interview questions}
\APS{下面把第 4 章压缩成一张二次曲面总表与四个口试题。表中的``截线''指用平行于坐标面的平面去截的截口类型，``直母线''指是否有实直母线。}{Below is a summary table of the quadrics of Chapter 4 and four interview questions. In the table ``sections'' means traces by planes parallel to coordinate planes, and ``generators'' asks for real rulings.}

\begin{center}\small
\begin{tabular}{@{}>{\raggedright\arraybackslash}p{2.6cm}>{\raggedright\arraybackslash}p{3.7cm}>{\raggedright\arraybackslash}p{3.2cm}>{\raggedright\arraybackslash}p{2.2cm}>{\raggedright\arraybackslash}p{2.5cm}@{}}
\toprule
曲面 Surface & 标准方程 Equation & 水平截线 $z=h$ & 对称性 & 直母线 \\
\midrule
椭球面 ellipsoid & $\frac{x^2}{a^2}+\frac{y^2}{b^2}+\frac{z^2}{c^2}=1$ & 椭圆（$|h|<c$） & 三平面、三轴、中心 & 无 \\
单叶双曲面 one sheet & $\frac{x^2}{a^2}+\frac{y^2}{b^2}-\frac{z^2}{c^2}=1$ & 椭圆（任意 $h$） & 三平面、三轴、中心 & 两族 \\
双叶双曲面 two sheets & $\frac{x^2}{a^2}+\frac{y^2}{b^2}-\frac{z^2}{c^2}=-1$ & 无（$|h|<c$）、椭圆（$|h|>c$） & 三平面、三轴、中心 & 无 \\
椭圆抛物面 elliptic paraboloid & $\frac{x^2}{a^2}+\frac{y^2}{b^2}=2z$ & 椭圆（$h>0$） & 两平面、$z$ 轴、无心 & 无 \\
双曲抛物面 hyperbolic paraboloid & $\frac{x^2}{a^2}-\frac{y^2}{b^2}=2z$ & 双曲线（$h\ne0$）、两直线（$h=0$） & 两平面、$z$ 轴、无心 & 两族 \\
柱面 cylinder & $F(x,y)=0$ & 直线（母线） & 随 $F$ 而定 & $F$ 的一次因子给出 \\
锥面 cone & 关于 $x,y,z$ 齐次 & 随次数而定 & 过原点 & 过顶点的直线 \\
\bottomrule
\end{tabular}
\end{center}

\ASK{从方程出发，怎样在最短时间内说出二次曲面的类型与主要特征？}{Given an equation, how do you identify a quadric and its main features quickly?}{先看平方项的符号与个数：三项全正或全负且右边为 $1$ 是椭球面；两正一负（右边 $1$）是单叶双曲面，一正两负或等价地右边 $-1$ 是双叶双曲面；只有一个平方项（另一坐标只出现一次）是抛物面：同号为椭圆抛物面，异号为双曲抛物面。再看缺项：只含两个坐标的是柱面；含 $xy,yz,zx$ 等交叉项时先转轴消去，再按上述判断。最后用平行截割法确认：椭球面的截口有界，单叶双曲面与双曲抛物面有直线截口，双叶双曲面有``无截口''的区间。}{Look first at the signs and number of squared terms: three positive (or all negative) with right side $1$ gives an ellipsoid; two positive one negative gives a one-sheeted hyperboloid, and the right side $-1$ gives two sheets; a single squared term plus one linear term gives a paraboloid, with equal signs elliptic and opposite signs hyperbolic. A missing variable means a cylinder. Cross terms are removed by rotation first. Finally confirm by sections: bounded sections for the ellipsoid, line sections for the ruled quadrics, and a range with no section for the two-sheeted hyperboloid.}

\ASK{平行截割法为什么能``决定''曲面形状？它的局限是什么？}{Why does the method of parallel sections determine the shape? What are its limits?}{若对每个 $h$，截口曲线的方程、形状与位置都已知，那么曲面就是这些截口曲线的并，故形状完全确定；而且截口族的变化规律（半轴如何随 $h$ 变）给出``母线''（变动椭圆）的运动规律。局限是：截口只给出``水平切片''，不能直接看出沿其他方向的弯曲（如双曲抛物面的鞍形需要再看平行于 $xOz$ 的截线），且当截口退化（点、直线）时需要单独分析，否则会误判。}{If for every $h$ the section is known in shape and position, the surface is the union of the sections, so it is determined; the way the semi-axes vary gives the motion of the sweeping ellipse. The limitation is that sections only show horizontal slices, not the bending in other directions (the saddle needs the sections parallel to $xOz$), and degenerate sections (points, lines) must be analysed separately.}

\ASK{单叶双曲面与双曲抛物面都是直纹曲面，它们的直母线结构有何不同？}{Both are ruled; how does the structure of their generators differ?}{单叶双曲面的两族直母线中，异族必共面（定理 4.7.1），同族异面；双曲抛物面两族中异族必相交，同族异面且全体平行于同一平面（定理 4.7.2）。差别来自方程：单叶双曲面的两族由 $(x/a\pm z/c)$ 与 $(1\pm y/b)$ 配对，双曲抛物面由 $(x/a\pm y/b)$ 与 $z$ 配对，后者使两族直线落在同一族平行平面中。}{On the one-sheeted hyperboloid generators from different families are coplanar and those from the same family are skew; on the hyperbolic paraboloid generators from different families intersect, while one family is skew yet all parallel to a fixed plane. The reason is the pairing: $(x/a\pm z/c)$ with $(1\pm y/b)$ for the hyperboloid, $(x/a\pm y/b)$ with $z$ for the paraboloid.}

\ASK{为什么说柱面、锥面、旋转曲面、二次曲面这四类``统一''在``曲线按规律运动''这一观点下？}{How do the four kinds of surfaces unify under ``a curve moving by a rule''?}{柱面是直线沿准线平移，锥面是直线绕顶点转动并沿准线滑动，旋转曲面是曲线绕轴旋转，二次曲面则可由一族变动椭圆或抛物线扫出（如椭球面由 $z=h$ 的椭圆族生成）。统一的求法都是：写出含参数的母线族方程（或纬圆族），再写出参数满足的约束条件，消去参数即得曲面方程。}{A cylinder is a line translated along a directrix, a cone a line through a fixed vertex sliding along a directrix, a surface of revolution a curve rotated about an axis, and a quadric is swept by a varying ellipse or parabola. The unified method: write the parameter-dependent family of generators (or latitude circles), add the constraint satisfied by the parameter, and eliminate it.}

\section{\label{sec:9}二次曲线的一般方程、记号与直线相关位置 | General conics and lines}
\APS{第 5 章的统一工具是``矩阵 + 二次型''。把二次曲线写成 $F(x,y)=0$ 后，一切几何概念都可以用 $F$ 的系数表达：方向用只含二次项的二次型 $\Phi(X,Y)$，点用两个``半梯度'' $F_1,F_2$，位置关系用行列式 $I_1,I_2,I_3,K_1$。本节先建立这套记号与恒等式 $F=xF_1+yF_2+F_3$，再用参数方程把直线代入二次曲线，由所得一元二次方程的判别式得到``两个实交点、两个重合实交点、一对共轭虚交点、一个实交点、无交点、整条直线在曲线上''六种可能。这一节是整个第 5 章的``语法''。}{The unifying tool of Chapter 5 is matrices plus quadratic forms. Once a conic is written $F(x,y)=0$, every geometric notion can be expressed through the coefficients: directions by the quadratic form $\Phi(X,Y)$ of the quadratic terms, points by the ``half-gradients'' $F_1,F_2$, and relative positions by the determinants $I_1,I_2,I_3,K_1$. This section sets up the notation and the identity $F=xF_1+yF_2+F_3$, then substitutes a parametrised line into $F$; the discriminant of the resulting quadratic gives the six possible positions: two real intersections, two coincident real intersections, a pair of conjugate imaginary intersections, one real intersection, no intersection, or the whole line contained in the conic. This section is the grammar of the chapter.}

\APS{在平面上，由二元二次方程
\[
F(x,y)\equiv a_{11}x^2+2a_{12}xy+a_{22}y^2+2a_{13}x+2a_{23}y+a_{33}=0\tag{1}
\]
所表示的曲线，叫做\textbf{二次曲线 / conic}。其中 $xy$ 项的系数写成 $2a_{12}$，是为了以后配平方与求偏导数的方便；以后总假设系数为实数，但允许点、方向为复的（见前文``预备：虚元素''）。}{A curve given by the real quadratic equation (1) is a conic. The coefficient of $xy$ is written $2a_{12}$ to simplify completing squares and differentiating; the coefficients are real, but points and directions may be complex.}

\subsection{记号、矩阵与恒等式 | Notation, matrices, identity}
\APS{为了方便起见，引进下面的一些记号：
\[
F(x,y)\equiv a_{11}x^2+2a_{12}xy+a_{22}y^2+2a_{13}x+2a_{23}y+a_{33},
\]
\[
F_1(x,y)\equiv a_{11}x+a_{12}y+a_{13},\qquad F_2(x,y)\equiv a_{12}x+a_{22}y+a_{23},\qquad F_3(x,y)\equiv a_{13}x+a_{23}y+a_{33},
\]
\[
\Phi(x,y)\equiv a_{11}x^2+2a_{12}xy+a_{22}y^2.
\]
这样容易验证，下面的恒等式成立：
\[
F(x,y)\equiv xF_1(x,y)+yF_2(x,y)+F_3(x,y),\tag{2}
\]
(1) 式也就可以写成
\[
F(x,y)\equiv xF_1(x,y)+yF_2(x,y)+F_3(x,y)=0.\tag{3}
\]
今后常用记号：$F_1=\frac12F_x$，$F_2=\frac12F_y$（便于记忆，借用偏导数记号），$F_3$ 是去掉二次项后剩下的一次与常数部分。}{Introduce $F_1,F_2,F_3,\Phi$ as displayed. Then the identity (2) holds, so (1) can be written in the form (3). Memorise $F_1=\tfrac12F_x$, $F_2=\tfrac12F_y$, and that $F_3$ is the linear-plus-constant part.}

\APS{把 $F(x,y)$ 的系数所排成的矩阵
\[
A=\begin{pmatrix}a_{11}&a_{12}&a_{13}\\ a_{12}&a_{22}&a_{23}\\ a_{13}&a_{23}&a_{33}\end{pmatrix}
\]
叫做二次曲线 (1) 的\textbf{矩阵 / matrix}（或称 $F(x,y)$ 的矩阵），而将 $\Phi(x,y)$ 的系数所排成的矩阵
\[
A^*=\begin{pmatrix}a_{11}&a_{12}\\ a_{12}&a_{22}\end{pmatrix}
\]
叫做 $\Phi(x,y)$ 的矩阵。显然，二次曲线 (1) 的矩阵 $A$ 的第一、第二与第三行（或列）的元素分别是 $F_1(x,y)$、$F_2(x,y)$、$F_3(x,y)$ 的系数。}{The symmetric matrix $A$ of all coefficients is the matrix of the conic, and $A^*$ that of $\Phi$. The rows (or columns) of $A$ are the coefficients of $F_1,F_2,F_3$.}

\APS{今后还常常要引用下面的几个符号：
\[
I_1=a_{11}+a_{22},\qquad
I_2={\begin{vmatrix}a_{11}&a_{12}\\ a_{12}&a_{22}\end{vmatrix}},\qquad
I_3={\begin{vmatrix}a_{11}&a_{12}&a_{13}\\ a_{12}&a_{22}&a_{23}\\ a_{13}&a_{23}&a_{33}\end{vmatrix}},
\]
\[
K_1={\begin{vmatrix}a_{11}&a_{13}\\ a_{13}&a_{33}\end{vmatrix}}+{\begin{vmatrix}a_{22}&a_{23}\\ a_{23}&a_{33}\end{vmatrix}}.
\]
这里的 $I_1$ 是矩阵 $A^*$ 的主对角元素的和，$I_2$ 是 $A^*$ 的行列式，$I_3$ 是 $A$ 的行列式；而 $K_1$ 的两项是 $I_1$ 的两项分别添加上两条``边''而成的两个二阶行列式，这添加上的两条``边''的元素是矩阵 $A$ 中的第三行与第三列的对应元素，也就是说用二阶行列式 $\begin{vmatrix}a_{11}&a_{13}\\ a_{13}&a_{33}\end{vmatrix}$、$\begin{vmatrix}a_{22}&a_{23}\\ a_{23}&a_{33}\end{vmatrix}$ 分别代替 $I_1$ 中的 $a_{11},a_{22}$，就由 $I_1$ 得到 $K_1$。第 5.7 节将证明 $I_1,I_2,I_3$ 在直角坐标变换下不变，而 $K_1$ 只在转轴下不变（半不变量）。}{Define $I_1,I_2,I_3,K_1$ as displayed: $I_1$ is the trace of $A^*$, $I_2=\det A^*$, $I_3=\det A$, and $K_1$ consists of the two $2\times2$ minors obtained by adding to the two terms of $I_1$ the corresponding entries of the third row and column of $A$. Section 5.7 proves that $I_1,I_2,I_3$ are invariants under rectangular coordinate changes, while $K_1$ is only invariant under rotation (a seminvariant).}

\subsection{二次曲线与直线的相关位置 | A conic and a line}
\APS{现在来讨论二次曲线 (1) 与过点 $(x_0,y_0)$ 且具有方向 $X:Y$ 的直线
\[
\begin{cases}x=x_0+Xt,\\ y=y_0+Yt\end{cases}\tag{2$'$}
\]
的交点。把 (2$'$) 代入 (1)，经过整理得关于 $t$ 的方程
\[
\Phi(X,Y)t^2+2\bigl[F_1(x_0,y_0)X+F_2(x_0,y_0)Y\bigr]t+F(x_0,y_0)=0,\tag{3}
\]
利用前面的记号，(3) 可写成
\[
\Phi(X,Y)t^2+2\bigl[XF_1(x_0,y_0)+YF_2(x_0,y_0)\bigr]t+F(x_0,y_0)=0.\tag{4}
\]
方程 (3) 或 (4) 可分以下几种情况来讨论。}{Substitute the line (2$'$) into (1) to get the quadratic (3)--(4) in $t$. Discuss the cases by the sign of $\Phi(X,Y)$ and the discriminant.}

\APS{1) $\Phi(X,Y)\ne0$。这时 (4) 是关于 $t$ 的二次方程，它的判别式为
\[
\Delta=\bigl[XF_1(x_0,y_0)+YF_2(x_0,y_0)\bigr]^2-\Phi(X,Y)F(x_0,y_0),\tag{5}
\]
这又可分三种情况：(i) $\Delta>0$，方程 (4) 有两个不等的实根 $t_1$ 与 $t_2$，代入 (2$'$) 便得直线 (2$'$) 与二次曲线 (1) 的两个不同的实交点；(ii) $\Delta=0$，方程 (4) 有两个相等的实根 $t_1=t_2$，这时直线 (2$'$) 与二次曲线 (1) 有两个相互重合的实交点；(iii) $\Delta<0$，方程 (4) 有两个共轭的虚根，这时直线 (2$'$) 与二次曲线交于两个共轭的虚点。}{If $\Phi(X,Y)\ne0$, the discriminant is (5): $\Delta>0$ gives two distinct real intersections, $\Delta=0$ two coincident real intersections, $\Delta<0$ a pair of conjugate imaginary intersections.}

\APS{2) $\Phi(X,Y)=0$。这时又可分三种情况：(i) $F_1(x_0,y_0)X+F_2(x_0,y_0)Y\ne0$，这时 (4) 是关于 $t$ 的一次方程，它有惟一的一个实根，所以直线与二次曲线有惟一的实交点；(ii) $F_1(x_0,y_0)X+F_2(x_0,y_0)Y=0$ 而 $F(x_0,y_0)\ne0$，这时 (4) 为矛盾方程，方程无解，所以直线与二次曲线没有交点；(iii) $F_1(x_0,y_0)X+F_2(x_0,y_0)Y=F(x_0,y_0)=0$，这时方程 (4) 成为一个恒等式，它能被任何值（实的或虚的）$t$ 所满足，所以直线 (2$'$) 上的一切点都是 (1) 与 (2$'$) 的公共点，也就是说直线 (2$'$) 全部在二次曲线上。}{If $\Phi(X,Y)=0$: a nonvanishing linear term gives exactly one real intersection; a vanishing linear term with $F\ne0$ gives none; if both vanish, the whole line lies on the conic.}

\EX{（补充算例，仿书中 §5.1 习题 2）求二次曲线 $x^2-2xy-3y^2-4x-6y+3=0$ 与直线 $5x-y-5=0$ 的交点。}{（Supplementary example, after Exercise 2 of §5.1）Find the intersections of the conic $x^2-2xy-3y^2-4x-6y+3=0$ with the line $5x-y-5=0$.}
直线的方向数可取 $X:Y=1:5$，且 $\Phi(1,5)=1-2\cdot1\cdot5-3\cdot25=-84\ne0$，故交点数由判别式决定。由直线方程得 $y=5x-5$，代入曲线方程：
\[
x^2-2x(5x-5)-3(5x-5)^2-4x-6(5x-5)+3=0.
\]
展开并合并同类项：
\[
(1-10-75)x^2+(10+150-4-30)x+(-75+30+3)=0,
\]
即
\[
-84x^2+126x-42=0\quad\Longleftrightarrow\quad 2x^2-3x+1=0\quad\Longleftrightarrow\quad (2x-1)(x-1)=0.
\]
所以 $x=1$ 或 $x=\frac12$，从而 $y=0$ 或 $y=-\frac52$。所求交点为 $(1,0)$ 与 $\bigl(\frac12,-\frac52\bigr)$，是两个不同的实交点。
\EN{Take $y=5x-5$; since $\Phi(1,5)=-84\ne0$ the number of intersections is decided by the discriminant. Substitution gives $-84x^2+126x-42=0$, i.e.\ $(2x-1)(x-1)=0$, so $x=1$ or $x=\tfrac12$ and the intersections are $(1,0)$ and $(\tfrac12,-\tfrac52)$.}

\ASK{为什么把直线写成参数方程而不是直接消元？判别式中的 $\Phi(X,Y)$ 起什么作用？}{Why parametrise the line instead of eliminating one variable? What is the role of $\Phi(X,Y)$?}{参数方程让交点的``重数''一目了然：$t$ 是直线上的坐标，方程的根就是交点，重根就是切点，且根的实、复性直接反映交点的实、复性；直接消元会把两种坐标的不对称性带进判别式。$\Phi(X,Y)$ 是二次齐次部分在方向 $(X,Y)$ 上的值，它决定方程 (3) 是否为二次：$\Phi\ne0$ 时弦有两个端点，$\Phi=0$ 时直线最多与曲线交于一点或整条在曲线上，所以 $\Phi=0$ 的方向就是下一节的``渐近方向''。}{The parametrisation makes multiplicities transparent: $t$ is a coordinate on the line, roots are intersections, double roots are points of tangency, and reality of roots reflects reality of intersections; elimination would destroy the symmetry. $\Phi(X,Y)$ is the quadratic part evaluated in the direction $(X,Y)$; it decides whether (3) is quadratic: if $\Phi\ne0$ a chord has two endpoints, while $\Phi=0$ means at most one intersection or the whole line, so $\Phi=0$ defines the ``asymptotic directions'' of the next section.}

\section{\label{sec:10}渐近方向、中心与渐近线 | Asymptotic directions, centres, asymptotes}
\APS{本节从``方向''与``点''两个方面拆解二次曲线。方向由 $\Phi(X,Y)$ 的零点（渐近方向）刻划，由此得到椭圆型、抛物型、双曲型三种类型；点由``是否为所有弦的中点''刻划（中心），由此得到中心曲线、无心曲线、线心曲线三种类型。渐近线是``过中心、以渐近方向为方向''的直线，它是双曲线最显著的几何特征。}{This section splits a conic into directions and points. Directions are controlled by the zeros of $\Phi(X,Y)$ (asymptotic directions), which yields the elliptic, parabolic and hyperbolic types; points are tested by ``being the midpoint of every chord'' (centre), which yields central, non-central and line-central conics. An asymptote is the line through the centre in an asymptotic direction, the most visible feature of a hyperbola.}

\subsection{二次曲线的渐近方向 | Asymptotic directions}
\APS{在 §5.1 中看到，二次曲线 (1) 和具有方向 $X:Y$ 的直线 (2$'$) 当满足条件
\[
\Phi(X,Y)=a_{11}X^2+2a_{12}XY+a_{22}Y^2=0\tag{3}
\]
时，或者只有一个实交点，或者没有交点，或者直线 (2$'$) 全部在二次曲线 (1) 上，成为二次曲线的组成部分。}{Recall: when $\Phi(X,Y)=0$, the line has at most one intersection, or lies on the conic.}

\begin{definition}[5.2.1 渐近方向 / asymptotic direction]\label{def:asymdir}
满足条件 $\Phi(X,Y)=0$ 的方向 $X:Y$ 叫做二次曲线 (1) 的\textbf{渐近方向 / asymptotic direction}，否则叫做\textbf{非渐近方向 / non-asymptotic direction}。
\EN{A direction $X:Y$ with $\Phi(X,Y)=0$ is an asymptotic direction of (1); otherwise it is non-asymptotic.}
\end{definition}

\APS{因为二次曲线 (1) 的二次项系数不能全为零，所以渐近方向 $X:Y$ 所满足的 (3) 总有确定的解。如果 $a_{22}\ne0$，把 (3) 改写成
\[
a_{11}\Bigl(\frac{X}{Y}\Bigr)^2+2a_{12}\frac{X}{Y}+a_{22}=0
\]
得
\[
\frac{X}{Y}=\frac{-a_{12}\pm\sqrt{a_{12}^2-a_{11}a_{22}}}{a_{22}}=\frac{-a_{12}\pm\sqrt{-I_2}}{a_{22}}.
\]
如果 $a_{22}=0$，把 (3) 改写成
\[
X(a_{11}X+2a_{12}Y)=0,
\]
得 $X:Y=1:0$ 或 $X:Y=-2a_{12}:a_{11}$。所以，当且仅当 $I_2>0$ 时，二次曲线 (1) 的渐近方向是一对共轭的虚方向；$I_2=0$ 时，(1) 有一个实渐近方向；$I_2<0$ 时，(1) 有两个实渐近方向。因此二次曲线的渐近方向最多有两个，显然二次曲线的非渐近方向有无数多个。}{If $a_{22}\ne0$ then $X/Y=(-a_{12}\pm\sqrt{-I_2})/a_{22}$; if $a_{22}=0$ then $X:Y=1:0$ or $-2a_{12}:a_{11}$. Hence $I_2>0$ gives a pair of conjugate imaginary asymptotic directions, $I_2=0$ one real one, $I_2<0$ two real ones; there are at most two asymptotic directions and infinitely many non-asymptotic ones.}

\begin{definition}[5.2.2 按渐近方向的类型 / types by asymptotic directions]\label{def:conictype}
没有实渐近方向的二次曲线叫做\textbf{椭圆型的 / elliptic}，有一个实渐近方向的二次曲线叫做\textbf{抛物型的 / parabolic}，有两个实渐近方向的二次曲线叫做\textbf{双曲型的 / hyperbolic}。
\EN{A conic with no real asymptotic direction is elliptic, with one real asymptotic direction parabolic, with two real asymptotic directions hyperbolic.}
\end{definition}

\APS{因此二次曲线 (1) 按其渐近方向可以分为三种类型，即
\[
\text{椭圆型：}I_2>0;\qquad \text{抛物型：}I_2=0;\qquad \text{双曲型：}I_2<0.
\]}{Thus the type is read off from $I_2$: elliptic if $I_2>0$, parabolic if $I_2=0$, hyperbolic if $I_2<0$.}

\subsection{二次曲线的中心与渐近线 | Centres and asymptotes}
\APS{在 §5.1 中又看到，当直线 (2$'$) 的方向 $X:Y$ 为二次曲线 (1) 的非渐近方向时，直线与二次曲线总交于两个点（两不同实的、两重合实的或一对共轭虚的），我们把由这两点决定的线段叫做二次曲线的\textbf{弦 / chord}。}{A non-asymptotic line meets the conic in two points (distinct real, coincident real, or conjugate imaginary); the segment joining them is a chord.}

\begin{definition}[5.2.3 中心 / centre]\label{def:centre}
如果点 $C$ 是二次曲线的通过它的所有弦的中点（因而 $C$ 是二次曲线的对称中心），那么点 $C$ 叫做二次曲线的\textbf{中心 / centre}。
\EN{A point $C$ that is the midpoint of every chord passing through it (hence the centre of symmetry) is the centre of the conic.}
\end{definition}

\APS{根据这个定义，当点 $(x_0,y_0)$ 为二次曲线 (1) 的中心时，过 $(x_0,y_0)$ 以任意非渐近方向 $X:Y$ 为方向的直线与 (1) 交于两点 $M_1,M_2$，点 $(x_0,y_0)$ 就是弦 $M_1M_2$ 的中点。因此把直线代入 (1) 得
\[
\Phi(X,Y)t^2+2\bigl[XF_1(x_0,y_0)+YF_2(x_0,y_0)\bigr]t+F(x_0,y_0)=0.\tag{4}
\]
因为 $X:Y$ 为任意非渐近方向，所以 (4) 式是关于 $X,Y$ 的恒等式，从而有
\[
F_1(x_0,y_0)=0,\qquad F_2(x_0,y_0)=0.
\]
反过来，适合上面两式的点 $(x_0,y_0)$，显然是二次曲线的中心。}{If $(x_0,y_0)$ is a centre, the two roots in $t$ are opposite for every non-asymptotic direction, so (4) is an identity in $X,Y$, forcing $F_1=F_2=0$; conversely these equations make $(x_0,y_0)$ a centre.}

\begin{theorem}[5.2.1 中心坐标 / coordinates of the centre]\label{thm:centre}
点 $C(x_0,y_0)$ 是二次曲线 (1) 的中心，其充要条件是
\[
\begin{cases}F_1(x_0,y_0)\equiv a_{11}x_0+a_{12}y_0+a_{13}=0,\\ F_2(x_0,y_0)\equiv a_{12}x_0+a_{22}y_0+a_{23}=0.\end{cases}\tag{5.2-1}
\]
\EN{$C(x_0,y_0)$ is a centre of (1) if and only if (5.2-1) holds.}
\end{theorem}

\begin{corollary}[原点为中心 / origin as centre]\label{cor:origin}
坐标原点是二次曲线的中心，其充要条件是曲线方程里不含 $x$ 与 $y$ 的一次项。
\EN{The origin is the centre if and only if the equation contains no linear terms in $x$ and $y$.}
\end{corollary}

\APS{所以，二次曲线 (1) 的中心坐标由下列方程组决定：
\[
\begin{cases}F_1(x,y)\equiv a_{11}x+a_{12}y+a_{13}=0,\\ F_2(x,y)\equiv a_{12}x+a_{22}y+a_{23}=0.\end{cases}\tag{5.2-2}
\]
如果 $I_2=\begin{vmatrix}a_{11}&a_{12}\\ a_{12}&a_{22}\end{vmatrix}\ne0$，那么 (5.2-2) 有惟一解，这时二次曲线 (1) 将有惟一中心，(5.2-2) 的解即为中心的坐标。如果 $I_2=0$，即 $\frac{a_{11}}{a_{12}}=\frac{a_{12}}{a_{22}}$，那么当 $\frac{a_{11}}{a_{12}}=\frac{a_{12}}{a_{22}}\ne\frac{a_{13}}{a_{23}}$ 时，(5.2-2) 无解，二次曲线 (1) 没有中心；而当 $\frac{a_{11}}{a_{12}}=\frac{a_{12}}{a_{22}}=\frac{a_{13}}{a_{23}}$ 时，(5.2-2) 有无数多解，这时直线 $a_{11}x+a_{12}y+a_{13}=0$（或 $a_{12}x+a_{22}y+a_{23}=0$）上的所有点都是二次曲线 (1) 的中心，这条直线叫做\textbf{中心直线 / line of centres}。}{The centre solves (5.2-2). If $I_2\ne0$ there is a unique centre; if $I_2=0$ with $a_{11}/a_{12}=a_{12}/a_{22}\ne a_{13}/a_{23}$ there is no centre; if all three ratios are equal, every point of the line $a_{11}x+a_{12}y+a_{13}=0$ is a centre, the line of centres.}

\begin{definition}[5.2.4 按中心的分类 / types by centres]\label{def:centretype}
有惟一中心的二次曲线叫做\textbf{中心二次曲线 / central conic}，没有中心的二次曲线叫做\textbf{无心二次曲线 / non-central conic}，有一条中心直线的二次曲线叫做\textbf{线心二次曲线 / line-central conic}，无心二次曲线与线心二次曲线统称为\textbf{非中心二次曲线 / non-central conics}。
\EN{A conic with a unique centre is central; with no centre, non-central; with a line of centres, line-central; the last two are called non-central.}
\end{definition}

\APS{根据这个定义与 (5.2-2)，二次曲线 (1) 按其中心的分类为：
\[
\text{1) 中心曲线：}I_2\ne0;\qquad
\text{2) 非中心曲线：}I_2=0,\ \text{即}\ \frac{a_{11}}{a_{12}}=\frac{a_{12}}{a_{22}},
\]
其中 $1^\circ$ \textbf{无心曲线}：$\frac{a_{11}}{a_{12}}=\frac{a_{12}}{a_{22}}\ne\frac{a_{13}}{a_{23}}$；$2^\circ$ \textbf{线心曲线}：$\frac{a_{11}}{a_{12}}=\frac{a_{12}}{a_{22}}=\frac{a_{13}}{a_{23}}$。从按渐近方向与按中心的两种初步分类中容易看出，椭圆型曲线与双曲型曲线都是中心曲线，而抛物型曲线是非中心曲线，它包括无心曲线与线心曲线。}{Classification: central if $I_2\ne0$; non-central if $I_2=0$, i.e.\ $a_{11}/a_{12}=a_{12}/a_{22}$; among these, non-central (no centre) if the constant ratios differ, line-central if all three agree. Elliptic and hyperbolic conics are central; parabolic conics are non-central and split into the two kinds.}

\begin{definition}[5.2.5 渐近线 / asymptote]\label{def:asymptote}
通过二次曲线的中心，而且以渐近方向为方向的直线叫做这二次曲线的\textbf{渐近线 / asymptote}。
\EN{The line through the centre of a conic in an asymptotic direction is an asymptote.}
\end{definition}

\APS{显然，椭圆型曲线只有两条虚渐近线而无实渐近线，双曲型曲线有两条实渐近线，抛物型曲线中的无心曲线无渐近线，而线心曲线有一条实渐近线，就是它的中心直线。}{An elliptic conic has two imaginary asymptotes and no real one; a hyperbolic conic has two real asymptotes; a non-central parabolic conic has none; a line-central conic has one real asymptote, its line of centres.}

\begin{theorem}[5.2.2 渐近线与曲线的位置关系 / asymptotes and the conic]\label{thm:asymptote}
二次曲线的渐近线与这二次曲线或者没有交点，或者整条直线在这二次曲线上，成为二次曲线的组成部分。
\EN{An asymptote of a conic either misses the conic or lies entirely on it.}
\end{theorem}

\begin{proof}[证明 / Proof]
设直线 (2$'$) 是二次曲线 (1) 的渐近线，这里 $(x_0,y_0)$ 为二次曲线的中心，$X:Y$ 为二次曲线的渐近方向，那么有 $F_1(x_0,y_0)=F_2(x_0,y_0)=0$ 与 $\Phi(X,Y)=0$。因此根据 §5.1 中直线与二次曲线相交情况的讨论，当点 $(x_0,y_0)$ 不在二次曲线 (1) 上，即 $F(x_0,y_0)\ne0$ 时，渐近线与二次曲线没有交点；当点 $(x_0,y_0)$ 在二次曲线 (1) 上，即 $F(x_0,y_0)=0$ 时，渐近线全部在二次曲线上，成为二次曲线的组成部分。
\EN{For an asymptote the centre satisfies $F_1=F_2=0$ and the direction satisfies $\Phi=0$. By §5.1, if $F(x_0,y_0)\ne0$ there is no intersection, and if $F(x_0,y_0)=0$ the whole line lies on the conic.}
\end{proof}

\EX{（补充算例，仿书中 §5.2 习题）求二次曲线 $x^2-2xy-3y^2-4x-6y+3=0$ 的渐近方向、中心与渐近线，并指出它的类型。}{（Supplementary example, after Exercise of §5.2）Find the asymptotic directions, centre and asymptotes of $x^2-2xy-3y^2-4x-6y+3=0$, and identify its type.}
这里 $a_{11}=1$，$a_{12}=-1$，$a_{22}=-3$，$a_{13}=-2$，$a_{23}=-3$，$a_{33}=3$。
\[
I_2=\begin{vmatrix}1&-1\\ -1&-3\end{vmatrix}=-3-1=-4<0,
\]
所以曲线是\textbf{双曲型}，有两个实渐近方向。由 $\Phi(X,Y)=X^2-2XY-3Y^2=(X-3Y)(X+Y)=0$ 得渐近方向 $X:Y=1:3$ 与 $X:Y=1:(-1)$。中心由 $F_1=x-y-2=0$，$F_2=-x-3y-3=0$ 决定，解得 $x=\frac34$，$y=-\frac54$，即中心 $C\bigl(\frac34,-\frac54\bigr)$。过 $C$ 沿方向 $1:3$ 的渐近线为
\[
\frac{x-\frac34}{1}=\frac{y+\frac54}{3}\quad\Longleftrightarrow\quad 6x-2y-7=0;
\]
过 $C$ 沿方向 $1:(-1)$ 的渐近线为
\[
\frac{x-\frac34}{1}=\frac{y+\frac54}{-1}\quad\Longleftrightarrow\quad 2x+2y+1=0.
\]
\EN{Here $I_2=-4<0$, so the conic is hyperbolic. $\Phi=(X-3Y)(X+Y)$ gives asymptotic directions $1:3$ and $1:-1$; the centre solves $x-y-2=0$, $-x-3y-3=0$, i.e.\ $(\tfrac34,-\tfrac54)$, and the two asymptotes through it are $6x-2y-7=0$ and $2x+2y+1=0$.}

\EX{（补充算例）讨论二次曲线 $x^2-2xy+y^2+2x-2y-3=0$ 的类型、中心与渐近线。}{（Supplementary example）Discuss the type, centre and asymptotes of $x^2-2xy+y^2+2x-2y-3=0$.}
这里 $a_{11}=1$，$a_{12}=-1$，$a_{22}=1$，$a_{13}=1$，$a_{23}=-1$，$a_{33}=-3$，故 $I_2=\begin{vmatrix}1&-1\\ -1&1\end{vmatrix}=0$，曲线是\textbf{抛物型}。又 $\frac{a_{11}}{a_{12}}=\frac{a_{12}}{a_{22}}=\frac{a_{13}}{a_{23}}=-1$，故它是\textbf{线心曲线}，有中心直线
\[
x-y+1=0.
\]
渐近方向由 $\Phi(X,Y)=(X-Y)^2=0$ 给出，只有 $X:Y=1:1$；线心曲线只有一条渐近线，就是它的中心直线 $x-y+1=0$（事实上方程可分解为 $(x-y+3)(x-y-1)=0$，是两条平行直线，渐近线恰是与之平行的一条）。
\EN{$I_2=0$ and all three ratios equal, so the conic is parabolic and line-central with line of centres $x-y+1=0$. The only asymptotic direction is $1:1$; the single asymptote is the line of centres. Indeed the equation factors as $(x-y+3)(x-y-1)=0$, two parallel lines.}

\ASK{渐近方向、渐近线、中心三者之间是什么逻辑关系？}{What is the logical relation between asymptotic directions, asymptotes and centres?}{顺序是：先用 $\Phi(X,Y)=0$ 定义渐近方向（纯代数条件），再由 $F_1=F_2=0$ 定义中心（``所有弦中点''的几何条件），最后把二者结合：过中心、以渐近方向为方向作直线，就是渐近线。所以中心曲线（$I_2\ne0$）才有完整的渐近线理论；无心曲线没有中心，谈不上渐近线；线心曲线的中心退化为一条直线，渐近线就是这条中心直线。}{The order is: $\Phi(X,Y)=0$ defines asymptotic directions (algebraic), $F_1=F_2=0$ defines the centre (the midpoint condition), and an asymptote is the line through the centre in an asymptotic direction. Hence a full theory of asymptotes needs a central conic; a non-central parabolic conic without centre has none, while a line-central conic has its line of centres as its single asymptote.}

\section{\label{sec:11}二次曲线的切线 | Tangents to a conic}
\APS{切线是``与曲线交于两个重合点的直线''。用 §5.1 的判别式立刻可以翻译成代数条件：把直线代入 $F$，所得二次方程的判别式为零。对曲线上的点，条件简化为一个一次方程 (5.3-2)，于是得到切线方向与切线方程 (5.3-3)；再把 (5.3-3) 恒等变形，得到对称美观、便于记忆的切线方程 (5.3-4)——它把 $x^2$ 换成 $x_0x$、$xy$ 换成 $\frac12(x_0y+xy_0)$ 等。若 $F_1=F_2=0$，切线方向不确定，这样的点是奇点。}{A tangent is a line meeting the conic in two coincident points. By the discriminant of §5.1 this becomes: the quadratic in $t$ has a double root. For a point on the conic the condition collapses to the linear equation (5.3-2), which gives the direction of the tangent and the equation (5.3-3); rewriting it symmetrically gives the memorable form (5.3-4), obtained by replacing $x^2$ by $x_0x$, $xy$ by $\frac12(x_0y+xy_0)$, and so on. If $F_1=F_2=0$ the direction is undetermined: such a point is singular.}

\begin{definition}[5.3.1 切线 / tangent]\label{def:tangent}
如果直线与二次曲线相交于相互重合的两个点，那么这条直线就叫做二次曲线的\textbf{切线 / tangent}，这个重合的交点叫做\textbf{切点 / point of tangency}；如果直线全部在二次曲线上，我们也称它为二次曲线的切线，直线上的每一个点都可以看做切点。
\EN{A line meeting a conic in two coincident points is a tangent, and the common point is the point of tangency; if the whole line lies on the conic it is also called a tangent, and every point of it is a point of tangency.}
\end{definition}

\APS{现在来求经过二次曲线
\[
F(x,y)\equiv a_{11}x^2+2a_{12}xy+a_{22}y^2+2a_{13}x+2a_{23}y+a_{33}=0\tag{1}
\]
上的点 $(x_0,y_0)$ 的切线方程。因为通过 $(x_0,y_0)$ 的直线总可写成
\[
\begin{cases}x=x_0+Xt,\\ y=y_0+Yt,\end{cases}\tag{2}
\]
那么根据 §5.1 的讨论，容易知道直线 (2) 成为二次曲线 (1) 的切线的条件，当 $\Phi(X,Y)\ne0$ 时，是
\[
\bigl[XF_1(x_0,y_0)+YF_2(x_0,y_0)\bigr]^2-\Phi(X,Y)F(x_0,y_0)=0.\tag{5.3-1}
\]
因为 $(x_0,y_0)$ 在 (1) 上，所以 $F(x_0,y_0)=0$，因而 (5.3-1) 可以化为
\[
XF_1(x_0,y_0)+YF_2(x_0,y_0)=0.\tag{5.3-2}
\]
当 $\Phi(X,Y)=0$ 时，直线 (2) 成为二次曲线 (1) 的切线的条件除了 $F(x_0,y_0)=0$ 外，惟一的条件仍然是 (5.3-2)。}{For the line (2) through a point of the conic, tangency means (5.3-1); since $F(x_0,y_0)=0$ this becomes the linear condition (5.3-2), which also covers the case $\Phi=0$.}

\APS{如果 $F_1(x_0,y_0)$ 与 $F_2(x_0,y_0)$ 不全为零，那么由 (5.3-2) 得
\[
X:Y=F_2(x_0,y_0):\bigl(-F_1(x_0,y_0)\bigr),
\]
因此过 $(x_0,y_0)$ 的切线方程为
\[
(x-x_0)F_1(x_0,y_0)+(y-y_0)F_2(x_0,y_0)=0,\tag{5.3-3}
\]
或写成
\[
\frac{x-x_0}{F_2(x_0,y_0)}=\frac{y-y_0}{-F_1(x_0,y_0)}.
\]
如果 $F_1(x_0,y_0)=F_2(x_0,y_0)=0$，那么 (5.3-2) 变为恒等式，切线的方向 $X:Y$ 不能惟一地被确定，从而切线不确定；这时通过 $(x_0,y_0)$ 的任何直线都和二次曲线 (1) 相交于相互重合的两点，我们把这样的直线也看成是二次曲线 (1) 的切线。}{If $(F_1,F_2)\ne(0,0)$, then $X:Y=F_2:-F_1$ and the tangent is (5.3-3). If $F_1=F_2=0$ the condition is an identity, every line through the point counts as a tangent, and the tangent is not unique.}

\begin{definition}[5.3.2 奇点与正则点 / singular and regular points]\label{def:singular}
二次曲线 (1) 上满足条件 $F_1(x_0,y_0)=F_2(x_0,y_0)=0$ 的点 $(x_0,y_0)$ 叫做二次曲线的\textbf{奇异点 / singular point}，简称\textbf{奇点}；二次曲线的非奇异点叫做二次曲线的\textbf{正则点 / regular point}。
\EN{A point of the conic at which $F_1=F_2=0$ is a singular point; a non-singular point is regular.}
\end{definition}

\begin{theorem}[5.3.1 切线方程 / equation of the tangent]\label{thm:tangent}
如果 $(x_0,y_0)$ 是二次曲线 (1) 的正则点，那么通过 $(x_0,y_0)$ 的切线方程是 (5.3-3)，$(x_0,y_0)$ 是它的切点。如果 $(x_0,y_0)$ 是二次曲线 (1) 的奇异点，那么通过 $(x_0,y_0)$ 的切线不确定，或者说通过点 $(x_0,y_0)$ 的每一条直线都是二次曲线 (1) 的切线。
\EN{At a regular point the tangent is (5.3-3) with the point as point of tangency; at a singular point every line through it is a tangent.}
\end{theorem}

\begin{corollary}[对称形式的切线方程 / symmetric form]\label{cor:tangent-sym}
如果 $(x_0,y_0)$ 是二次曲线 (1) 的正则点，那么通过 $(x_0,y_0)$ 的切线方程是
\[
a_{11}x_0x+a_{12}(x_0y+xy_0)+a_{22}y_0y+a_{13}(x+x_0)+a_{23}(y+y_0)+a_{33}=0.\tag{5.3-4}
\]
\EN{At a regular point $(x_0,y_0)$ the tangent is (5.3-4).}
\end{corollary}

\begin{proof}[证明 / Proof]
把 (5.3-3) 改写为
\[
xF_1(x_0,y_0)+yF_2(x_0,y_0)-\bigl[x_0F_1(x_0,y_0)+y_0F_2(x_0,y_0)\bigr]=0,
\]
再根据恒等式 $F=xF_1+yF_2+F_3$ 与 $F(x_0,y_0)=0$，上式又可写为
\[
x(a_{11}x_0+a_{12}y_0+a_{13})+y(a_{12}x_0+a_{22}y_0+a_{23})+(a_{13}x_0+a_{23}y_0+a_{33})=0,
\]
展开合并即得 (5.3-4)。
\EN{Rewrite (5.3-3) using the identity $F=xF_1+yF_2+F_3$ and $F(x_0,y_0)=0$, then collect terms to obtain (5.3-4).}
\end{proof}

\APS{为便于记忆公式 (5.3-4)，记忆的方法是：在原方程 (1) 中，把
\[
x^2\ \text{写成}\ x_0x,\qquad 2xy\ \text{写成}\ x_0y+xy_0,\qquad y^2\ \text{写成}\ y_0y,\qquad 2x\ \text{写成}\ x+x_0,\qquad 2y\ \text{写成}\ y+y_0,
\]
即每一项中一个 $x$（或 $y$）用 $x_0$（或 $y_0$）代入后再写成对称形式，就得出 (5.3-4)。}{Mnemonic for (5.3-4): replace $x^2$ by $x_0x$, $2xy$ by $x_0y+xy_0$, $y^2$ by $y_0y$, $2x$ by $x+x_0$, $2y$ by $y+y_0$.}

\EX{（例 1，书中第 136 页）求二次曲线 $x^2-xy+y^2+2x-4y-3=0$ 在点 $(2,1)$ 的切线方程。}{（Example 1, p.~136）Find the tangent to $x^2-xy+y^2+2x-4y-3=0$ at $(2,1)$.}
\textbf{解法一}：因为 $F(2,1)=4-2+1+4-4-3=0$，且
\[
F_1(2,1)=2-\frac12+1=\frac52\ne0,\qquad F_2(2,1)=-1+1-2=-2\ne0,
\]
所以 $(2,1)$ 是二次曲线上的正则点。由 (5.3-3) 得在点 $(2,1)$ 的切线方程为
\[
\frac52(x-2)-2(y-1)=0,
\]
即
\[
5x-4y-6=0.
\]
\textbf{解法二}：因为 $(2,1)$ 是曲线上的正则点，所以直接利用 (5.3-4)：这里 $a_{11}=1$，$a_{12}=-\frac12$，$a_{22}=1$，$a_{13}=1$，$a_{23}=-2$，$a_{33}=-3$，$x_0=2$，$y_0=1$，得
\[
2x-\frac12(2y+x)+y+(x+2)-2(y+1)-3=0,
\]
即
\[
5x-4y-6=0.
\]
\EN{Both methods give the same tangent. Method 1 uses (5.3-3) with $F_1(2,1)=\tfrac52$, $F_2(2,1)=-2$; method 2 uses the symmetric form (5.3-4) directly. The tangent is $5x-4y-6=0$.}

\EX{（例 2，书中第 136--137 页）求二次曲线 $x^2-xy+y^2-1=0$ 通过点 $(0,2)$ 的切线方程。}{（Example 2, pp.~136--137）Find the tangents to $x^2-xy+y^2-1=0$ passing through $(0,2)$.}
\textbf{解法一}：因为 $F(0,2)=3\ne0$，所以点 $(0,2)$ 不在曲线上，不能直接应用公式 (5.3-3) 或 (5.3-4)。过点 $(0,2)$ 的直线可以写成
\[
\begin{cases}x=Xt,\\ y=2+Yt,\end{cases}
\]
其中 $X:Y$ 为直线的方向数，$t$ 为参数。又因为
\[
F_1(0,2)=-1,\qquad F_2(0,2)=2,
\]
所以根据直线与二次曲线的相切条件 (5.3-1) 得
\[
\bigl(-X+2Y\bigr)^2=3\bigl(X^2-XY+Y^2\bigr),
\]
化简得
\[
2X^2+XY-Y^2=0,
\]
即 $(2X-Y)(X+Y)=0$，从而有
\[
X:Y=1:2\quad\text{或}\quad X:Y=1:(-1).
\]
再由过点 $(0,2)$ 的直线方程得两条切线
\[
2x-y+2=0\quad\text{与}\quad x+y-2=0.
\]
这两直线的方向分别为 $1:2$ 与 $1:(-1)$，显然它们都不是已知二次曲线的渐近方向，所以这两直线就是所求的过点 $(0,2)$ 的切线。\\[2pt]
\textbf{解法二}：设过 $(0,2)$ 的切线与已知二次曲线相切于点 $(x_0,y_0)$，那么切线方程为 (5.3-4)：
\[
x_0x-\frac12(x_0y+xy_0)+y_0y-1=0.\tag{3}
\]
因为它通过 $(0,2)$，所以将 $(0,2)$ 代入化简得
\[
x_0-2y_0+1=0.\tag{4}
\]
另一方面，点 $(x_0,y_0)$ 在曲线上，所以又有
\[
x_0^2-x_0y_0+y_0^2-1=0.\tag{5}
\]
联立解 (4)、(5)：由 (4) 得 $x_0=2y_0-1$，代入 (5) 得 $3y_0^2-3y_0=0$，所以 $y_0=0$ 或 $y_0=1$，切点坐标为 $(-1,0)$ 与 $(1,1)$。把切点坐标代入 (3) 得所求的切线方程为
\[
2x-y+2=0\quad\text{与}\quad x+y-2=0.
\]
\EN{Method 1: the point is not on the conic; write the line through it as $x=Xt$, $y=2+Yt$. With $F_1(0,2)=-1$, $F_2(0,2)=2$, the tangency condition (5.3-1) gives $(-X+2Y)^2=3(X^2-XY+Y^2)$, i.e.\ $(2X-Y)(X+Y)=0$, hence directions $1:2$ and $1:-1$ and tangents $2x-y+2=0$, $x+y-2=0$. Method 2: let $(x_0,y_0)$ be the point of tangency; its tangent (5.3-4) passes through $(0,2)$, giving $x_0-2y_0+1=0$, while $(x_0,y_0)$ satisfies the conic. Solving gives tangency points $(-1,0)$ and $(1,1)$, hence the same two tangents.}

\ASK{为什么曲线上的奇点处切线不确定？怎样快速判断一个点是奇点？}{Why is the tangent undetermined at a singular point? How do you spot a singular point?}{因为在正则点处 $F_1,F_2$ 不全为零，切线方向被方程 $XF_1+YF_2=0$ 唯一确定；在奇点处 $F_1=F_2=0$，该方程对任何 $X:Y$ 都成立，切线的方向条件失效，所以任何过该点的直线都与曲线``交于两个重合点''（在代数重数意义下），切线不确定。快速判断：解 $F_1=F_2=0$ 得候选点，再检查它是否满足 $F=0$；满足则为奇点。奇点通常出现在曲线退化（分解为两条直线且相交）的情形，例如 $(x-y+3)(x-y-1)=0$ 没有奇点，而 $x^2-y^2=0$ 的原点是奇点。}{At a regular point $(F_1,F_2)\ne(0,0)$ and the direction is fixed by $XF_1+YF_2=0$; at a singular point both vanish, the condition is void for every direction, and every line through the point meets the conic with multiplicity two, so the tangent is undetermined. To find singular points, solve $F_1=F_2=0$ and test $F=0$; they occur typically on degenerate conics consisting of two intersecting lines, e.g.\ $x^2-y^2=0$ has a singular point at the origin.}

\section{\label{sec:12}二次曲线的直径 | Diameters of a conic}
\APS{弦的中点轨迹是解析几何中最漂亮的``二次变一次''现象：一族平行弦的中点竟然总在一条直线上，这条直线就是直径。证明只用到韦达定理：两交点参数 $t_1,t_2$ 满足 $t_1+t_2=0$（中点对应参数 $0$）。由此得到直径方程 $XF_1+YF_2=0$；反过来，给定直径方向可以求出与它共轭的弦方向，于是有``共轭方向''与``共轭直径''的概念。定理 5.4.2 说明直径的走向由曲线类型决定：中心曲线的所有直径过中心，无心曲线的直径平行于渐近方向，线心曲线只有一条直径。}{The locus of midpoints of parallel chords is the most elegant ``quadratic to linear'' phenomenon in analytic geometry: the midpoints of a family of parallel chords always lie on a line, the diameter. The proof uses only Viète's formula: the two parameters $t_1,t_2$ of the intersections satisfy $t_1+t_2=0$ (the midpoint corresponds to parameter $0$). This gives $XF_1+YF_2=0$; conversely a diameter direction determines the direction of its conjugate chords, leading to conjugate directions and conjugate diameters. Theorem 5.4.2 says the behaviour depends on the type: all diameters of a central conic pass through the centre, diameters of a non-central conic are parallel to the asymptotic direction, and a line-central conic has only one diameter.}

\begin{theorem}[5.4.1 平行弦中点轨迹 / midpoints of parallel chords]\label{thm:diameter}
二次曲线的一族平行弦的中点轨迹是一条直线。
\EN{The midpoints of a family of parallel chords of a conic lie on a straight line.}
\end{theorem}

\begin{proof}[证明 / Proof]
设 $X:Y$ 是二次曲线的一个非渐近方向，即 $\Phi(X,Y)\ne0$，而 $(x_0,y_0)$ 是平行于方向 $X:Y$ 的弦的中点，那么过 $(x_0,y_0)$ 的弦为
\[
\begin{cases}x=x_0+Xt,\\ y=y_0+Yt.\end{cases}
\]
它与二次曲线 $F(x,y)=0$ 的两交点（即弦的两端点）由下列二次方程
\[
\Phi(X,Y)t^2+2\bigl[XF_1(x_0,y_0)+YF_2(x_0,y_0)\bigr]t+F(x_0,y_0)=0
\]
的两根 $t_1$ 与 $t_2$ 所决定。因为 $(x_0,y_0)$ 为弦的中点，所以 $t_1+t_2=0$，从而有
\[
XF_1(x_0,y_0)+YF_2(x_0,y_0)=0,
\]
这就是说平行于方向 $X:Y$ 的弦的中点 $(x_0,y_0)$ 的坐标满足方程
\[
XF_1(x,y)+YF_2(x,y)=0,\tag{5.4-1}
\]
或
\[
X(a_{11}x+a_{12}y+a_{13})+Y(a_{12}x+a_{22}y+a_{23})=0,\tag{5.4-2}
\]
或
\[
(a_{11}X+a_{12}Y)x+(a_{12}X+a_{22}Y)y+(a_{13}X+a_{23}Y)=0.\tag{5.4-3}
\]
反过来，如果点 $(x_0,y_0)$ 满足方程 (5.4-1)（或 (5.4-2)、(5.4-3)），那么方程中 $t$ 的两个根将有绝对值相等而符号相反的值，点 $(x_0,y_0)$ 就是具有方向 $X:Y$ 的弦的中点，因此方程 (5.4-1) 或 (5.4-2) 或 (5.4-3) 为一族平行于某一非渐近方向 $X:Y$ 的弦的中点轨迹方程。方程 (5.4-3) 的一次项系数不能全为零，因为当 $a_{11}X+a_{12}Y=a_{12}X+a_{22}Y=0$ 时将有
\[
\Phi(X,Y)=(a_{11}X+a_{12}Y)X+(a_{12}X+a_{22}Y)Y=0,
\]
这与 $X:Y$ 是非渐近方向的假设矛盾，所以 (5.4-3) 是一个二元一次方程，它是一条直线，于是定理得到了证明。
\EN{For a non-asymptotic direction, the two roots in $t$ have sum zero for the midpoint, which is exactly (5.4-1)--(5.4-3); conversely these equations force $t_1+t_2=0$. The coefficients of the linear equation cannot both vanish, since that would make $\Phi(X,Y)=0$. Hence the locus is a line.}
\end{proof}

\begin{definition}[5.4.1 直径 / diameter]\label{def:diameter}
二次曲线的平行弦中点的轨迹叫做这个二次曲线的\textbf{直径 / diameter}，它所对应的平行弦叫做\textbf{共轭于这条直径的共轭弦 / the chords conjugate to it}，而直径也叫做\textbf{共轭于平行弦方向的直径 / the diameter conjugate to that direction}。
\EN{The locus of midpoints of a family of parallel chords is called a diameter; the chords are conjugate to it, and the diameter is conjugate to their direction.}
\end{definition}

\begin{corollary}[斜率形式的直径方程 / diameter in slope form]\label{cor:diameter}
如果二次曲线的一族平行弦的斜率为 $k$，那么共轭于这族平行弦的直径方程是
\[
(a_{11}k+a_{12})x+(a_{12}k+a_{22})y+(a_{13}k+a_{23})=0.\tag{5.4-4}
\]
\EN{For chords of slope $k$ the conjugate diameter is (5.4-4).}
\end{corollary}

\APS{从方程 (5.4-3) 容易看出：当 $I_2\ne0$ 即二次曲线为中心曲线时，(5.4-1) 或 (5.4-4) 构成一个中心直线束，其中心就是二次曲线的中心；当 $I_2=0$ 即二次曲线为非中心曲线时，全部直径属于一个平行直线束，它的方向为二次曲线的渐近方向 $X:Y$（当 $a_{22}\ne0$ 时即为 $-a_{12}:a_{22}$）。当二次曲线为线心曲线时，这时二次曲线只有一条直径，它的方程是
\[
a_{11}x+a_{12}y+a_{13}=0\quad(\text{或 } a_{12}x+a_{22}y+a_{23}=0),
\]
即线心二次曲线的中心直线。因此我们有：}{If $I_2\ne0$ the diameters form a pencil through the centre; if $I_2=0$ they form a parallel pencil in the asymptotic direction; for a line-central conic there is a single diameter, namely the line of centres.}

\begin{theorem}[5.4.2 直径与曲线类型 / diameters and types]\label{thm:diamtype}
中心二次曲线的直径通过曲线的中心，无心二次曲线的直径平行于曲线的渐近方向，线心二次曲线的直径只有一条，就是曲线的中心直线。
\EN{Diameters of a central conic pass through its centre; those of a non-central conic are parallel to its asymptotic direction; a line-central conic has a single diameter, its line of centres.}
\end{theorem}

\EX{（例 1，书中第 139--140 页）求椭圆 $\dfrac{x^2}{a^2}+\dfrac{y^2}{b^2}=1$（或双曲线 $\dfrac{x^2}{a^2}-\dfrac{y^2}{b^2}=1$）的直径。}{（Example 1, pp.~139--140）Find the diameters of the ellipse $x^2/a^2+y^2/b^2=1$ and of the hyperbola $x^2/a^2-y^2/b^2=1$.}
对椭圆，$a_{11}=\frac1{a^2}$，$a_{12}=0$，$a_{22}=\frac1{b^2}$，$a_{13}=a_{23}=0$，$F_1=\frac{x}{a^2}$，$F_2=\frac{y}{b^2}$。根据 (5.4-1)，共轭于非渐近方向 $X:Y$ 的直径方程是
\[
\frac{X}{a^2}x+\frac{Y}{b^2}y=0.
\]
显然，直径通过曲线的中心 $(0,0)$。对双曲线同理得 $\dfrac{X}{a^2}x-\dfrac{Y}{b^2}y=0$，也通过中心。
\EN{For the ellipse the diameter conjugate to $X:Y$ is $Xx/a^2+Yy/b^2=0$, passing through the origin; the hyperbola gives $Xx/a^2-Yy/b^2=0$.}

\EX{（例 2，书中第 140 页）求抛物线 $y^2=2px$ 的直径。}{（Example 2, p.~140）Find the diameters of the parabola $y^2=2px$.}
这里 $F(x,y)\equiv y^2-2px$，$a_{11}=0$，$a_{12}=0$，$a_{22}=1$，$a_{13}=-p$，$a_{23}=0$，$F_1=-p$，$F_2=y$。所以共轭于非渐近方向 $X:Y$ 的直径为
\[
-pX+Yy=0,
\]
即
\[
y=\frac{pX}{Y}.
\]
所以抛物线 $y^2=2px$ 的直径平行于它的渐近方向 $1:0$（即平行于 $x$ 轴），这与定理 5.4.2 一致。
\EN{With $F=y^2-2px$ we get $F_1=-p$, $F_2=y$, so the diameter is $y=pX/Y$, a line parallel to the asymptotic direction $1:0$, as Theorem 5.4.2 predicts.}

\EX{（例 3，书中第 140 页）求二次曲线 $x^2-2xy+y^2+2x-2y-3=0$ 的共轭于非渐近方向 $X:Y$ 的直径。}{（Example 3, p.~140）Find the diameter of $x^2-2xy+y^2+2x-2y-3=0$ conjugate to the non-asymptotic direction $X:Y$.}
因为
\[
F_1=x-y+1,\qquad F_2=-x+y-1,
\]
所以直径方程为
\[
X(x-y+1)+Y(-x+y-1)=0,
\]
即
\[
(X-Y)(x-y+1)=0.
\]
因为已知曲线 $F(x,y)=0$ 的渐近方向为 $X':Y'=1:1$，所以对于非渐近方向 $X:Y$ 一定有 $X\ne Y$，因此曲线的共轭于非渐近方向 $X:Y$ 的直径为
\[
x-y+1=0,
\]
它只有一条直径（线心曲线）。
\EN{Since $F_1=x-y+1$ and $F_2=-(x-y+1)$, the diameter equation is $(X-Y)(x-y+1)=0$; the asymptotic direction is $1:1$, so for a non-asymptotic direction $X\ne Y$ and the single diameter is $x-y+1=0$.}

\subsection{共轭方向与共轭直径 | Conjugate directions and diameters}
\APS{把二次曲线的与非渐近方向 $X:Y$ 共轭的直径方向
\[
X':Y'=-(a_{12}X+a_{22}Y):(a_{11}X+a_{12}Y)
\]
叫做非渐近方向 $X:Y$ 的\textbf{共轭方向 / conjugate direction}。由 (5.4-1) 可得
\[
\Phi(X',Y')=(a_{22}a_{11}-a_{12}^2)\bigl(a_{11}X^2+2a_{12}XY+a_{22}Y^2\bigr)=I_2\,\Phi(X,Y).
\]
因为 $X:Y$ 为非渐近方向，所以 $\Phi(X,Y)\ne0$，因此当 $I_2\ne0$ 即二次曲线为中心曲线时，$\Phi(X',Y')\ne0$；当 $I_2=0$ 即二次曲线为非中心曲线时，$\Phi(X',Y')=0$。这就是说，中心二次曲线的非渐近方向的共轭方向仍然是非渐近方向，而在非中心二次曲线的情形是渐近方向。}{The diameter direction $X':Y'=-(a_{12}X+a_{22}Y):(a_{11}X+a_{12}Y)$ is the conjugate direction of $X:Y$. Since $\Phi(X',Y')=I_2\Phi(X,Y)$, the conjugate of a non-asymptotic direction is again non-asymptotic for a central conic and asymptotic for a non-central one.}

\APS{由上式得二次曲线的非渐近方向 $X:Y$ 与它的共轭方向 $X':Y'$ 之间的关系
\[
a_{11}XX'+a_{12}(XY'+X'Y)+a_{22}YY'=0.\tag{5.4-5}
\]
从 (5.4-5) 看出，两个方向 $X:Y$ 与 $X':Y'$ 是对称的，因此对中心曲线来说，非渐近方向 $X:Y$ 的共轭方向为非渐近方向 $X':Y'$，而 $X':Y'$ 的共轭方向就是 $X:Y$。}{Relation (5.4-5) is symmetric in the two directions, so for a central conic conjugation is an involution.}

\begin{definition}[5.4.2 共轭直径 / conjugate diameters]\label{def:conjdiam}
中心曲线的一对具有相互共轭方向的直径叫做一对\textbf{共轭直径 / conjugate diameters}。
\EN{A pair of diameters of a central conic in mutually conjugate directions is a pair of conjugate diameters.}
\end{definition}

\APS{设 $k$ 与 $k'$ 为一对共轭直径的斜率，即 $k=Y/X$，$k'=Y'/X'$，代入 (5.4-5) 得
\[
a_{11}+a_{12}(k+k')+a_{22}kk'=0,\tag{5.4-6}
\]
这就是一对共轭直径的斜率满足的关系式。例如椭圆 $\dfrac{x^2}{a^2}+\dfrac{y^2}{b^2}=1$ 的一对共轭直径的斜率 $k$ 与 $k'$ 有着关系
\[
kk'=-\frac{b^2}{a^2},\tag{5.4-7}
\]
而双曲线 $\dfrac{x^2}{a^2}-\dfrac{y^2}{b^2}=1$ 的一对共轭直径的斜率 $k$ 与 $k'$ 有着关系
\[
kk'=\frac{b^2}{a^2}.\tag{5.4-8}
\]
如果在 (5.4-5) 中设 $X'=X$，$Y'=Y$，那么有
\[
a_{11}X^2+2a_{12}XY+a_{22}Y^2=0,
\]
显然此时 $X:Y$ 为二次曲线的渐近方向。因此如果对二次曲线的共轭方向从 (5.4-5) 作代数的推广，那么渐近方向可以看成与自己共轭的方向，从而渐近线也就可以看成与自己共轭的直径。}{With slopes $k=Y/X$, $k'=Y'/X'$ relation (5.4-5) becomes (5.4-6). For the ellipse $kk'=-b^2/a^2$ and for the hyperbola $kk'=b^2/a^2$. Taking $X'=X$, $Y'=Y$ in (5.4-5) gives $\Phi(X,Y)=0$: an asymptotic direction is self-conjugate, and an asymptote is a self-conjugate diameter.}

\ASK{为什么'共轭'这个词同时用于弦、直径、方向？它们的关系是什么？}{Why is the word ``conjugate'' used for chords, diameters and directions?}{一根直径是``与其共轭的平行弦''中点轨迹，所以直径与弦方向共轭；直径的``方向''就是其自身直线的方向，于是弦方向与直径方向共轭。两条直径若其方向互相共轭，则称为共轭直径。由于 (5.4-5) 关于两个方向对称，中心曲线的共轭关系是``对合''：$X:Y$ 的共轭是 $X':Y'$，$X':Y'$ 的共轭回到 $X:Y$。抛物线等非中心曲线中，非渐近方向的共轭方向恰是渐近方向，所以只有一个方向的直径（平行束）。}{A diameter is the locus of midpoints of its conjugate chords, so the diameter is conjugate to the chord direction; the direction of the diameter is conjugate to the chord direction. Two diameters whose directions are conjugate form conjugate diameters. Because (5.4-5) is symmetric, conjugation is an involution on a central conic; for a non-central conic the conjugate of a non-asymptotic direction is the asymptotic direction, so the diameters form a parallel pencil.}

\section{\label{sec:13}二次曲线的主直径与主方向 | Principal diameters and directions}
\APS{主直径就是``垂直于其共轭弦的直径''，也就是二次曲线的对称轴。求主直径的问题可以化成求矩阵 $A^*$ 的特征向量问题：主方向 $(X,Y)$ 满足 $a_{11}X+a_{12}Y=\lambda X$，$a_{12}X+a_{22}Y=\lambda Y$，其中 $\lambda$ 是特征方程 $\lambda^2-I_1\lambda+I_2=0$ 的根。由于 $A^*$ 是对称矩阵，特征根为实数且不同特征根的主方向正交（定理 5.5.1、5.5.4），这正是``对称轴互相垂直''的代数版本。}{A principal diameter is a diameter perpendicular to its conjugate chords, i.e.\ an axis of symmetry. Finding principal diameters reduces to the eigenproblem of the symmetric matrix $A^*$: a principal direction $(X,Y)$ satisfies $a_{11}X+a_{12}Y=\lambda X$, $a_{12}X+a_{22}Y=\lambda Y$, where $\lambda$ solves $\lambda^2-I_1\lambda+I_2=0$. Since $A^*$ is symmetric, the eigenvalues are real and eigenvectors for distinct eigenvalues are orthogonal (Theorems 5.5.1, 5.5.4), the algebraic form of ``axes are perpendicular''.}

\begin{definition}[5.5.1 主直径与主方向 / principal diameter and direction]\label{def:principal}
二次曲线的垂直于其共轭弦的直径叫做二次曲线的\textbf{主直径 / principal diameter}，主直径的方向与垂直于主直径的方向都叫做二次曲线的\textbf{主方向 / principal direction}。
\EN{A diameter perpendicular to its conjugate chords is a principal diameter; both its direction and the perpendicular direction are principal directions.}
\end{definition}

\APS{显然，主直径是二次曲线的对称轴，因此主直径也叫做二次曲线的\textbf{轴 / axis}，轴与曲线的交点叫做曲线的\textbf{顶点 / vertex}。}{Clearly a principal diameter is an axis of symmetry; its intersections with the conic are the vertices.}

\APS{现在我们在直角坐标系下来求二次曲线
\[
F(x,y)\equiv a_{11}x^2+2a_{12}xy+a_{22}y^2+2a_{13}x+2a_{23}y+a_{33}=0\tag{1}
\]
的主方向与主直径。如果二次曲线 (1) 为中心曲线，那么与二次曲线 (1) 的非渐近方向 $X:Y$ 共轭的直径为 (5.4-1) 或 (5.4-3)。设直径的方向为 $X':Y'$，那么
\[
X':Y'=-(a_{12}X+a_{22}Y):(a_{11}X+a_{12}Y).\tag{2}
\]
根据主方向的定义，$X:Y$ 成为主方向的条件是它垂直于它的共轭方向，在直角坐标系下，由 $XX'+YY'=0$ 得
\[
a_{11}X+a_{12}Y=\lambda X,\qquad a_{12}X+a_{22}Y=\lambda Y\tag{5.5-1}
\]
成立，其中 $\lambda\ne0$（$\lambda$ 为比例因子）。把它改写成
\[
\begin{cases}(a_{11}-\lambda)X+a_{12}Y=0,\\ a_{12}X+(a_{22}-\lambda)Y=0.\end{cases}\tag{5.5-1$'$}
\]
这是一个关于 $X,Y$ 的齐次线性方程组，而 $X,Y$ 不能全为零，所以
\[
\begin{vmatrix}a_{11}-\lambda&a_{12}\\ a_{12}&a_{22}-\lambda\end{vmatrix}=0,\tag{5.5-2}
\]
即
\[
\lambda^2-I_1\lambda+I_2=0.\tag{5.5-3}
\]
因此对于中心二次曲线来说，只要由 (5.5-3) 解出 $\lambda$，再代入 (5.5-1) 就能得到它的主方向。}{For a central conic, the conjugate diameter has direction (2); the principal direction condition $XX'+YY'=0$ gives (5.5-1), and a nonzero solution requires the characteristic equation (5.5-2), i.e.\ (5.5-3). Solving it and substituting into (5.5-1) gives the principal directions.}

\APS{如果二次曲线 (1) 为非中心二次曲线，那么它的任何直径的方向总是它的惟一的渐近方向
\[
X:Y=-a_{12}:a_{11}\quad(\text{或 } -a_{22}:a_{12}),
\]
而垂直于它的方向显然为 $X:Y=a_{11}:-a_{12}$，所以非中心二次曲线 (1) 的主方向为
\[
\text{渐近主方向：}-a_{12}:a_{11};\qquad \text{非渐近主方向：}a_{11}:-a_{12}.
\]
如果我们把 (5.5-2) 或 (5.5-3) 推广到非中心二次曲线，即式中的 $I_2$ 可取等于零，这样当 $I_2=0$ 时，方程 (5.5-3) 的两根为 $\lambda_1=a_{11}+a_{22}$，$\lambda_2=0$，把它们代入 (5.5-1) 或 (5.5-1$'$) 所得的主方向，正是非中心二次曲线的渐近主方向与非渐近主方向。因此，一个方向 $X:Y$ 成为二次曲线 (1) 的主方向的条件是 (5.5-1) 成立，这里的 $\lambda$ 是方程 (5.5-2) 或 (5.5-3) 的根。}{For a non-central conic every diameter is parallel to the unique asymptotic direction, and the perpendicular direction is the other principal direction; the same equations with $I_2=0$ give $\lambda_1=a_{11}+a_{22}$, $\lambda_2=0$, recovering the two principal directions. Thus (5.5-1) with $\lambda$ a root of (5.5-2) characterises principal directions in general.}

\begin{definition}[5.5.2 特征方程与特征根 / characteristic equation]\label{def:chareq}
方程 (5.5-2) 或 (5.5-3) 叫做二次曲线 (1) 的\textbf{特征方程 / characteristic equation}，特征方程的根叫做二次曲线的\textbf{特征根 / eigenvalues}。
\EN{Equations (5.5-2), (5.5-3) are the characteristic equation of (1), and its roots are the eigenvalues.}
\end{definition}

\APS{从二次曲线 (1) 的特征方程 (5.5-3) 求出特征根 $\lambda$，把它代入 (5.5-1) 或 (5.5-1$'$)，我们就得到相应的主方向。如果主方向为非渐近方向，那么根据 (5.4-1) 就能得到共轭于它的主直径。}{Solve (5.5-3) for $\lambda$, substitute into (5.5-1) to get a principal direction, and if it is non-asymptotic use (5.4-1) to get the corresponding principal diameter.}

\begin{theorem}[5.5.1 特征根为实 / real eigenvalues]\label{thm:real-eig}
二次曲线的特征根都是实数。
\EN{The eigenvalues of a conic are real.}
\end{theorem}

\begin{proof}[证明 / Proof]
因为特征方程的判别式
\[
\Delta=I_1^2-4I_2=(a_{11}+a_{22})^2-4(a_{11}a_{22}-a_{12}^2)=(a_{11}-a_{22})^2+4a_{12}^2\ge0,
\]
所以二次曲线的特征根都是实数。
\EN{The discriminant is $(a_{11}-a_{22})^2+4a_{12}^2\ge0$, so both roots are real.}
\end{proof}

\begin{theorem}[5.5.2 特征根不全为零 / not all eigenvalues vanish]\label{thm:nonzero-eig}
二次曲线的特征根不能全为零。
\EN{The eigenvalues of a conic cannot both be zero.}
\end{theorem}

\begin{proof}[证明 / Proof]
如果二次曲线的特征根全为零，那么由 (5.5-3) 得 $I_1=a_{11}+a_{22}=0$ 与 $I_2=a_{11}a_{22}-a_{12}^2=0$，从而得 $a_{22}=-a_{11}$ 与 $-a_{11}^2-a_{12}^2=0$，即 $a_{11}=a_{12}=a_{22}=0$，这与二次曲线的定义矛盾，所以二次曲线的特征根不能全为零。
\EN{If both roots vanished then $I_1=I_2=0$, forcing $a_{11}=a_{12}=a_{22}=0$, contradicting the definition of a conic.}
\end{proof}

\begin{theorem}[5.5.3 特征根与主方向的类型 / eigenvalues and principal directions]\label{thm:eig-dir}
由二次曲线 (1) 的特征根 $\lambda$ 确定的主方向 $X:Y$，当 $\lambda\ne0$ 时，为二次曲线的非渐近主方向；当 $\lambda=0$ 时，为二次曲线的渐近主方向。
\EN{The principal direction belonging to an eigenvalue $\lambda$ is non-asymptotic when $\lambda\ne0$ and asymptotic when $\lambda=0$.}
\end{theorem}

\begin{proof}[证明 / Proof]
因为
\[
\Phi(X,Y)=a_{11}X^2+2a_{12}XY+a_{22}Y^2=(a_{11}X+a_{12}Y)X+(a_{12}X+a_{22}Y)Y,
\]
所以由 (5.5-1) 得
\[
\Phi(X,Y)=\lambda X^2+\lambda Y^2=\lambda(X^2+Y^2).
\]
又因为 $X,Y$ 不全为零，所以当 $\lambda\ne0$ 时 $\Phi(X,Y)\ne0$，$X:Y$ 为二次曲线 (1) 的非渐近主方向；当 $\lambda=0$ 时 $\Phi(X,Y)=0$，$X:Y$ 为二次曲线 (1) 的渐近主方向。
\EN{Using (5.5-1), $\Phi(X,Y)=\lambda(X^2+Y^2)$; since $(X,Y)\ne(0,0)$ the sign of $\lambda$ decides non-asymptotic versus asymptotic.}
\end{proof}

\begin{theorem}[5.5.4 主直径的条数 / number of principal diameters]\label{thm:num-principal}
中心二次曲线至少有两条主直径，非中心二次曲线只有一条主直径。
\EN{A central conic has at least two principal diameters; a non-central conic has exactly one.}
\end{theorem}

\begin{proof}[证明 / Proof]
由二次曲线 (1) 的特征方程 (5.5-3) 解得两特征根为
\[
\lambda_{1,2}=\frac{I_1\pm\sqrt{I_1^2-4I_2}}{2}.
\]
$1^\circ$ 当二次曲线 (1) 为中心曲线时，$I_2\ne0$。如果特征方程的判别式 $\Delta=I_1^2-4I_2=(a_{11}-a_{22})^2+4a_{12}^2=0$，那么 $a_{11}=a_{22}$，$a_{12}=0$，这时中心曲线为圆（包括点圆和虚圆），它的特征根为一对二重根 $\lambda=a_{11}=a_{22}\ne0$。把它代入 (5.5-1) 或 (5.5-1$'$)，则得到两个恒等式，能被任何方向 $X:Y$ 所满足，所以任何实方向都是圆的非渐近主方向，从而通过圆心的任何直线不仅都是直径（见 §5.4 习题第 6 题），而且都是圆的主直径。如果特征方程的判别式 $\Delta=(a_{11}-a_{22})^2+4a_{12}^2>0$，那么特征根为两不等的非零实根 $\lambda_1,\lambda_2$，将它们分别代入 (5.5-1$'$) 得相应的两非渐近主方向为
\[
X_1:Y_1=a_{12}:(\lambda_1-a_{11})=(\lambda_1-a_{22}):a_{12},
\]
\[
X_2:Y_2=a_{12}:(\lambda_2-a_{11})=(\lambda_2-a_{22}):a_{12}.
\]
这两主方向相互垂直（见本节习题第 4 题），从而它们又互相共轭，因此非圆的中心二次曲线有且只有一对互相垂直从而又互相共轭的主直径。$2^\circ$ 当二次曲线 (1) 为非中心曲线时，$I_2=0$，这时两特征根为 $\lambda_1=a_{11}+a_{22}$，$\lambda_2=0$。所以它只有一个非渐近主方向，即与 $\lambda_1$ 相应的主方向，从而非中心二次曲线只有一条主直径。
\EN{The roots are $\lambda_{1,2}=(I_1\pm\sqrt{I_1^2-4I_2})/2$. For a central conic, if the discriminant vanishes the conic is a circle and every diameter is principal; otherwise the two distinct nonzero roots give two perpendicular (hence conjugate) principal directions, i.e.\ two principal diameters. For a non-central conic $I_2=0$ and the roots are $a_{11}+a_{22}$ and $0$; only the nonzero root gives a non-asymptotic principal direction, hence one principal diameter.}
\end{proof}

\EX{（例 1，书中第 144--145 页）求 $F(x,y)\equiv x^2-xy+y^2-1=0$ 的主方向与主直径。}{（Example 1, pp.~144--145）Find the principal directions and principal diameters of $F=x^2-xy+y^2-1=0$.}
因为
\[
I_1=1+1=2,\qquad I_2=\begin{vmatrix}1&-\frac12\\ -\frac12&1\end{vmatrix}=1-\frac14=\frac34\ne0,
\]
所以曲线为中心曲线，它的特征方程为
\[
\lambda^2-2\lambda+\frac34=0,
\]
解这方程得两特征根为
\[
\lambda_1=\frac12,\qquad \lambda_2=\frac32.
\]
由特征根 $\lambda_1=\frac12$ 确定的主方向为（由 $a_{12}X+(a_{22}-\lambda_1)Y=0$，即 $-\frac12X+\frac12Y=0$）
\[
X_1:Y_1=1:1;
\]
由特征根 $\lambda_2=\frac32$ 确定的主方向为（$-\frac12X-\frac12Y=0$）
\[
X_2:Y_2=-1:1.
\]
又因为
\[
F_1(x,y)=x-\frac12y,\qquad F_2(x,y)=-\frac12x+y,
\]
所以曲线的主直径为
\[
\Bigl(x-\frac12y\Bigr)+\Bigl(-\frac12x+y\Bigr)=0\quad\text{与}\quad -\Bigl(x-\frac12y\Bigr)+\Bigl(-\frac12x+y\Bigr)=0,
\]
即
\[
x+y=0\quad\text{与}\quad x-y=0.
\]
\EN{Here $I_1=2$, $I_2=3/4$; the characteristic equation is $\lambda^2-2\lambda+3/4=0$ with roots $1/2$ and $3/2$. The corresponding directions are $1:1$ and $-1:1$, and using $F_1=x-y/2$, $F_2=-x/2+y$ the principal diameters are $x+y=0$ and $x-y=0$.}

\EX{（例 2，书中第 145 页）求曲线 $F(x,y)\equiv x^2-2xy+y^2-4x=0$ 的主方向与主直径。}{（Example 2, p.~145）Find the principal directions and diameter of $F=x^2-2xy+y^2-4x=0$.}
因为
\[
I_1=1+1=2,\qquad I_2=\begin{vmatrix}1&-1\\ -1&1\end{vmatrix}=0,
\]
所以曲线为非中心曲线，它的特征方程为
\[
\lambda^2-2\lambda=0,
\]
因此两特征根为
\[
\lambda_1=2,\qquad \lambda_2=0.
\]
由这两特征根所确定的主方向为
\[
\text{非渐近主方向：}X_1:Y_1=-1:(2-1)=-1:1;\qquad \text{渐近主方向：}X_2:Y_2=-1:(0-1)=1:1.
\]
又因为 $F_1(x,y)=x-y-2$，$F_2(x,y)=-x+y$，所以曲线的惟一主直径为
\[
-(x-y-2)+(-x+y)=0,
\]
即
\[
x-y-1=0.
\]
\EN{Here $I_1=2$, $I_2=0$, so the conic is non-central; the characteristic equation is $\lambda^2-2\lambda=0$ with roots $2$ and $0$. The principal directions are $-1:1$ (non-asymptotic) and $1:1$ (asymptotic); with $F_1=x-y-2$, $F_2=-x+y$ the single principal diameter is $x-y-1=0$.}

\ASK{为什么主方向恰好是矩阵 $A^*$ 的特征向量？主直径为什么一定是对称轴？}{Why are principal directions eigenvectors of $A^*$? Why is a principal diameter an axis of symmetry?}{直径方向与共轭弦方向共轭，条件是 (5.4-5)；主方向要求它与共轭方向垂直，即 $XX'+YY'=0$。把共轭方向公式代入并整理，正得到 $A^*(X,Y)^{\mathsf T}=\lambda(X,Y)^{\mathsf T}$，所以主方向就是 $A^*$ 的特征向量。主直径是垂直于自己共轭弦的直径，而垂直于弦方向的直径上的弦中点都在这条直径上，所以直径两侧的弦对称，即它是曲线的对称轴。}{The diameter direction is conjugate to the chord direction by (5.4-5); a principal direction is perpendicular to its conjugate direction, $XX'+YY'=0$. Substituting the conjugate-direction formula yields $A^*(X,Y)^{\mathsf T}=\lambda(X,Y)^{\mathsf T}$, hence principal directions are eigenvectors. A principal diameter is perpendicular to its conjugate chords and consists of their midpoints, so the chords are symmetric about it: it is an axis of symmetry.}

\section{\label{sec:14}二次曲线方程的化简与分类 | Reduction and classification}
\APS{本节把``几何''与``代数''接起来：选取适当的坐标系（移轴使原点落在中心或顶点，转轴使坐标轴与主方向平行），把二次曲线方程化为三种简化方程 (I)(II)(III)，进而化为九种标准方程。关键事实有三条：移轴下二次项系数不变、一次项系数变为 $(2F_1,2F_2)$、常数项变为 $F(x_0,y_0)$；转轴下二次项按二次型合同变换改变、一次项按向量旋转改变、常数项不变；总可以选转角使 $xy$ 项消失。几何上，这一套化简等价于``把坐标轴搬到主直径上''。}{This section joins geometry and algebra: choose a convenient coordinate system (translate the origin to the centre or vertex, rotate the axes to the principal directions), reduce the conic to one of the three simplified equations (I)(II)(III), and then to the nine standard equations. Three facts matter: under translation the quadratic coefficients are unchanged, the linear coefficients become $(2F_1,2F_2)$ and the constant becomes $F(x_0,y_0)$; under rotation the quadratic part changes by congruence, the linear part rotates as a vector and the constant is unchanged; and the angle can always be chosen to kill the $xy$ term. Geometrically this is exactly moving the axes onto the principal diameters.}

\subsection{平面直角坐标变换 | Rectangular coordinate changes}
\APS{如果平面内一点的旧坐标与新坐标分别为 $(x,y)$ 与 $(x',y')$，那么\textbf{移轴 / translation}公式为
\[
\begin{cases}x=x'+x_0,\\ y=y'+y_0\end{cases}\qquad\text{或}\qquad
\begin{cases}x'=x-x_0,\\ y'=y-y_0,\end{cases}\tag{5.6-1}
\]
式中 $(x_0,y_0)$ 为新坐标系原点在旧坐标系里的坐标。\textbf{转轴 / rotation}公式为
\[
\begin{cases}x=x'\cos\alpha-y'\sin\alpha,\\ y=x'\sin\alpha+y'\cos\alpha\end{cases}\tag{5.6-2}
\]
或
\[
\begin{cases}x'=x\cos\alpha+y\sin\alpha,\\ y'=-x\sin\alpha+y\cos\alpha,\end{cases}\tag{5.6-2$'$}
\]
式中的 $\alpha$ 为坐标轴的旋转角。}{Translation and rotation formulas are (5.6-1) and (5.6-2), (5.6-2$'$).}

\APS{而在一般情形，由旧坐标系 $O-xy$ 变成新坐标系 $O'-x'y'$，总可以分两步来完成：先移轴使坐标原点与新坐标系的原点 $O'$ 重合，变成辅助坐标系 $O'-x''y''$，然后再由辅助坐标系转轴 $\alpha$ 而成新坐标系 $O'-x'y'$。于是得一般坐标变换公式
\[
\begin{cases}x=x'\cos\alpha-y'\sin\alpha+x_0,\\ y=x'\sin\alpha+y'\cos\alpha+y_0,\end{cases}\tag{5.6-3}
\]
由 (5.6-3) 解出 $x',y'$ 便得逆变换公式
\[
\begin{cases}x'=(x-x_0)\cos\alpha+(y-y_0)\sin\alpha,\\ y'=-(x-x_0)\sin\alpha+(y-y_0)\cos\alpha.\end{cases}\tag{5.6-4}
\]
平面直角坐标变换公式 (5.6-3) 是由新坐标系原点的坐标 $(x_0,y_0)$ 与坐标轴的旋转角 $\alpha$ 决定的。}{A general change is translation followed by rotation, giving (5.6-3) and the inverse (5.6-4); it is determined by the new origin and the rotation angle.}

\APS{确定坐标变换公式还有另一种常用方法：给出新坐标系两坐标轴在旧坐标系里的方程。设在直角坐标系 $O-xy$ 里给定了两条相互垂直的直线
\[
l_1:A_1x+B_1y+C_1=0,\qquad l_2:A_2x+B_2y+C_2=0,
\]
其中 $A_1A_2+B_1B_2=0$。如果取直线 $l_1$ 为新坐标系中的横轴 $O'x'$，而直线 $l_2$ 为纵轴 $O'y'$，设平面上任意点 $M$ 的旧坐标与新坐标分别是 $(x,y)$ 与 $(x',y')$，因为 $O'y'$ 是点 $M$ 到 $O'y'$ 轴的距离，也就是 $M$ 点到 $l_2$ 的距离，因此有
\[
|x'|=\frac{|A_2x+B_2y+C_2|}{\sqrt{A_2^2+B_2^2}},\qquad |y'|=\frac{|A_1x+B_1y+C_1|}{\sqrt{A_1^2+B_1^2}}.
\]
去掉绝对值符号后，为使新坐标系仍是右手坐标系，符号的选取要使得 (5.6-3) 与 (5.6-4) 相比较时，第一式右端 $x$ 的系数与第二式右端 $y$ 的系数相等（都等于 $\cos\alpha$）。}{If the new axes are given by two perpendicular lines, $|x'|$ and $|y'|$ are the distances from $M$ to $l_2$ and $l_1$; the signs are chosen so that the resulting formulas are (5.6-3) with a right-handed frame.}

\EX{（例 1，书中第 147--148 页，按书中方法整理）已知两垂直的直线 $l_1:x-2y+3=0$ 与 $l_2:2x+y-2=0$，取 $l_1$ 为 $x'$ 轴、$l_2$ 为 $y'$ 轴，求坐标变换公式。}{（Example 1, pp.~147--148）Given perpendicular lines $l_1:x-2y+3=0$, $l_2:2x+y-2=0$, take $l_1$ as the $x'$-axis and $l_2$ as the $y'$-axis; find the coordinate transformation.}
设 $M(x,y)$ 的新坐标为 $(x',y')$，则
\[
|x'|=\frac{|2x+y-2|}{\sqrt5},\qquad |y'|=\frac{|x-2y+3|}{\sqrt5}.
\]
取定一组符号，例如
\[
x'=\frac{2x+y-2}{\sqrt5},\qquad y'=\frac{-x+2y-3}{\sqrt5},
\]
这组符号使变换是右手系：将其与 (5.6-4) 比较，得 $\cos\alpha=\frac{2}{\sqrt5}$，$\sin\alpha=\frac{1}{\sqrt5}$，旋转角 $\alpha$ 是小于 $\frac{\pi}{2}$ 的正角；同时 $x_0\cos\alpha+y_0\sin\alpha=\frac{2}{\sqrt5}$ 与 $x_0\sin\alpha-y_0\cos\alpha=-\frac{3}{\sqrt5}$ 给出新原点 $(x_0,y_0)=\bigl(\frac15,\frac85\bigr)$，它正是 $l_1$ 与 $l_2$ 的交点。若取相反的符号，则得到旋转角为负、新原点相同的另一组公式（新 $y'$ 轴方向相反）。\\[2pt]
从这组公式解出 $x,y$ 即得逆变换
\[
x=\frac{2x'-y'}{\sqrt5}+\frac15,\qquad y=\frac{x'+2y'}{\sqrt5}+\frac85.
\]
\EN{Since $x'$ and $y'$ are distances to $l_2$ and $l_1$, $|x'|=|2x+y-2|/\sqrt5$ and $|y'|=|x-2y+3|/\sqrt5$. Choosing the signs $x'=(2x+y-2)/\sqrt5$, $y'=(-x+2y-3)/\sqrt5$ gives a right-handed frame with $\cos\alpha=2/\sqrt5$, $\sin\alpha=1/\sqrt5$ and new origin $(1/5,8/5)$, the intersection of $l_1,l_2$; the inverse is displayed. The opposite choice reverses the $y'$ direction.}

\APS{在坐标变换下，平面上曲线的方程将改变，但是如果曲线方程 $F(x,y)=0$ 的左端 $F(x,y)$ 是一个多项式，其次数为 $n$，那么通过坐标变换 (5.6-3)，它的新方程 $G(x',y')=0$ 的左端 $G(x',y')$ 将仍然是一个多项式，而且它的次数 $n'$ 不变，即 $n'=n$（因为变换式右端是一次式，代入后次数不增；逆变换同理）。我们把多项式 $F(x,y)$ 构成的方程 $F(x,y)=0$ 叫做\textbf{代数方程 / algebraic equation}，而由它表示的曲线叫做\textbf{代数曲线 / algebraic curve}，方程的次数叫做曲线的\textbf{次数 / degree}。这个性质与坐标系的选择无关。}{Under a coordinate change a polynomial equation stays polynomial of the same degree $n$; this degree is intrinsic to the algebraic curve.}

\subsection{二次方程的化简与分类 | Reduction and classification}
\APS{设二次曲线的方程为
\[
F(x,y)\equiv a_{11}x^2+2a_{12}xy+a_{22}y^2+2a_{13}x+2a_{23}y+a_{33}=0,\tag{1}
\]
现在要选取一个适当的坐标系，使得曲线 (1) 在新坐标系下的方程最为简单，这就是二次曲线方程的化简。为此先考察在移轴与转轴下，二次曲线方程 (1) 的系数的变换规律。}{We seek a coordinate system in which (1) takes its simplest form; first record how the coefficients change under translation and rotation.}

\APS{在移轴 (5.6-1) 下，二次曲线 (1) 的新方程为
\[
a_{11}x'^2+2a_{12}x'y'+a_{22}y'^2+2a'_{13}x'+2a'_{23}y'+a'_{33}=0,
\]
这里
\[
a'_{13}=a_{11}x_0+a_{12}y_0+a_{13}=F_1(x_0,y_0),\qquad a'_{23}=a_{12}x_0+a_{22}y_0+a_{23}=F_2(x_0,y_0),
\]
\[
a'_{33}=a_{11}x_0^2+2a_{12}x_0y_0+a_{22}y_0^2+2a_{13}x_0+2a_{23}y_0+a_{33}=F(x_0,y_0).\tag{5.6-6}
\]
因此在移轴 (5.6-1) 下，二次曲线方程系数的变换规律为：$1^\circ$ 二次项系数不变；$2^\circ$ 一次项系数变为 $2F_1(x_0,y_0)$ 与 $2F_2(x_0,y_0)$；$3^\circ$ 常数项变为 $F(x_0,y_0)$。因为当 $(x_0,y_0)$ 为二次曲线 (1) 的中心时有 $F_1(x_0,y_0)=F_2(x_0,y_0)=0$，所以当二次曲线有中心时，作移轴使原点与中心重合，那么在新坐标系下二次曲线的新方程中一次项消失。}{Under translation the quadratic coefficients are unchanged, the linear coefficients become $2F_1,2F_2$ and the constant becomes $F(x_0,y_0)$; translating the origin to the centre kills the linear terms.}

\APS{把转轴公式 (5.6-2) 代入 (1)，得在转轴下二次曲线 (1) 的新方程
\[
a'_{11}x'^2+2a'_{12}x'y'+a'_{22}y'^2+2a'_{13}x'+2a'_{23}y'+a'_{33}=0,
\]
这里
\[
a'_{11}=a_{11}\cos^2\alpha+2a_{12}\sin\alpha\cos\alpha+a_{22}\sin^2\alpha,
\]
\[
a'_{12}=(a_{22}-a_{11})\sin\alpha\cos\alpha+a_{12}(\cos^2\alpha-\sin^2\alpha),
\]
\[
a'_{22}=a_{11}\sin^2\alpha-2a_{12}\sin\alpha\cos\alpha+a_{22}\cos^2\alpha,
\]
\[
a'_{13}=a_{13}\cos\alpha+a_{23}\sin\alpha,\qquad a'_{23}=-a_{13}\sin\alpha+a_{23}\cos\alpha,\qquad a'_{33}=a_{33}.\tag{5.6-7}
\]
因此，在转轴下，二次曲线方程 (1) 的系数的变换规律为：$1^\circ$ 二次项系数一般要改变，新方程的二次项系数仅与原方程的二次项系数及旋转角有关，而与一次项系数及常数项无关；$2^\circ$ 一次项系数一般要改变，新方程的一次项系数仅与原方程的一次项系数及旋转角有关，与二次项系数及常数项无关；$3^\circ$ 常数项不变。}{Under rotation the quadratic coefficients change by congruence, the linear coefficients rotate as a vector, and the constant is unchanged (5.6-7).}

\APS{二次曲线方程 (1) 里，如果 $a_{12}\ne0$，我们往往使用转轴使新方程中的 $a'_{12}=0$。为此，我们只要取旋转角 $\alpha$，使得
\[
a'_{12}=(a_{22}-a_{11})\sin\alpha\cos\alpha+a_{12}(\cos^2\alpha-\sin^2\alpha)=0,
\]
即
\[
(a_{22}-a_{11})\sin2\alpha+2a_{12}\cos2\alpha=0,
\]
所以
\[
\cot2\alpha=\frac{a_{11}-a_{22}}{2a_{12}}.\tag{5.6-8}
\]
因为余切的值可以是任意的实数，所以总有 $\alpha$ 满足 (5.6-8)，也就是说总可以经过适当的转轴消去 (1) 的 $xy$ 项。}{To kill the $xy$ term choose $\alpha$ with $\cot2\alpha=(a_{11}-a_{22})/(2a_{12})$ (5.6-8); since cot takes all real values such an angle always exists.}

\EX{（例 3，书中第 151--152 页）化简二次曲线方程
\[
x^2-xy+y^2+2x-4y=0,
\]
并画出它的图形。}{（Example 3, pp.~151--152）Reduce the conic $x^2-xy+y^2+2x-4y=0$ and describe its graph.}
这里 $a_{11}=1$，$a_{12}=-\frac12$，$a_{22}=1$，$a_{13}=1$，$a_{23}=-2$，$a_{33}=0$，
\[
I_2=\begin{vmatrix}1&-\frac12\\ -\frac12&1\end{vmatrix}=\frac34\ne0,
\]
所以曲线为中心二次曲线。解中心方程组
\[
\begin{cases}F_1=x-\frac12y+1=0,\\ F_2=-\frac12x+y-2=0,\end{cases}
\]
得中心的坐标为 $x_0=0$，$y_0=2$。取 $(0,2)$ 为新原点，作移轴
\[
x=x',\qquad y=y'+2,
\]
原方程变为
\[
x'^2-x'y'+y'^2-4=0.
\]
再转轴消去 $x'y'$ 项，由 (5.6-8) 得 $\cot2\alpha=\frac{1-1}{2\cdot(-\frac12)}=0$，从而可取 $\alpha=\frac{\pi}{4}$，故转轴公式为
\[
x'=\frac{\sqrt2}{2}(x''-y''),\qquad y'=\frac{\sqrt2}{2}(x''+y'').
\]
经转轴后曲线的方程化为最简形式
\[
\frac12x''^2+\frac32y''^2-4=0,
\]
或写成标准形式
\[
\frac{x''^2}{8}+\frac{y''^2}{\frac83}=1,
\]
这是一个椭圆。
\EN{Since $I_2=3/4\ne0$ the conic is central; the centre solves $x-y/2+1=0$, $-x/2+y-2=0$, giving $(0,2)$. Translating gives $x'^2-x'y'+y'^2-4=0$; $\cot2\alpha=0$ so $\alpha=\pi/4$, and the rotated equation is $x''^2/2+3y''^2/2-4=0$, i.e.\ the ellipse $x''^2/8+y''^2/(8/3)=1$.}

\EX{（例 6，书中第 155--156 页）化简 $x^2-2xy+y^2+2x-2y-3=0$。}{（Example 6, pp.~155--156）Reduce $x^2-2xy+y^2+2x-2y-3=0$.}
已知曲线的矩阵为
\[
\begin{pmatrix}1&-1&1\\ -1&1&-1\\ 1&-1&-3\end{pmatrix},
\]
它的第一、第二两行成比例，曲线为线心曲线，它有惟一的直径即中心线，也是曲线的主直径，其方程是
\[
x-y+1=0.
\]
取它为新坐标系的 $x'$ 轴，再任意取垂直于这中心线的直线，比如 $x+y=0$ 为 $y'$ 轴，作坐标变换，代入已知方程，经过整理得
\[
2y'^2-4=0\quad\text{即}\quad y'^2=2,
\]
这是两条平行直线。对于线心曲线，也可以直接从原方程分解为两个一次因式：把方程改写为
\[
(x-y)^2+2(x-y)-3=0,
\]
所以
\[
(x-y+3)(x-y-1)=0,
\]
因此原方程表示两条直线
\[
x-y+3=0\quad\text{与}\quad x-y-1=0.
\]
\EN{The rows of the matrix being proportional, the conic is line-central with line of centres $x-y+1=0$, which is its principal diameter. Taking it as the $x'$-axis and $x+y=0$ as the $y'$-axis gives $2y'^2-4=0$, i.e.\ $y'^2=2$: two parallel lines. Factoring directly, $(x-y+3)(x-y-1)=0$.}

\begin{theorem}[5.6.1 三种简化方程 / three simplified equations]\label{thm:simplified}
适当选取坐标系，二次曲线的方程
\[
F(x,y)\equiv a_{11}x^2+2a_{12}xy+a_{22}y^2+2a_{13}x+2a_{23}y+a_{33}=0
\]
总可以化成下列三个简化方程中的一个：
\[
\text{(I) }a_{11}x^2+a_{22}y^2+a_{33}=0,\quad a_{11}a_{22}\ne0;
\]
\[
\text{(II) }a_{22}y^2+2a_{13}x=0,\quad a_{22}\ne0;
\]
\[
\text{(III) }a_{22}y^2+a_{33}=0,\quad a_{22}\ne0.
\]
\EN{By a suitable choice of coordinates the conic can be reduced to one of the three simplified equations (I), (II), (III).}
\end{theorem}

\begin{proof}[证明 / Proof]
根据二次曲线是中心曲线、无心曲线与线心曲线三种情况来讨论。

$1^\circ$ 当已知二次曲线为中心曲线时，我们取它的一对既共轭又互相垂直的主直径作为坐标轴建立直角坐标系。设二次曲线在这样的坐标系下的方程为
\[
a_{11}x^2+2a_{12}xy+a_{22}y^2+2a_{13}x+2a_{23}y+a_{33}=0,
\]
因为这时原点就是曲线的中心，所以根据定理 5.2.1 的推论知道 $a_{13}=a_{23}=0$。其次，二次曲线的两条主直径（即坐标轴）的方向为 $1:0$ 与 $0:1$，它们互相共轭，因此根据 (5.4-5) 有 $a_{12}=0$，所以曲线的方程为
\[
a_{11}x^2+a_{22}y^2+a_{33}=0,
\]
又因为它是中心曲线，所以 $I_2=a_{11}a_{22}-a_{12}^2=a_{11}a_{22}\ne0$，即为 (I)。

$2^\circ$ 当已知二次曲线为无心曲线时，取它的惟一主直径为 $x$ 轴，而过顶点（即主直径与曲线的交点）且以非渐近主方向为方向的直线（即过顶点垂直于主直径的直线）为 $y$ 轴建立坐标系，这时的曲线方程假设为
\[
a_{11}x^2+2a_{12}xy+a_{22}y^2+2a_{13}x+2a_{23}y+a_{33}=0,
\]
因为这时主直径的共轭方向为 $X:Y=0:1$，所以主直径的方程为 $a_{12}x+a_{22}y+a_{23}=0$，它就是 $x$ 轴，即与直线 $y=0$ 重合，所以有 $a_{12}=a_{23}=0$，$a_{22}\ne0$。又因为顶点与坐标原点重合，所以 $(0,0)$ 满足曲线方程，从而又有 $a_{33}=0$。其次，由于二次曲线为无心曲线，所以 $I_2=a_{11}a_{22}-a_{12}^2=0$，而 $a_{22}\ne0$，所以有 $a_{11}=0$，$a_{13}\ne0$，因而曲线的方程为
\[
a_{22}y^2+2a_{13}x=0,\quad a_{22}\ne0,
\]
即为 (II)。

$3^\circ$ 当已知二次曲线为线心曲线时，我们取它的中心直线（即曲线的惟一直径也是主直径）为 $y$ 轴，任意垂直于它的直线为 $x$ 轴建立坐标系。设曲线的方程为
\[
a_{11}x^2+2a_{12}xy+a_{22}y^2+2a_{13}x+2a_{23}y+a_{33}=0,
\]
因为线心二次曲线的中心直线的方程是方程 $a_{11}x+a_{12}y+a_{13}=0$ 与 $a_{12}x+a_{22}y+a_{23}=0$ 中的任何一个，第二个方程表示 $y$ 轴的条件为 $a_{12}=a_{23}=0$，$a_{22}\ne0$；而第一个方程在 $a_{12}=0$ 的条件下，不可能再表示 $y$ 轴，所以它必须是恒等式，因而有 $a_{11}=a_{13}=0$。所以线心二次曲线的方程为
\[
a_{22}y^2+a_{33}=0,\quad a_{22}\ne0,
\]
即为 (III)。定理证毕。
\EN{For a central conic take the pair of perpendicular conjugate principal diameters as axes: the origin is the centre so $a_{13}=a_{23}=0$, and conjugate axes give $a_{12}=0$, hence (I) with $a_{11}a_{22}=I_2\ne0$. For a non-central conic take the single principal diameter as the $x$-axis and the perpendicular through the vertex as the $y$-axis: then $a_{12}=a_{23}=0$, the vertex at the origin gives $a_{33}=0$, and $I_2=0$ forces $a_{11}=0$, $a_{13}\ne0$, hence (II). For a line-central conic take the line of centres as the $y$-axis: then $a_{12}=a_{23}=0$ and the first equation must be an identity, so $a_{11}=a_{13}=0$, hence (III).}
\end{proof}

\APS{现在可以根据二次曲线三种简化方程系数的各种不同情况，写出二次曲线的各种标准方程，从而得出二次曲线的分类。$1^\circ$ \textbf{中心曲线}：
\[
a_{11}x^2+a_{22}y^2+a_{33}=0,\quad a_{11}a_{22}\ne0.
\]
当 $a_{33}\ne0$ 时，方程可化为
\[
Ax^2+By^2=1,
\]
其中 $A=-\frac{a_{11}}{a_{33}}$，$B=-\frac{a_{22}}{a_{33}}$。如果 $A>0$，$B>0$，那么设 $A=\frac1{a^2}$，$B=\frac1{b^2}$，于是得方程 $\frac{x^2}{a^2}+\frac{y^2}{b^2}=1$（椭圆）；如果 $A<0$，$B<0$，那么设 $A=-\frac1{a^2}$，$B=-\frac1{b^2}$，于是得 $\frac{x^2}{a^2}+\frac{y^2}{b^2}=-1$（虚椭圆）；如果 $A$ 与 $B$ 异号，那么不失一般性可设 $A>0$，$B<0$，于是得 $\frac{x^2}{a^2}-\frac{y^2}{b^2}=1$（双曲线）。当 $a_{33}=0$ 时，如果 $a_{11}$ 与 $a_{22}$ 异号，得 $\frac{x^2}{a^2}-\frac{y^2}{b^2}=0$（两相交直线）；如果 $a_{11}$ 与 $a_{22}$ 同号，得 $\frac{x^2}{a^2}+\frac{y^2}{b^2}=0$（点或称两相交于实点的共轭虚直线）。

$2^\circ$ \textbf{无心曲线}：$a_{22}y^2+2a_{13}x=0$，$a_{22}a_{13}\ne0$。设 $-\frac{2a_{13}}{a_{22}}=2p$，于是得
\[
y^2=2px\quad\text{（抛物线）}.
\]

$3^\circ$ \textbf{线心曲线}：$a_{22}y^2+a_{33}=0$，$a_{22}\ne0$。当 $a_{33}$ 与 $a_{22}$ 异号，设 $-\frac{a_{33}}{a_{22}}=a^2$，得 $y^2=a^2$（两平行直线）；当 $a_{33}$ 与 $a_{22}$ 同号，设 $\frac{a_{33}}{a_{22}}=a^2$，得 $y^2=-a^2$（两平行共轭虚直线）；当 $a_{33}=0$ 时得 $y^2=0$（两重合直线）。}{Discussing the signs in (I)--(III) yields the nine standard forms: ellipse, imaginary ellipse, hyperbola, a point (pair of conjugate imaginary lines), a pair of intersecting lines, parabola, a pair of parallel lines, a pair of parallel conjugate imaginary lines, and a double line.}

\begin{theorem}[5.6.2 九种标准方程 / nine standard forms]\label{thm:nine}
通过适当地选取坐标系，二次曲线的方程总可以写成下面九种标准方程的一种形式：
\[
{[1]}\ \frac{x^2}{a^2}+\frac{y^2}{b^2}=1\ \text{（椭圆）};\qquad
{[2]}\ \frac{x^2}{a^2}+\frac{y^2}{b^2}=-1\ \text{（虚椭圆）};
\]
\[
{[3]}\ \frac{x^2}{a^2}-\frac{y^2}{b^2}=1\ \text{（双曲线）};\qquad
[4]\ \frac{x^2}{a^2}+\frac{y^2}{b^2}=0\ \text{（点或称两相交于实点的共轭虚直线）};
\]
\[
{[5]}\ \frac{x^2}{a^2}-\frac{y^2}{b^2}=0\ \text{（两相交直线）};\qquad
[6]\ y^2=2px\ \text{（抛物线）};
\]
\[
{[7]}\ y^2=a^2\ \text{（两平行直线）};\qquad
[8]\ y^2=-a^2\ \text{（两平行共轭虚直线）};\qquad
[9]\ y^2=0\ \text{（两重合直线）}.
\]
\EN{By a suitable choice of coordinates every conic can be written as one of the nine displayed standard equations: ellipse, imaginary ellipse, hyperbola, a point, two intersecting lines, parabola, two parallel lines, two parallel conjugate imaginary lines, two coincident lines.}
\end{theorem}

\ASK{化简二次曲线有哪两条路线？各自的优缺点是什么？}{What are the two routes to reducing a conic? What are their merits?}{路线一（坐标变换法）：先移轴到中心（或顶点），再用 $\cot2\alpha=\frac{a_{11}-a_{22}}{2a_{12}}$ 转轴消去 $xy$ 项，得到简化方程与标准方程，同时得到坐标变换公式与图形；优点是几何意义清楚（坐标轴即主直径），并能写出变换后的坐标；缺点是计算量大。路线二（不变量法）：只算 $I_1,I_2,I_3,K_1$ 就得到简化方程（定理 5.7.3--5.7.5），不必求变换公式；优点是快，适合判别类型；缺点是不能直接给出图形的位置与坐标变换。两条路线合起来就是``定性判别 + 定量化简''。}{Route 1 (coordinate changes): translate to the centre or vertex, then rotate by $\cot2\alpha=(a_{11}-a_{22})/(2a_{12})$ to kill $xy$; this gives the simplified and standard equations together with the transformation and the position of the graph, at the cost of computation. Route 2 (invariants): compute $I_1,I_2,I_3,K_1$ to get the simplified equation (Theorems 5.7.3--5.7.5) without finding the transformation; fast for classification, but gives neither position nor transformation. Together they give qualitative classification plus quantitative reduction.}

\section{\label{sec:15}应用不变量化简二次曲线的方程 | Reduction by invariants}
\APS{本节是全章的``压缩''：二次曲线的形状与大小完全可以由四个数 $I_1,I_2,I_3,K_1$ 决定，而不必解方程组、求转角。定义 5.7.1 给出不变量的严格定义；定理 5.7.1 说明 $I_1,I_2,I_3$ 在任意直角坐标变换下不变、$K_1$ 只在转轴下不变（半不变量）；定理 5.7.2 补上 $K_1$ 对线心曲线在移轴下也不变；定理 5.7.3--5.7.5 分别给出中心、无心、线心曲线的简化方程；定理 5.7.6 给出九类曲线的判定表。}{This section compresses the chapter: the shape and size of a conic are completely determined by four numbers $I_1,I_2,I_3,K_1$, with no need to solve systems or find angles. Definition 5.7.1 makes ``invariant'' precise; Theorem 5.7.1 proves that $I_1,I_2,I_3$ are invariant under every rectangular change while $K_1$ is only a seminvariant (invariant under rotation); Theorem 5.7.2 adds that for a line-central conic $K_1$ is also invariant under translation; Theorems 5.7.3--5.7.5 give the simplified equations for the three types; Theorem 5.7.6 gives the classification table.}

\begin{definition}[5.7.1 不变量与半不变量 / invariants and seminvariants]\label{def:invariant}
由 $x,y$ 的系数组成的一个非常数函数 $f$，如果经过直角坐标变换 (5.6-3)，$F(x,y)$ 变为 $F'(x',y')$ 时，有
\[
f(a_{11},a_{12},\dots,a_{33})=f(a'_{11},a'_{12},\dots,a'_{33}),
\]
那么这个函数 $f$ 叫做二次曲线 (1) 在直角坐标变换 (5.6-3) 下的\textbf{不变量 / invariant}。如果这个函数 $f$ 的值，只是经过转轴变换不变，那么这个函数叫做二次曲线 (1) 在直角坐标变换下的\textbf{半不变量 / seminvariant}。
\EN{A non-constant function $f$ of the coefficients that takes the same value for the transformed coefficients under every rectangular change is an invariant; if it is unchanged only under rotations it is a seminvariant.}
\end{definition}

\begin{theorem}[5.7.1 三个不变量与一个半不变量 / three invariants and a seminvariant]\label{thm:invariants}
二次曲线 (1) 在直角坐标变换下，有三个不变量 $I_1,I_2,I_3$ 与一个半不变量 $K_1$。
\EN{Under rectangular changes a conic has three invariants $I_1,I_2,I_3$ and one seminvariant $K_1$.}
\end{theorem}

\begin{proof}[证明 / Proof]
因为直角坐标变换 (5.6-3) 总可以分成移轴 (5.6-1) 与转轴 (5.6-2) 两步来完成，因此本定理的证明也就分成移轴与转轴两步来完成。

先证明在移轴 (5.6-1) 下 $I_1,I_2,I_3$ 不变，而 $K_1$ 一般要改变。根据 (5.6-6) 知，在移轴下二次曲线 (1) 的二次项系数不变，所以 $I_1=a_{11}+a_{22}$ 与 $I_2=a_{11}a_{22}-a_{12}^2$ 不变；而移轴后的新矩阵为
\[
A'=\begin{pmatrix}a_{11}&a_{12}&F_1(x_0,y_0)\\ a_{12}&a_{22}&F_2(x_0,y_0)\\ F_1(x_0,y_0)&F_2(x_0,y_0)&F(x_0,y_0)\end{pmatrix},
\]
它是由原矩阵 $A$ 先作列变换``第三列加上第一列的 $x_0$ 倍与第二列的 $y_0$ 倍''、再作相应的行变换得到的，而这两种初等变换都不改变行列式，故 $I'_3=\det A'=\det A=I_3$。$K_1$ 在移轴下一般要改变：例如 $F(x,y)\equiv 2xy$ 的 $K_1=0$，而经移轴 $x=x'+x_0$，$y=y'+y_0$ 后 $F=2x'y'+2y_0x'+2x_0y'+2x_0y_0$，此时 $K'_1=-(x_0^2+y_0^2)$，一般不为零。

再证明在转轴 (5.6-2) 下 $I_1,I_2,I_3$ 与 $K_1$ 都不变。转轴时整张矩阵作正交合同变换 $A\mapsto P^{\mathsf T}AP$，其中 $P=\begin{pmatrix}\cos\alpha&-\sin\alpha&0\\ \sin\alpha&\cos\alpha&0\\ 0&0&1\end{pmatrix}$ 为正交矩阵，故 $\det A$ 不变，即 $I_3$ 不变；而 $A^*$ 同时作 $A^*\mapsto P_0^{\mathsf T}A^*P_0$（$P_0$ 为左上角二阶块），故迹 $I_1$ 与行列式 $I_2$ 也不变。$K_1$ 在转轴下不变可用 (5.6-7) 直接代入验证（书中的做法是把 $K_1$ 展开为 $I_1a_{33}-(a_{13}^2+a_{23}^2)$ 的形式，再用 $I_1,a_{33}$ 不变与 $(a_{13},a_{23})$ 按旋转矩阵变换的事实）。定理证毕。
\EN{Since a rectangular change is a translation followed by a rotation, check both steps. Translation keeps the quadratic coefficients, hence $I_1,I_2$; the translated matrix is obtained from $A$ by elementary row and column operations with multipliers $x_0,y_0$, which preserve the determinant, so $I_3$ is unchanged. $K_1$ generally changes under translation: for $F=2xy$ we have $K_1=0$, but after the translation $K'_1=-(x_0^2+y_0^2)$. Under rotation $A\mapsto P^{\mathsf T}AP$ with $P$ orthogonal, so $I_3=\det A$ is unchanged, and $A^*\mapsto P_0^{\mathsf T}A^*P_0$ leaves $I_1$ and $I_2$ unchanged; a direct substitution from (5.6-7) shows $K_1$ is unchanged.}
\end{proof}

\begin{theorem}[5.7.2 线心曲线的 $K_1$ / $K_1$ for line-central conics]\label{thm:k1inv}
当二次曲线 (1) 为线心曲线时，在直角坐标变换下 $K_1$ 是不变量。
\EN{For a line-central conic $K_1$ is invariant under rectangular changes.}
\end{theorem}

\APS{证明思路（书中第 163--164 页）：先对简化方程 (III) $a_{22}y^2+a_{33}=0$ 证明，此时 $K_1$ 是半不变量，故只需证它在移轴下不变；在移轴下 (III) 的左端变为 $a_{22}(y'+y_0)^2+a_{33}=a_{22}y'^2+2a_{22}y_0y'+(a_{22}y_0^2+a_{33})$，直接计算得 $K'_1=K_1$。再对一般的线心曲线 $F(x,y)=0$：因它总可通过某个直角坐标变换 $\tau_1$ 化为 (III)，而 $(III)$ 又可通过 $\tau_2$ 变为 $F'$，利用坐标变换的复合与 $K_1$ 在移轴下的不变性，即得 $K_1$ 不变。}{Sketch: for (III) the computation under translation gives $K'_1=K_1$; for a general line-central conic, compose the transformation to (III) with the transformation from (III) to the given one, using invariance under translation.}

\begin{theorem}[5.7.3 中心曲线的简化方程 / central conics]\label{thm:inv-central}
如果二次曲线 (1) 是中心曲线，那么它的简化方程为
\[
\lambda_1x^2+\lambda_2y^2+\frac{I_3}{I_2}=0,\tag{5.7-1}
\]
其中 $\lambda_1,\lambda_2$ 是二次曲线特征方程的两个根（方程中的撇号已略去）。
\EN{For a central conic the simplified equation is (5.7-1), where $\lambda_1,\lambda_2$ are the roots of the characteristic equation.}
\end{theorem}

\APS{证明思路：中心曲线经移轴到中心后化为 $a_{11}x^2+2a_{12}xy+a_{22}y^2+a_{33}=0$，再转轴到主方向化为 $\lambda_1x^2+\lambda_2y^2+a_{33}=0$。此时 $I_1=\lambda_1+\lambda_2$，$I_2=\lambda_1\lambda_2$，$I_3=\lambda_1\lambda_2a_{33}=I_2a_{33}$，故 $a_{33}=I_3/I_2$，即得 (5.7-1)。}{After translating to the centre and rotating to the principal directions, $I_2=\lambda_1\lambda_2$ and $I_3=\lambda_1\lambda_2a_{33}$, so $a_{33}=I_3/I_2$.}

\EX{（例 1，书中第 165--166 页）求二次曲线
\[
5x^2-6xy+5y^2-6\sqrt2\,x+2\sqrt2\,y-4=0
\]
的简化方程与标准方程。}{（Example 1, pp.~165--166）Find the simplified and standard equations of the displayed conic.}
这里 $a_{11}=5$，$a_{12}=-3$，$a_{22}=5$，$a_{13}=-3\sqrt2$，$a_{23}=\sqrt2$，$a_{33}=-4$。计算
\[
I_1=5+5=10,\qquad I_2=\begin{vmatrix}5&-3\\ -3&5\end{vmatrix}=25-9=16,
\]
\[
I_3=\begin{vmatrix}5&-3&-3\sqrt2\\ -3&5&\sqrt2\\ -3\sqrt2&\sqrt2&-4\end{vmatrix}=-128.
\]
所以
\[
\frac{I_3}{I_2}=\frac{-128}{16}=-8,
\]
而特征方程 $\lambda^2-10\lambda+16=0$ 的两根为
\[
\lambda_1=2,\qquad \lambda_2=8.
\]
所以曲线的简化方程（略去撇号）为
\[
2x^2+8y^2-8=0,
\]
曲线的标准方程（略去撇号）为
\[
\frac{x^2}{4}+\frac{y^2}{1}=1,
\]
这是一个椭圆。
\EN{With $I_1=10$, $I_2=16$, $I_3=-128$ we get $I_3/I_2=-8$; the characteristic roots are $2$ and $8$, so the simplified equation is $2x^2+8y^2-8=0$ and the standard one is $x^2/4+y^2/1=1$, an ellipse.}

\begin{theorem}[5.7.4 无心曲线的简化方程 / non-central conics]\label{thm:inv-noncentral}
如果二次曲线 (1) 是无心曲线，那么它的简化方程为
\[
I_1y^2\pm2\sqrt{-\frac{I_3}{I_1}}\,x=0,\tag{5.7-2}
\]
这里的正负号可以任意选取（方程中的撇号已略去）。
\EN{For a non-central conic the simplified equation is (5.7-2), with either sign.}
\end{theorem}

\APS{证明思路：无心曲线经坐标变换化为 (II) $a'_{22}y'^2+2a'_{13}x'=0$。此时 $I'_1=a'_{22}=I_1$，
\[
I'_3=\begin{vmatrix}0&0&a'_{13}\\ 0&a'_{22}&0\\ a'_{13}&0&0\end{vmatrix}=-a'_{22}a'^2_{13}=-I_1a'^2_{13},
\]
又 $I'_3=I_3$，所以 $a'_{13}=\pm\sqrt{-I_3/I_1}$，即得 (5.7-2)。}{In the simplified form (II) one has $I'_1=a'_{22}=I_1$ and $I'_3=-I_1a'^2_{13}=I_3$, hence $a'_{13}=\pm\sqrt{-I_3/I_1}$.}

\EX{（例 2，书中第 166--167 页）求二次曲线
\[
x^2-2xy+y^2-10x-6y+25=0
\]
的简化方程与标准方程。}{（Example 2, pp.~166--167）Find the simplified and standard equations of the displayed conic.}
因为
\[
I_1=1+1=2,\qquad I_2=\begin{vmatrix}1&-1\\ -1&1\end{vmatrix}=0,
\]
\[
I_3=\begin{vmatrix}1&-1&-5\\ -1&1&-3\\ -5&-3&25\end{vmatrix}=-64,
\]
所以曲线的简化方程（略去撇号）为
\[
2y^2-2\sqrt{\frac{64}{2}}\,x=0\quad\text{或}\quad 2y^2+2\sqrt{\frac{64}{2}}\,x=0,
\]
即
\[
2y^2-2\sqrt{32}\,x=0\quad\text{或}\quad 2y^2+2\sqrt{32}\,x=0.
\]
它的标准方程（略去撇号）为
\[
y^2=4\sqrt2\,x\quad\text{或}\quad y^2=-4\sqrt2\,x,
\]
曲线是一条抛物线。
\EN{Here $I_1=2$, $I_2=0$, $I_3=-64$, so (5.7-2) gives $2y^2\pm2\sqrt{32}\,x=0$, i.e.\ the standard form $y^2=\pm4\sqrt2\,x$: a parabola.}

\begin{theorem}[5.7.5 线心曲线的简化方程 / line-central conics]\label{thm:inv-linecentral}
如果二次曲线 (1) 是线心曲线，那么它的简化方程为
\[
I_1y^2+\frac{K_1}{I_1}=0,\tag{5.7-3}
\]
（方程中的撇号已略去）。
\EN{For a line-central conic the simplified equation is (5.7-3).}
\end{theorem}

\APS{证明思路：线心曲线经坐标变换化为 (III) $a'_{22}y'^2+a'_{33}=0$。此时 $I'_1=a'_{22}=I_1$，而 $K'_1=a'_{22}a'_{33}=I_1a'_{33}$；又 $K_1$ 是线心曲线的不变量（定理 5.7.2），所以 $a'_{33}=K_1/I_1$，即得 (5.7-3)。}{For (III), $I'_1=a'_{22}=I_1$ and $K'_1=a'_{22}a'_{33}=I_1a'_{33}$; by Theorem 5.7.2 $K_1$ is invariant, hence $a'_{33}=K_1/I_1$.}

\APS{从 (5.7-1)、(5.7-2) 与 (5.7-3) 又可以得到下面的分类定理。}{From (5.7-1)--(5.7-3) follows the classification theorem.}

\begin{theorem}[5.7.6 用不变量判定曲线类型 / classification by invariants]\label{thm:inv-class}
如果给出了二次曲线 (1)，那么用它的不变量来判断已知曲线为何种曲线的条件是：
\[
\begin{gathered}
{[1]}\ \text{椭圆：}I_2>0,\ I_1I_3<0;\qquad
[2]\ \text{虚椭圆：}I_2>0,\ I_1I_3>0;\\
[3]\ \text{点（或称一对交于实点的共轭虚直线）：}I_2>0,\ I_3=0;\\
[4]\ \text{双曲线：}I_2<0,\ I_3\ne0;\qquad
[5]\ \text{一对相交直线：}I_2<0,\ I_3=0;\\
[6]\ \text{抛物线：}I_2=0,\ I_3\ne0;\qquad
[7]\ \text{一对平行直线：}I_2=I_3=0,\ K_1<0;\\
[8]\ \text{一对平行的共轭虚直线：}I_2=I_3=0,\ K_1>0;\qquad
[9]\ \text{一对重合的直线：}I_2=I_3=K_1=0.
\end{gathered}
\]
\EN{The invariants classify the conic as: ellipse ($I_2>0$, $I_1I_3<0$), imaginary ellipse ($I_2>0$, $I_1I_3>0$), a point ($I_2>0$, $I_3=0$), hyperbola ($I_2<0$, $I_3\ne0$), two intersecting lines ($I_2<0$, $I_3=0$), parabola ($I_2=0$, $I_3\ne0$), two parallel lines ($I_2=I_3=0$, $K_1<0$), two parallel conjugate imaginary lines ($I_2=I_3=0$, $K_1>0$), two coincident lines ($I_2=I_3=K_1=0$).}
\end{theorem}

\APS{这个定理的证明与定理 5.6.2 十分类似，它的证明留给读者。推论：二次曲线 (1) 表示两条直线（实的或虚的，不同的或重合的）的充要条件为 $I_3=0$。}{The proof is analogous to that of Theorem 5.6.2 and is left to the reader. Corollary: a conic represents two lines (real or imaginary, distinct or coincident) if and only if $I_3=0$.}

\EX{（补充算例，仿书中 §5.7 习题 1）利用不变量判断下列二次曲线为何种曲线，并求出它们的简化方程与标准方程：(1) $x^2+6xy+y^2+6x+2y-1=0$；(2) $4x^2-4xy+y^2+6x-8y+3=0$。}{（Supplementary examples, after Exercise 1 of §5.7）Classify by invariants and find the simplified and standard equations: (1) $x^2+6xy+y^2+6x+2y-1=0$; (2) $4x^2-4xy+y^2+6x-8y+3=0$.}
(1) $a_{11}=1$，$a_{12}=3$，$a_{22}=1$，$a_{13}=3$，$a_{23}=1$，$a_{33}=-1$：
\[
I_1=2,\qquad I_2=1-9=-8,\qquad
I_3=\begin{vmatrix}1&3&3\\ 3&1&1\\ 3&1&-1\end{vmatrix}=1(-1-1)-3(-3-3)+3(3-3)=-2+18=16.
\]
$I_2<0$，$I_3\ne0$，所以是\textbf{双曲线}。特征方程 $\lambda^2-2\lambda-8=0$ 的根为 $\lambda_1=4$，$\lambda_2=-2$，$I_3/I_2=-2$，所以简化方程为
\[
4x^2-2y^2-2=0\quad\text{即}\quad 2x^2-y^2=1,
\]
标准方程为
\[
\frac{x^2}{\frac12}-y^2=1.
\]
(2) $a_{11}=4$，$a_{12}=-2$，$a_{22}=1$，$a_{13}=3$，$a_{23}=-4$，$a_{33}=3$：
\[
I_1=5,\qquad I_2=4-4=0,
\]
\[
I_3=\begin{vmatrix}4&-2&3\\ -2&1&-4\\ 3&-4&3\end{vmatrix}=4(3-16)+2(-6+12)+3(8-3)=-52+12+15=-25\ne0,
\]
所以是\textbf{抛物线}。由 (5.7-2)，
\[
5y^2\pm2\sqrt{\frac{25}{5}}\,x=0\quad\text{即}\quad 5y^2\pm2\sqrt5\,x=0,
\]
标准方程为
\[
y^2=\pm\frac{2\sqrt5}{5}x.
\]
\EN{(1) $I_2=-8<0$, $I_3=16\ne0$: hyperbola; the characteristic roots are $4,-2$ and $I_3/I_2=-2$, so the simplified equation is $4x^2-2y^2-2=0$, i.e.\ $2x^2-y^2=1$. (2) $I_2=0$, $I_3=-25\ne0$: parabola; (5.7-2) gives $5y^2\pm2\sqrt5\,x=0$, i.e.\ $y^2=\pm\frac{2\sqrt5}{5}x$.}

\EX{（补充算例）用不变量判断 $x^2-2xy+y^2+2x-2y-3=0$ 的类型，并求其简化方程。}{（Supplementary example）Classify $x^2-2xy+y^2+2x-2y-3=0$ by invariants and find its simplified equation.}
$a_{11}=1$，$a_{12}=-1$，$a_{22}=1$，$a_{13}=1$，$a_{23}=-1$，$a_{33}=-3$：
\[
I_1=2,\qquad I_2=0,\qquad
I_3=\begin{vmatrix}1&-1&1\\ -1&1&-1\\ 1&-1&-3\end{vmatrix}=1(-3-1)+1(3+1)+1(1-1)=0,
\]
\[
K_1=\begin{vmatrix}1&1\\ 1&-3\end{vmatrix}+\begin{vmatrix}1&-1\\ -1&-3\end{vmatrix}=(-3-1)+(-3-1)=-8<0.
\]
$I_2=I_3=0$，$K_1<0$，所以是\textbf{一对平行直线}。由 (5.7-3)，
\[
I_1y^2+\frac{K_1}{I_1}=0\ \Longrightarrow\ 2y^2-4=0\ \Longrightarrow\ y^2=2,
\]
即两条平行直线 $y=\pm\sqrt2$（与原方程分解式 $(x-y+3)(x-y-1)=0$ 一致，只是坐标系不同）。
\EN{Here $I_1=2$, $I_2=I_3=0$, $K_1=-8<0$: a pair of parallel lines; (5.7-3) gives $2y^2-4=0$, i.e.\ $y^2=2$, two parallel lines.}

\ASK{$I_1,I_2,I_3$ 为什么在直角坐标变换下不变？$K_1$ 为什么只是半不变量？}{Why are $I_1,I_2,I_3$ invariant under rectangular changes, while $K_1$ is only a seminvariant?}{$I_1$ 与 $I_2$ 是二次型 $A^*$ 的迹与行列式，转轴对应二次型的合同变换，迹与行列式在合同变换下不变；$I_3$ 是 $A$ 的行列式，移轴相当于对 $A$ 作``第三行减去前两行的线性组合''与相应的列变换，行列式不变，转轴同理。$K_1$ 由 $a_{13},a_{23}$ 参与构成，而 $a_{13},a_{23}$ 在移轴下变为 $F_1,F_2$（一般不为零），所以 $K_1$ 在移轴下一般改变；只有在``一次项系数与 $K_1$ 的相关组合''恰好抵消的特殊情形（线心曲线）才保持不变。因此 $K_1$ 只是转轴不变量（半不变量），对线心曲线才升级为不变量。}{$I_1,I_2$ are the trace and determinant of the quadratic form $A^*$, invariant under congruence; $I_3=\det A$ is unchanged by translation because translation only adds multiples of the first two rows/columns to the third. $K_1$ involves $a_{13},a_{23}$, which under translation become $F_1,F_2$ and generally change, so $K_1$ changes under translation; only for line-central conics do the terms cancel, making $K_1$ an invariant there. Hence $K_1$ is a rotation invariant (seminvariant) in general.}

\ASK{给一个二次曲线方程，如何用不变量完成``判别--化简--标准方程''一条龙？}{Given a conic, how do invariants give classification, simplified equation and standard equation in one go?}{第一步算 $I_1,I_2,I_3,K_1$；第二步查定理 5.7.6 的九类判定表确定类型；第三步按类型套简化方程：中心曲线用 $\lambda_1x^2+\lambda_2y^2+I_3/I_2=0$（$\lambda$ 为特征根），无心曲线用 $I_1y^2\pm2\sqrt{-I_3/I_1}\,x=0$，线心曲线用 $I_1y^2+K_1/I_1=0$；第四步两边除以常数化为标准方程，读出半轴、焦距、开口方向等度量信息。注意 $I_2\ne0$ 时 $I_3$ 的符号决定实虚：$I_1I_3<0$ 为实椭圆，$I_1I_3>0$ 为虚椭圆。}{Compute $I_1,I_2,I_3,K_1$; read the type from Theorem 5.7.6; write the simplified equation according to the type: $\lambda_1x^2+\lambda_2y^2+I_3/I_2=0$ for central, $I_1y^2\pm2\sqrt{-I_3/I_1}x=0$ for non-central, $I_1y^2+K_1/I_1=0$ for line-central; divide by constants to get the standard form and read off the semi-axes, focal distance and opening direction. When $I_2\ne0$ the sign of $I_1I_3$ decides real versus imaginary ellipse.}

\section{\label{sec:16}第五章口试问答 | Chapter 5 interview questions}
\ASK{二次曲线的几何理论可以按哪几个``层次''组织？}{At which levels is the geometric theory of conics organised?}{四个层次：第一层``方向''——用 $\Phi(X,Y)$ 定义渐近方向，得椭圆型/抛物型/双曲型；第二层``点''——用 $F_1=F_2=0$ 定义中心，得中心/无心/线心；第三层``线与弦''——切线（判别式为零）、直径（平行弦中点轨迹）、主直径（垂直于共轭弦的直径）；第四层``方程''——坐标变换化简、九种标准方程、不变量分类。每一层都是前一层代数条件的几何翻译，最终由 $I_1,I_2,I_3,K_1$ 四个数统一。}{Four levels: directions ($\Phi=0$ gives elliptic/parabolic/hyperbolic); points ($F_1=F_2=0$ gives central/non-central/line-central); lines and chords (tangents via vanishing discriminant, diameters as loci of midpoints, principal diameters as perpendicular to conjugate chords); equations (reduction, nine standard forms, invariants). Each level translates the algebraic conditions of the previous one, and everything is unified by $I_1,I_2,I_3,K_1$.}

\ASK{为什么说''直径''概念把二次曲线的``线性化''体现得最充分？}{Why does the notion of diameter best display the ``linearisation'' of a conic?}{二次曲线是二次的，但一族平行弦的中点轨迹却由一次方程 $XF_1+YF_2=0$ 给出，这是因为二次方程两根之和由一次项系数决定（韦达定理），而``中点为参数 $0$''正好把两根之和置零。更进一步，主直径由特征向量给出，把二次型对角化；于是二次曲线的所有对称信息都被压缩到矩阵 $A^*$ 的特征值与特征向量上——这正是``用线性代数研究二次曲线''的核心。}{A conic is quadratic, yet the midpoints of parallel chords satisfy the linear equation $XF_1+YF_2=0$, because the sum of the two roots is governed by the linear coefficient (Viète), and the midpoint corresponds to parameter zero. Moreover principal diameters come from eigenvectors, diagonalising the quadratic form; all symmetry information is encoded in the eigenvalues and eigenvectors of $A^*$---the core of studying conics by linear algebra.}

\ASK{退化二次曲线（两条直线）有哪些类型？不变量如何识别？}{What are the degenerate conics (two lines), and how do invariants recognise them?}{退化情形共五类：两条相交实直线（$I_2<0,I_3=0$）、两条平行实直线（$I_2=I_3=0,K_1<0$）、两条平行共轭虚直线（$I_2=I_3=0,K_1>0$）、两条重合直线（$I_2=I_3=K_1=0$）、一点（$I_2>0,I_3=0$，即两条相交于实点的共轭虚直线）。共同特征是 $I_3=0$（曲线分解为两条直线，实或虚）；再用 $I_2$ 的符号区分相交/平行，用 $K_1$ 的符号区分平行直线的实虚。}{There are five degenerate cases: two intersecting real lines ($I_2<0$, $I_3=0$); two parallel real lines ($I_2=I_3=0$, $K_1<0$); two parallel conjugate imaginary lines ($I_2=I_3=0$, $K_1>0$); two coincident lines ($I_2=I_3=K_1=0$); a point ($I_2>0$, $I_3=0$, a pair of conjugate imaginary lines meeting in a real point). All have $I_3=0$; the sign of $I_2$ separates intersecting from parallel, and $K_1$ separates real from imaginary parallel pairs.}

\ASK{二次曲线的特征根与椭圆的半轴、双曲线的半轴有什么关系？}{How do the eigenvalues relate to the semi-axes of an ellipse or hyperbola?}{对中心曲线，简化方程 $\lambda_1x^2+\lambda_2y^2+I_3/I_2=0$ 即 $\frac{x^2}{-I_3/(I_2\lambda_1)}+\frac{y^2}{-I_3/(I_2\lambda_2)}=1$。若 $\lambda_1,\lambda_2$ 同号且 $I_3/I_2<0$，得椭圆，半轴为 $\sqrt{-I_3/(I_2\lambda_1)}$ 与 $\sqrt{-I_3/(I_2\lambda_2)}$；若异号则得双曲线，实半轴与虚半轴由同一公式给出（取正量）。可见特征根决定两条主直径上的``伸缩率''，$I_3/I_2$ 决定整体大小。}{For a central conic the simplified equation gives $\frac{x^2}{-I_3/(I_2\lambda_1)}+\frac{y^2}{-I_3/(I_2\lambda_2)}=1$. If the eigenvalues have the same sign and $I_3/I_2<0$ this is an ellipse with semi-axes $\sqrt{-I_3/(I_2\lambda_1)}$, $\sqrt{-I_3/(I_2\lambda_2)}$; if they have opposite signs it is a hyperbola with the same formulas for its semi-axes. Thus the eigenvalues give the stretching along the principal diameters and $I_3/I_2$ the overall size.}

\section*{中英术语表 | Bilingual glossary}\label{sec:glossary}
\begin{small}
\begin{longtable}{@{}p{3.4cm}p{4.6cm}p{7.4cm}@{}}
\hline
中文 Chinese & English & 说明 / Note \\
\hline
\endhead
柱面 & cylinder & 一族平行直线沿准线滑动所生成的曲面（定义 4.1.1） \\
准线 & directrix & 母线族所沿的定曲线 \\
母线 & generator & 生成曲面的直线（或曲线）族中的一条 \\
柱面的方向 & direction of a cylinder & 母线的公共方向 \\
二次柱面 & quadric cylinder & 方程二次的柱面：椭圆柱面、双曲柱面、抛物柱面 \\
椭圆柱面 & elliptic cylinder & $x^2/a^2+y^2/b^2=1$ \\
双曲柱面 & hyperbolic cylinder & $x^2/a^2-y^2/b^2=1$ \\
抛物柱面 & parabolic cylinder & $y^2=2px$ \\
射影柱面 & projecting cylinder & 空间曲线对坐标面射影所得的柱面 \\
射影曲线 & projection & 射影柱面与坐标面的交线 \\
锥面 & cone & 一族过定点的直线沿准线滑动所生成的曲面（定义 4.2.1） \\
顶点 & vertex & 锥面的公共点；抛物线与轴的切点；轴与圆锥曲线的交点 \\
二次锥面 & quadratic cone & $x^2/a^2+y^2/b^2-z^2/c^2=0$ \\
齐次方程 & homogeneous equation & 每项次数相同的方程；齐次方程表示以原点为顶点的锥面（定理 4.2.1） \\
虚锥面 & imaginary cone & 只有实顶点、无其他实点的齐次方程曲面 \\
旋转曲面 & surface of revolution & 一条曲线绕定直线旋转一周生成的曲面（定义 4.3.1） \\
旋转轴 & axis of revolution & 旋转所绕的定直线 \\
纬圆 & latitude circle & 母线上的点旋转形成的圆 \\
经线 & meridian & 含轴半平面与旋转曲面的交线 \\
替换规律 & substitution rule & 保留与轴同名的坐标，另一坐标换成其余两坐标平方和的开方 \\
长形旋转椭球面 & prolate spheroid & $x^2/a^2+(y^2+z^2)/b^2=1$（$a>b$） \\
扁形旋转椭球面 & oblate spheroid & $(x^2+z^2)/a^2+y^2/b^2=1$（$a>b$） \\
单叶旋转双曲面 & one-sheeted hyperboloid of revolution & $(x^2+y^2)/b^2-z^2/c^2=1$ \\
双叶旋转双曲面 & two-sheeted hyperboloid of revolution & $y^2/b^2-(x^2+z^2)/c^2=1$ \\
旋转抛物面 & paraboloid of revolution & $x^2+y^2=2pz$ \\
环面 & torus & $(x^2+y^2+z^2+b^2-a^2)^2=4b^2(x^2+y^2)$ \\
椭球面 & ellipsoid & $x^2/a^2+y^2/b^2+z^2/c^2=1$（定义 4.4.1） \\
主平面 & principal plane & 椭球面的对称平面（坐标面） \\
主轴 & principal axis & 椭球面的对称轴（坐标轴） \\
半轴 & semi-axis & 中心到顶点的线段长度 $a,b,c$ \\
平行截割法 & method of parallel sections & 用一组平行平面截曲面，由截口曲线变化读出曲面形状 \\
主截线 & principal section & 用坐标面截椭球面所得的椭圆 \\
参数方程 & parametric equation & 用参数表示曲面上点的坐标，如 (4.4-2) \\
单叶双曲面 & hyperboloid of one sheet & $x^2/a^2+y^2/b^2-z^2/c^2=1$（定义 4.5.1） \\
双叶双曲面 & hyperboloid of two sheets & $x^2/a^2+y^2/b^2-z^2/c^2=-1$（定义 4.5.2） \\
腰椭圆 & waist ellipse & 单叶双曲面在 $z=0$ 上的截线 \\
双曲面 & hyperboloid & 单叶与双叶双曲面的统称 \\
椭圆抛物面 & elliptic paraboloid & $x^2/a^2+y^2/b^2=2z$（定义 4.6.1） \\
双曲抛物面 & hyperbolic paraboloid & $x^2/a^2-y^2/b^2=2z$（定义 4.6.2） \\
主抛物线 & principal parabola & 用坐标面截抛物面所得的抛物线 \\
马鞍曲面 & saddle surface & 双曲抛物面的俗称 \\
无心二次曲面 & quadric without centre & 抛物面（无对称中心） \\
直纹曲面 & ruled surface & 由一族直线生成的曲面 \\
直母线 & generator (of a ruled surface) & 直纹曲面上属于该族的直线 \\
族 & family & 含一个或两个参数的直线集合 \\
共面 & coplanar & 两直线在同一平面内（异族直母线的性质） \\
异面 & skew & 两直线不共面（同族直母线的性质） \\
二次曲线 & conic & 二元二次方程所表示的平面曲线 \\
二次曲线的矩阵 & matrix of a conic & 系数对称矩阵 $A$（含 $a_{33}$） \\
不变量 & invariant & 直角坐标变换下不变的系数函数（定义 5.7.1） \\
半不变量 & seminvariant & 仅在转轴下不变的系数函数，如 $K_1$ \\
特征方程 & characteristic equation & $\lambda^2-I_1\lambda+I_2=0$（定义 5.5.2） \\
特征根 & eigenvalue & 特征方程的根 $\lambda_1,\lambda_2$ \\
渐近方向 & asymptotic direction & 满足 $\Phi(X,Y)=0$ 的方向（定义 5.2.1） \\
椭圆型 & elliptic & 无实渐近方向，$I_2>0$ \\
抛物型 & parabolic & 一个实渐近方向，$I_2=0$ \\
双曲型 & hyperbolic & 两个实渐近方向，$I_2<0$ \\
弦 & chord & 直线与二次曲线两交点间的线段 \\
中心 & centre & 所有过它的弦的中点（定义 5.2.3） \\
中心直线 & line of centres & 线心曲线的中心轨迹 \\
渐近线 & asymptote & 过中心、以渐近方向为方向的直线（定义 5.2.5） \\
中心曲线 & central conic & 有惟一中心，$I_2\ne0$ \\
无心曲线 & non-central conic (no centre) & $I_2=0$ 且中心方程组矛盾 \\
线心曲线 & line-central conic & $I_2=0$ 且中心为一条直线 \\
切线 & tangent & 与二次曲线交于两个重合点的直线（定义 5.3.1） \\
切点 & point of tangency & 重合的交点 \\
奇点 & singular point & 曲线上满足 $F_1=F_2=0$ 的点（定义 5.3.2） \\
正则点 & regular point & 非奇点 \\
直径 & diameter & 一族平行弦中点的轨迹（定义 5.4.1） \\
共轭弦 & conjugate chords & 与某直径对应的那族平行弦 \\
共轭方向 & conjugate direction & 满足 (5.4-5) 的两个方向 \\
共轭直径 & conjugate diameters & 方向互相共轭的一对直径（定义 5.4.2） \\
主直径 & principal diameter & 垂直于其共轭弦的直径，即对称轴（定义 5.5.1） \\
主方向 & principal direction & 主直径方向及其垂直方向 \\
坐标变换 & coordinate transformation & 移轴与转轴的复合（5.6-3） \\
移轴 & translation of axes & 只改变原点位置（5.6-1） \\
转轴 & rotation of axes & 只改变轴的方向（5.6-2） \\
化简 & reduction & 用坐标变换把方程化为最简形式 \\
标准方程 & standard equation & 九种标准形式之一（定理 5.6.2） \\
分类 & classification & 按渐近方向、按中心或按不变量分类 \\
退化 & degenerate & 曲线分解为直线、点或虚图形 \\
虚椭圆 & imaginary ellipse & $x^2/a^2+y^2/b^2=-1$ \\
两平行直线 & pair of parallel lines & $y^2=a^2$ \\
两重合直线 & double line & $y^2=0$ \\
点圆 & point circle & 退化为一点的圆 \\
\end{longtable}
\end{small}

\section*{参考文献 | References}\label{sec:refs}
\noindent 吕林根，许子道．解析几何（第五版）．北京：高等教育出版社．第 4 章（柱面、锥面、旋转曲面与二次曲面，书内第 96--127 页）、第 5 章（二次曲线的一般理论，书内第 128--169 页）；第 4.1 节柱面与射影柱面、第 4.2 节锥面与齐次方程、第 4.3 节旋转曲面与替换规律、第 4.4--4.6 节椭球面/双曲面/抛物面的平行截割法与标准方程、第 4.7 节单叶双曲面与双曲抛物面的直母线；第 5.1 节记号与直线相关位置、第 5.2 节渐近方向/中心/渐近线、第 5.3 节切线、第 5.4 节直径与共轭方向、第 5.5 节主直径与主方向、第 5.6 节坐标变换与九种标准方程、第 5.7 节不变量与分类．本讲义的定义、定理编号均与教材一致，书内页码加 $8$ 即本扫描 PDF 的物理页码．\par
\noindent 尤承业．解析几何．北京：北京大学出版社．第 3--4 章（二次曲面、二次曲线的一般理论）可作平行阅读，尤其是不变量与主方向的计算．\par
\noindent 丘维声．解析几何．北京：北京大学出版社．第 4 章（二次曲面与二次曲线的方程化简）、第 5 章（二次曲线的一般理论）．其对二次型与矩阵方法的处理更偏代数，可与本讲义的矩阵记号 $A,A^*$ 对照．\par
\noindent 同济大学数学系．高等数学（第七版）下册．北京：高等教育出版社．空间解析几何一章（曲面方程、二次曲面）可作为工科背景的快速回顾．\par
\EN{Main source: Lü Lingen and Xu Zidao, \emph{Analytic Geometry}, 5th edition, Higher Education Press, Chapters 4--5 (pp.~96--127 and 128--169 in the book; add 8 for the PDF of the scan). Parallel reading: You Chengye, \emph{Analytic Geometry}, Peking University Press; Qiu Weisheng, \emph{Analytic Geometry}, Peking University Press; Tongji University, \emph{Advanced Mathematics}, 7th ed., Vol.~2, for a quick engineering review of surfaces. References are chapter-level; the numbering of definitions and theorems in these notes agrees with the textbook.}

\end{document}
